% Transformation et recyclage des déchets — SVT Tle D — Cours signé OnBuch+
% Programme officiel MINESEC. Compiler avec : tectonic N20-transformation-recyclage-dechets.tex
\newcommand{\DOCMATIERE}{SVT}
\newcommand{\DOCNIVEAU}{Tle D}
\newcommand{\DOCMODULE}{Module IV : La Biotechnologie — Valorisation des déchets de l'environnement}
\newcommand{\DOCLECON}{N20}
\newcommand{\DOCTITRE}{Transformation et recyclage des déchets}
% OnBuch+ — préambule « cours signés » ENRICHI (compilé avec tectonic / XeLaTeX)
% Système visuel « Pop » d'OnBuch+ : fond crème (jamais blanc), bordures encre
% épaisses, ombres « dures » décalées, titres Archivo Black en capitales.
% Version MATHÉMATIQUES : graphes (pgfplots/tikz), boîtes pédagogiques
% (définition, propriété, méthode, attention, le savais-tu, exercice résolu,
% exercice, TP, bilan), page de garde et signature OnBuch+ ancrée en bas.
%
% Macros attendues dans main.tex AVANT \input{preamble} :
%   \DOCMATIERE  \DOCNIVEAU  \DOCMODULE  \DOCLECON (numéro)  \DOCTITRE
\documentclass[11pt]{article}

\usepackage{fontspec}
\usepackage[french]{babel}
\usepackage[a4paper,margin=2cm,top=3.4cm,bottom=2.8cm,headheight=1.4cm,footskip=1.3cm]{geometry}
\usepackage{amsmath,amssymb,amsfonts}
\usepackage{mathtools} % \xmapsto, \coloneqq... souvent utilisés par les modèles
\usepackage{cancel} % simplifications barrées (bilans de forces)
% \abs{z} (module, valeur absolue) et \norm{u} (norme), utilisés par les modèles
\providecommand{\abs}{}\let\abs\relax\DeclarePairedDelimiter{\abs}{\lvert}{\rvert}
\providecommand{\norm}{}\let\norm\relax\DeclarePairedDelimiter{\norm}{\lVert}{\rVert}
\usepackage[table]{xcolor} % [table] : fournit aussi \rowcolors (alternance de lignes)
\usepackage{pagecolor}
\usepackage{tikz}
\usetikzlibrary{arrows.meta,positioning,calc,shapes.geometric,shapes.misc,shapes.arrows,decorations.pathmorphing,decorations.pathreplacing,decorations.markings,patterns,shadows,mindmap,trees,fit,backgrounds,matrix,intersections,angles,quotes}
\usepackage{pgfplots}
\usepackage[version=4]{mhchem} % équations nucléaires \ce{^{235}_{92}U}
\usepackage[european,siunitx]{circuitikz} % schémas électriques
\pgfplotsset{compat=1.18}
\usepgfplotslibrary{fillbetween,groupplots} % aires entre courbes (calcul intégral)
\usepackage{enumitem}
\usepackage{fancyhdr}
% [most] charge toutes les bibliothèques tcolorbox (listings, minted, raster,
% documentation...) et gonfle inutilement la mémoire TeX sur un document long
% (~300 boîtes) : on ne charge que ce qui est réellement utilisé.
\usepackage[skins,breakable]{tcolorbox}
\usepackage{graphicx}
\usepackage{colortbl}
\usepackage{booktabs}
\usepackage{tabularx}
\usepackage{diagbox} % case d'en-tête diagonale des tableaux à double entrée
% Colonnes extensibles centrée / gauche / droite (C, L, R), souvent utilisées par les modèles
\newcolumntype{C}{>{\centering\arraybackslash}X}
\newcolumntype{L}{>{\raggedright\arraybackslash}X}
\newcolumntype{R}{>{\raggedleft\arraybackslash}X}
\usepackage{array}
\usepackage{multicol}
\usepackage{siunitx}
% parse-numbers=false : accepte aussi \SI{2,5 \times 10^{-3}}{...}
\sisetup{locale=FR,output-decimal-marker={,},group-minimum-digits=4,parse-numbers=false}
\DeclareSIUnit\ohm{\ensuremath{\symup{\Omega}}} % U+2126 absent de la police
\DeclareSIUnit\division{div} % réglages d'oscilloscope (ms/div, V/div)
\DeclareSIUnit\tr{tr} % tours (vitesse de rotation, ex. \SI{1500}{\tr\per\minute})
\DeclareSIUnit\molar{M} % concentration molaire (ex. \SI{3}{\molar}, \SI{1}{\micro\molar})
\DeclareSIUnit\an{an} % année (le modèle confond parfois avec \year, un compteur TeX) — ex. \SI{5}{\kilo\meter\per\an}
\DeclareSIUnit\FCFA{FCFA} % franc CFA (ex. \SI{500}{\FCFA\per\kilo\gram})
\DeclareSIUnit\degreeBrix{\textdegree Brix} % teneur en sucre d'un sirop/jus (réfractométrie)
\DeclareSIUnit\month{mois} % le modèle utilise parfois \month, un compteur TeX, comme unité de durée
\usepackage{qrcode}
% \degree : commande fréquemment utilisée par les modèles pour un angle ou une
% température en dehors de \SI{}{} (non fournie par les paquets chargés).
\providecommand{\degree}{^{\circ}}
% \qed (fin de démonstration), utilisé par les modèles sans amsthm
\providecommand{\qed}{\hfill$\square$}
% \barre : double barre des valeurs interdites dans un tableau de variations (tkz-tab)
\providecommand{\barre}{\ensuremath{\|}}
% environnement proof (amsthm non chargé) utilisé par les modèles
\ifdefined\proof\else\newenvironment{proof}[1][Démonstration]{\par\noindent\textit{#1.}\ }{\qed\par}\fi
\usepackage{lastpage}
\usepackage{hyperref}
\hypersetup{hidelinks}

% ============================================================
% Palette « Pop » OnBuch+ — mêmes hex que la classe P (plus_ui.dart)
% ============================================================
\definecolor{poporange}{HTML}{F59321}
\definecolor{popdark}{HTML}{E07A0C}
\definecolor{popcream}{HTML}{FFF6E8}
\definecolor{popink}{HTML}{1C1714}
\definecolor{poppurple}{HTML}{7A5AE0}
\definecolor{popgreen}{HTML}{2E9E6A}
\definecolor{poppink}{HTML}{E0457B}
\definecolor{popblue}{HTML}{2F7DE1}
\definecolor{popgold}{HTML}{C98A12}
\definecolor{popmuted}{HTML}{8A7F76}
\definecolor{popline}{HTML}{EADFD2}
\definecolor{popgray}{HTML}{8A7F76}
\colorlet{popgrey}{popgray}
% Alias tolérants (le modèle nomme parfois les couleurs différemment)
\colorlet{popwhite}{white}
\colorlet{popblack}{popink}
\colorlet{popred}{poppink}
\colorlet{popyellow}{popgold}
% Teintes claires (fonds de tableaux/encadrés) — le modèle nomme parfois
% les couleurs avec un suffixe L (ex. popblueL) sans les avoir définies.
\colorlet{poporangeL}{poporange!20}
\colorlet{popdarkL}{popdark!20}
\colorlet{poppurpleL}{poppurple!20}
\colorlet{popgreenL}{popgreen!20}
\colorlet{poppinkL}{poppink!20}
\colorlet{popblueL}{popblue!20}
\colorlet{popgoldL}{popgold!20}
\colorlet{popredL}{popred!20}
\colorlet{popgrayL}{popgray!20}
% Teintes claires pour fonds de tableaux / zones
\colorlet{popblueL}{popblue!10!white}
\colorlet{poporangeL}{poporange!14!white}
\colorlet{popgreenL}{popgreen!12!white}
\colorlet{poppurpleL}{poppurple!12!white}
\colorlet{poppinkL}{poppink!12!white}
\colorlet{popgoldL}{popgold!14!white}

\pagecolor{popcream}

% ============================================================
% Typographie
% ============================================================
\defaultfontfeatures{Ligatures=TeX}
\setmainfont{PlusJakartaSans-Medium.ttf}[Path=../../../fonts/,BoldFont=PlusJakartaSans-Bold.ttf]
% Maths en sans-serif (Fira Math) avec chiffres et lettres droites de Plus
% Jakarta, pour que les formules se fondent dans le texte.
\usepackage[math-style=ISO,bold-style=ISO]{unicode-math}
\setmathfont{FiraMath-Regular.otf}[Path=../../../fonts/]
\setmathfont{PlusJakartaSans-Medium.ttf}[Path=../../../fonts/,range={up/num,up/{Latin,latin},\mathup}]
\usepackage{icomma} % virgule décimale sans espace en mode math
% Caractères mathématiques Unicode absents des polices de texte (Plus Jakarta,
% Archivo Black) : rendus par leur équivalent mathématique, en texte comme en
% formule — sinon ils s'affichent en carré vide (ex. « Arithmétique dans ℤ »).
\usepackage{newunicodechar}
\newunicodechar{ℝ}{\ensuremath{\mathbb{R}}}
\newunicodechar{ℤ}{\ensuremath{\mathbb{Z}}}
\newunicodechar{ℂ}{\ensuremath{\mathbb{C}}}
\newunicodechar{ℕ}{\ensuremath{\mathbb{N}}}
\newunicodechar{ℚ}{\ensuremath{\mathbb{Q}}}
\newunicodechar{𝔻}{\ensuremath{\mathbb{D}}}
\newunicodechar{α}{\ensuremath{\alpha}}
\newunicodechar{β}{\ensuremath{\beta}}
\newunicodechar{γ}{\ensuremath{\gamma}}
\newunicodechar{δ}{\ensuremath{\delta}}
\newunicodechar{ε}{\ensuremath{\varepsilon}}
\newunicodechar{θ}{\ensuremath{\theta}}
\newunicodechar{λ}{\ensuremath{\lambda}}
\newunicodechar{μ}{\ensuremath{\mu}}
\newunicodechar{σ}{\ensuremath{\sigma}}
\newunicodechar{τ}{\ensuremath{\tau}}
\newunicodechar{φ}{\ensuremath{\varphi}}
\newunicodechar{ψ}{\ensuremath{\psi}}
\newunicodechar{ω}{\ensuremath{\omega}}
\newunicodechar{Σ}{\ensuremath{\Sigma}}
% majuscules produites par \MakeUppercase dans les titres (α-aminés -> Α)
\newunicodechar{Α}{\ensuremath{\alpha}}
\newunicodechar{Β}{\ensuremath{\beta}}
\newunicodechar{Γ}{\ensuremath{\Gamma}}
\newunicodechar{Δ}{\ensuremath{\Delta}}
\newunicodechar{Ε}{\ensuremath{\varepsilon}}
\newunicodechar{Θ}{\ensuremath{\Theta}}
\newunicodechar{Λ}{\ensuremath{\Lambda}}
\newunicodechar{Μ}{\ensuremath{\mu}}
\newunicodechar{Τ}{\ensuremath{\tau}}
\newunicodechar{Φ}{\ensuremath{\Phi}}
\newunicodechar{Ψ}{\ensuremath{\Psi}}
\newunicodechar{Ω}{\ensuremath{\Omega}}
\newunicodechar{↦}{\ensuremath{\mapsto}}
\newunicodechar{∈}{\ensuremath{\in}}
\newunicodechar{∉}{\ensuremath{\notin}}
\newunicodechar{∧}{\ensuremath{\wedge}}
\newunicodechar{⇔}{\ensuremath{\Leftrightarrow}}
\newunicodechar{⟺}{\ensuremath{\Longleftrightarrow}}
\newunicodechar{⇒}{\ensuremath{\Rightarrow}}
\newunicodechar{⟹}{\ensuremath{\Longrightarrow}}
\newunicodechar{∘}{\ensuremath{\circ}}
\newunicodechar{∪}{\ensuremath{\cup}}
\newunicodechar{∩}{\ensuremath{\cap}}
\newunicodechar{≡}{\ensuremath{\equiv}}
\newunicodechar{⊂}{\ensuremath{\subset}}
\newunicodechar{⊥}{\ensuremath{\perp}}
\newunicodechar{∃}{\ensuremath{\exists}}
\newunicodechar{∀}{\ensuremath{\forall}}
\newunicodechar{∛}{\ensuremath{\sqrt[3]{\,}}}
\newunicodechar{∜}{\ensuremath{\sqrt[4]{\,}}}
\newunicodechar{⃗}{\ensuremath{{}^{\to}}}
\newunicodechar{ⁿ}{\ifmmode^{n}\else\textsuperscript{\ensuremath{n}}\fi}
\newunicodechar{ˣ}{\ifmmode^{x}\else\textsuperscript{\ensuremath{x}}\fi}
\newunicodechar{ᵇ}{\ifmmode^{b}\else\textsuperscript{\ensuremath{b}}\fi}
\newunicodechar{ʳ}{\ifmmode^{r}\else\textsuperscript{\ensuremath{r}}\fi}
\newunicodechar{ᵘ}{\ifmmode^{u}\else\textsuperscript{\ensuremath{u}}\fi}
\newunicodechar{ʸ}{\ifmmode^{y}\else\textsuperscript{\ensuremath{y}}\fi}
\newunicodechar{ᵃ}{\ifmmode^{a}\else\textsuperscript{\ensuremath{a}}\fi}
\newunicodechar{ᵉ}{\ifmmode^{e}\else\textsuperscript{\ensuremath{e}}\fi}
\newunicodechar{ᵏ}{\ifmmode^{k}\else\textsuperscript{\ensuremath{k}}\fi}
\newunicodechar{ᵖ}{\ifmmode^{p}\else\textsuperscript{\ensuremath{p}}\fi}
\newunicodechar{ᴺ}{\ifmmode^{N}\else\textsuperscript{\ensuremath{N}}\fi}
\newunicodechar{ᶜ}{\ifmmode^{c}\else\textsuperscript{\ensuremath{c}}\fi}
\newunicodechar{ʰ}{\ifmmode^{h}\else\textsuperscript{\ensuremath{h}}\fi}
\newunicodechar{ᵠ}{\ifmmode^{q}\else\textsuperscript{\ensuremath{q}}\fi}
\newunicodechar{ₙ}{\ifmmode_{n}\else\textsubscript{\ensuremath{n}}\fi}
\newunicodechar{ₐ}{\ifmmode_{a}\else\textsubscript{\ensuremath{a}}\fi}
\newunicodechar{ᵢ}{\ifmmode_{i}\else\textsubscript{\ensuremath{i}}\fi}
\newunicodechar{ᵣ}{\ifmmode_{r}\else\textsubscript{\ensuremath{r}}\fi}
\newunicodechar{ₖ}{\ifmmode_{k}\else\textsubscript{\ensuremath{k}}\fi}
\newunicodechar{ᵦ}{\ifmmode_{\beta}\else\textsubscript{\ensuremath{\beta}}\fi}
\newunicodechar{ₓ}{\ifmmode_{x}\else\textsubscript{\ensuremath{x}}\fi}
\newfontfamily\popdisplay{ArchivoBlack-Regular.ttf}[Path=../../../fonts/]
\newfontfamily\poplogo{SpaceGrotesk-Bold.ttf}[Path=../../../fonts/]
\newfontfamily\popsemi{PlusJakartaSans-SemiBold.ttf}[Path=../../../fonts/]
\newfontfamily\popxbold{PlusJakartaSans-ExtraBold.ttf}[Path=../../../fonts/]
\newfontfamily\popxboldls{PlusJakartaSans-ExtraBold.ttf}[Path=../../../fonts/,LetterSpace=3.0]

\setlist[enumerate]{leftmargin=1.7em,itemsep=5pt,topsep=4pt}
\setlist[itemize]{leftmargin=1.7em,itemsep=4pt,topsep=4pt,label={\color{poporange}\textbullet}}
\setlength{\parindent}{0pt}
\setlength{\parskip}{5pt}
\renewcommand{\arraystretch}{1.25}

% pgfplots : style Pop par défaut
\pgfplotsset{
  every axis/.append style={axis line style={popink,line width=1pt},tick style={popink},
    grid style={popline,line width=0.5pt},label style={font=\small},tick label style={font=\footnotesize}},
  popcurve/.style={poporange,line width=1.8pt,no markers,smooth},
}
% Même style disponible dans un \draw TikZ simple (hors pgfplots)
\tikzset{popcurve/.style={poporange,line width=1.8pt}}
\tikzset{popgrid/.style={popline,very thin}}
\colorlet{popcurve}{poporange} % le nom du style est parfois utilisé comme couleur

% ============================================================
% Pill / squiggle
% ============================================================
\newcommand{\poppill}[3]{%
  \tikz[baseline=-0.55ex]\node[draw=popink,line width=1.2pt,rounded corners=2.6pt,%
    fill=#2,inner xsep=6.5pt,inner ysep=3pt]{%
    \popxboldls\fontsize{7}{7}\selectfont\color{#3}\MakeUppercase{#1}};%
}
\newcommand{\popsquiggle}[1]{%
  \tikz[baseline=-0.4ex]{
    \draw[#1,line width=2.6pt,line cap=round]
      plot [smooth] coordinates {(0,0.09) (0.34,0.01) (0.68,0.21) (1.02,0.01) (1.36,0.10)};
  }%
}
% Mot-clé surligné (marqueur orange sous le texte)
\newcommand{\cle}[1]{{\popxbold\color{popdark}#1}}

% ============================================================
% En-tête / pied
% ============================================================
\pagestyle{fancy}
\fancyhf{}
\renewcommand{\headrulewidth}{0pt}
\renewcommand{\footrulewidth}{0pt}
\lhead{\raisebox{-0.15ex}{\poplogo\fontsize{13}{13}\selectfont\color{popink}OnBuch\color{poporange}+}}
\rhead{\poppill{\DOCMATIERE\ \textbullet\ \DOCNIVEAU\ \textbullet\ Leçon \DOCLECON}{white}{popink}}
\lfoot{\popxbold\fontsize{7.6}{7.6}\selectfont\color{popmuted}\MakeUppercase{Cours signé OnBuch+}\ \textbullet\ {\popsemi\DOCTITRE}}
\rfoot{\tikz[baseline=-0.6ex]\node[rounded corners=8.5pt,fill=popink,minimum height=17pt,minimum width=17pt,inner xsep=5pt,inner sep=0]{\popxbold\fontsize{8}{8}\selectfont\color{popcream}\thepage\,/\,\pageref*{LastPage}};}

% ============================================================
% Titres
% ============================================================
\newcommand{\chaptitre}[2]{%
  \begin{tcolorbox}[enhanced,colback=poporange,colframe=popink,boxrule=2.6pt,arc=11pt,
      drop shadow=popink,left=15pt,right=15pt,top=12pt,bottom=11pt,before skip=3pt,after skip=18pt]
    \poppill{#2}{popink}{popcream}\par\vspace{9pt}
    {\raggedright\hyphenpenalty=10000\exhyphenpenalty=10000\popdisplay\fontsize{22}{24}\selectfont\color{popcream}\MakeUppercase{#1}\par}
  \end{tcolorbox}
}
% Section numérotée : pastille orange avec le numéro + titre Archivo
\newcounter{popsec}
\newcommand{\coursec}[1]{%
  \stepcounter{popsec}\setcounter{popsub}{0}%
  \par\vspace{16pt}\noindent
  \begin{minipage}{\linewidth}
  \tikz[baseline=-0.7ex]\node[draw=popink,line width=1.6pt,fill=poporange,rounded corners=4pt,minimum size=22pt,inner sep=0]{\popdisplay\fontsize{12}{12}\selectfont\color{popcream}\thepopsec};%
  \hspace{8pt}{\popdisplay\fontsize{15}{17}\selectfont\color{popink}\MakeUppercase{#1}}\par
  \vspace{4pt}\noindent{\color{poporange}\rule{36pt}{3.6pt}}
  \end{minipage}\par\nopagebreak\vspace{8pt}
}
\newcounter{popsub}
\newcommand{\courssub}[1]{%
  \stepcounter{popsub}%
  \par\vspace{9pt}\noindent{\popxbold\fontsize{11.5}{13}\selectfont\color{popink}{\color{poporange}\thepopsec.\thepopsub}\hspace{6pt}#1}\par\nopagebreak\vspace{4pt}
}

% ============================================================
% Boîtes pédagogiques — même recette : fond blanc, bordure épaisse colorée,
% ombre dure encre, badge plein. Titre optionnel [#1].
% ============================================================
\tcbset{popbase/.style={breakable,enhanced,colback=white,boxrule=2.3pt,arc=9pt,
  shadow={3pt}{-3pt}{0pt}{popink},left=14pt,right=14pt,top=11pt,bottom=11pt,before skip=11pt,after skip=15pt,
  fontupper=\popsemi\fontsize{10.6}{14.5}\selectfont\color{popink},
  fontlower=\popsemi\fontsize{10.6}{14.5}\selectfont\color{popink}}}
% Coche dessinée (le glyphe ✓ n'existe pas dans les polices du document)
\newcommand{\popcheck}[1]{\tikz[baseline=-0.5ex]\draw[#1,line width=1.6pt,line cap=round,line join=round](-0.13,0.0)--(-0.04,-0.09)--(0.13,0.1);}
\providecommand{\coche}{\popcheck{popgreen}}
\providecommand{\resultat}[1]{\par\smallskip\noindent\fcolorbox{popgreen}{popgreenL}{\parbox{\dimexpr\linewidth-2\fboxsep-2\fboxrule\relax}{\textbf{Résultat.} #1}}\par\smallskip}
\newcommand{\popboxhead}[3]{\poppill{#1}{#2}{white}\ifx\relax#3\relax\else\hspace{6pt}{\popxbold\fontsize{10.6}{12}\selectfont\color{popink}#3}\fi\par\vspace{7pt}}

\newtcolorbox{aretenir}[1][]{popbase,colframe=popgreen,before upper={\popboxhead{À retenir}{popgreen}{#1}}}
\newtcolorbox{exemplebox}[1][]{popbase,colframe=poporange,before upper={\popboxhead{Exemple}{poporange}{#1}}}
\newtcolorbox{definition}[1][]{popbase,colframe=popblue,before upper={\popboxhead{Définition}{popblue}{#1}}}
\newtcolorbox{propriete}[1][]{popbase,colframe=poppurple,before upper={\popboxhead{Propriété}{poppurple}{#1}}}
\newtcolorbox{methode}[1][]{popbase,colframe=popgold,before upper={\popboxhead{Méthode}{popgold}{#1}}}
\newtcolorbox{attention}[1][]{popbase,colframe=poppink,before upper={\popboxhead{Attention}{poppink}{#1}}}
\newtcolorbox{experience}[1][]{popbase,colframe=popblue,colback=popblueL,before upper={\popboxhead{Expérience}{popblue}{#1}}}
% « Le savais-tu ? » : carte violette pleine (contexte, culture, Cameroun)
\newtcolorbox{savaistu}[1][]{popbase,colback=poppurple,colframe=popink,
  fontupper=\popsemi\fontsize{10.4}{14.2}\selectfont\color{white},
  before upper={\poppill{Le savais-tu\,?}{white}{poppurple}\ifx\relax#1\relax\else\hspace{6pt}{\popxbold\color{white}#1}\fi\par\vspace{7pt}}}
% Exercice résolu : énoncé en haut, \tcblower puis la solution sur fond crème
\newcounter{popexo}\newcounter{popexr}
\newtcolorbox{exoresolu}[1][]{popbase,colframe=poporange,colbacklower=popcream,
  segmentation style={popink,line width=1.2pt,dashed},
  before upper={\stepcounter{popexr}\popboxhead{Exercice résolu \thepopexr}{poporange}{#1}},
  before lower={\poppill{Solution}{popink}{popcream}\par\vspace{6pt}}}
% Exercice (non résolu en place) — la correction est dans la partie Corrigés
\newtcolorbox{exercice}[1][]{popbase,colframe=popink,
  before upper={\stepcounter{popexo}\popboxhead{Exercice \thepopexo}{popink}{#1}}}
% Corrigé
\newtcolorbox{corrige}[1][]{popbase,colframe=popgreen,colback=popgreenL,
  before upper={\popboxhead{Corrigé}{popgreen}{#1}}}
% Objectifs / prérequis / bilan
\newtcolorbox{objectifs}{popbase,colframe=popink,colback=poporangeL,
  before upper={\popboxhead{Objectifs de la leçon}{popink}{}}}
\newtcolorbox{prerequis}{popbase,colframe=popink,colback=popblueL,
  before upper={\popboxhead{Prérequis}{popblue}{}}}
\newtcolorbox{bilan}{popbase,colframe=popink,colback=white,boxrule=2.8pt,
  before upper={\popboxhead{Fiche bilan}{poporange}{L'essentiel à connaître pour l'examen}}}
% Figure encadrée avec légende
\newtcolorbox{popfigure}[1][]{enhanced,breakable=false,colback=white,colframe=popline,boxrule=1.4pt,arc=8pt,
  left=10pt,right=10pt,top=10pt,bottom=8pt,before skip=10pt,after skip=12pt,halign=center,
  lower separated=false,halign lower=center,
  fontlower=\popsemi\fontsize{9}{11.5}\selectfont\color{popmuted},
  #1}
\newcommand{\legende}[1]{\par\vspace{6pt}{\popsemi\fontsize{9}{11.5}\selectfont\color{popmuted}#1}}

% ============================================================
% Signature OnBuch+
% ============================================================
\newcommand{\poplogoinline}[1]{{\poplogo\fontsize{#1}{#1}\selectfont\color{popink}OnBuch\color{poporange}+}}
\newcommand{\signatureonbuch}{%
  \begin{tcolorbox}[enhanced,colback=popink,colframe=popink,boxrule=2.2pt,arc=10pt,
      drop shadow=poporange,left=14pt,right=14pt,top=10pt,bottom=10pt,before skip=0pt,after skip=0pt]
    \tikz[baseline=-0.7ex]\node[circle,fill=popgreen,draw=popcream,line width=1.3pt,minimum size=22pt,inner sep=0]{\popcheck{white}};%
    \hspace{9pt}\begin{minipage}[c]{0.62\linewidth}
      {\popxbold\fontsize{12}{13}\selectfont\color{popcream}Signé\ \textbullet\ {\poplogo OnBuch\color{poporange}+}}\par\vspace{2pt}
      {\raggedright\popsemi\fontsize{8}{10.5}\selectfont\color{popline}\MakeUppercase{Équipe pédagogique OnBuch+}\\\MakeUppercase{Conforme au programme officiel MINESEC}\par}
    \end{minipage}\hfill
    {\poplogo\fontsize{22}{22}\selectfont\color{popcream}OnBuch\color{poporange}+}
  \end{tcolorbox}%
}

% ============================================================
% Page de garde
% #1 titre · #2 matière · #3 niveau · #4 module · #5 n° leçon · #6 durée ·
% #7 description / objectifs (texte ou itemize)
% ============================================================
\newcommand{\pagegarde}[7]{%
  \thispagestyle{empty}
  \newgeometry{a4paper,margin=1.7cm,top=1.6cm,bottom=1.4cm}%
  \begin{tikzpicture}[remember picture,overlay]
    \fill[poporange!18] ([xshift=-3.2cm,yshift=-3.6cm]current page.north east) circle (4.6cm);
    \fill[poppurple!14] ([xshift=2.4cm,yshift=7.6cm]current page.south west) circle (3.8cm);
  \end{tikzpicture}%
  {\poplogo\fontsize{30}{30}\selectfont\color{popink}OnBuch\color{poporange}+}\hfill
  \raisebox{0.6ex}{\poppill{Cours signé \textbullet\ Premium}{popink}{popcream}}\par
  \vspace{1pt}\popsquiggle{poppurple}\par
  \vspace{8pt}
  \begin{tcolorbox}[enhanced,colback=poporange,colframe=popink,boxrule=2.8pt,arc=13pt,
      drop shadow=popink,left=18pt,right=18pt,top=12pt,bottom=13pt,after skip=10pt]
    \poppill{#2 \textbullet\ #3}{popink}{popcream}\hspace{5pt}\poppill{#4}{white}{popink}\par\vspace{8pt}
    {\popdisplay\fontsize{14}{14}\selectfont\color{popink}LEÇON #5}\par\vspace{4pt}
    {\raggedright\hyphenpenalty=10000\exhyphenpenalty=10000\popdisplay\fontsize{25}{27}\selectfont\color{popcream}\MakeUppercase{#1}\par}
    \vspace{8pt}
    {\popxbold\fontsize{9.5}{9.5}\selectfont\color{popink}\MakeUppercase{Durée conseillée : #6}}
  \end{tcolorbox}
  \begin{tcolorbox}[enhanced,colback=white,colframe=popink,boxrule=2.2pt,arc=10pt,
      drop shadow=popink,left=16pt,right=16pt,top=10pt,bottom=10pt,after skip=8pt]
    \poppill{Dans cette leçon}{poppurple}{white}\par\vspace{5pt}
    \popsemi\fontsize{9.3}{12.6}\selectfont\color{popink}#7
  \end{tcolorbox}
  \noindent\begin{minipage}[t]{0.487\linewidth}
    \begin{tcolorbox}[enhanced,colback=popgreen,colframe=popink,boxrule=2.2pt,arc=10pt,
        drop shadow=popink,left=11pt,right=11pt,top=10pt,bottom=10pt,height=3.1cm,valign=center]
      \tcbox[colback=white,colframe=popink,boxrule=1.3pt,arc=5pt,boxsep=0pt,nobeforeafter,
        left=3pt,right=3pt,top=3pt,bottom=3pt,valign=center]{\qrcode[height=1.9cm]{https://play.google.com/store/apps/details?id=com.luvvix.onbuch&hl=fr}}%
      \hspace{8pt}\begin{minipage}[c]{0.5\linewidth}
        {\popdisplay\fontsize{11}{12.5}\selectfont\color{white}\MakeUppercase{Télécharge OnBuch}}\par\vspace{4pt}
        {\raggedright\popsemi\fontsize{8.4}{11}\selectfont\color{white}Gratuit \textbullet\ Google Play\\Tuteur IA, annales, concours, résultats\par}
      \end{minipage}
    \end{tcolorbox}
  \end{minipage}\hfill
  \begin{minipage}[t]{0.487\linewidth}
    \begin{tcolorbox}[enhanced,colback=poppurple,colframe=popink,boxrule=2.2pt,arc=10pt,
        drop shadow=popink,left=11pt,right=11pt,top=10pt,bottom=10pt,height=3.1cm,valign=center]
      \tikz[baseline=-0.7ex]\node[circle,fill=white,draw=popink,line width=1.3pt,minimum size=1.6cm,inner sep=0]{\popdisplay\fontsize{24}{24}\selectfont\color{poppurple}+};%
      \hspace{8pt}\begin{minipage}[c]{0.6\linewidth}
        {\popdisplay\fontsize{11}{12.5}\selectfont\color{white}\MakeUppercase{Passe à OnBuch+}}\par\vspace{4pt}
        {\raggedright\popsemi\fontsize{8.4}{11}\selectfont\color{white}Tous les cours signés, Tuteur IA illimité, fiches PDF — onglet OnBuch+ de l'app\par}
      \end{minipage}
    \end{tcolorbox}
  \end{minipage}
  \vspace{8pt}
  \popcontenu
  \vfill
  \signatureonbuch
  \restoregeometry
  \newpage
}

% Rangée « Ce cours contient » (page de garde)
\newcommand{\poptile}[3]{% #1 couleur #2 gros texte #3 légende
  \begin{tcolorbox}[enhanced,colback=#1,colframe=popink,boxrule=2pt,arc=9pt,shadow={2.5pt}{-2.5pt}{0pt}{popink},
      left=6pt,right=6pt,top=9pt,bottom=9pt,halign=center,valign=center,height=1.75cm,width=\linewidth,nobeforeafter]
    {\popdisplay\fontsize{12.5}{13}\selectfont\color{white}\MakeUppercase{#2}}\par\vspace{3pt}
    {\centering\popsemi\fontsize{7.8}{9.5}\selectfont\color{white}#3\par}
  \end{tcolorbox}}
\newcommand{\popcontenu}{%
  \par\noindent{\popxboldls\fontsize{7.5}{7.5}\selectfont\color{popmuted}CE COURS SIGNÉ CONTIENT}\par\vspace{7pt}
  \noindent\begin{minipage}[t]{0.235\linewidth}\poptile{popblue}{Cours}{complet, illustré, pas à pas}\end{minipage}\hfill
  \begin{minipage}[t]{0.235\linewidth}\poptile{poporange}{Méthodes}{et exercices résolus}\end{minipage}\hfill
  \begin{minipage}[t]{0.235\linewidth}\poptile{poppink}{Exercices}{type examen + corrigés}\end{minipage}\hfill
  \begin{minipage}[t]{0.235\linewidth}\poptile{popgreen}{Fiche bilan}{l'essentiel à retenir}\end{minipage}\par}

% Sommaire
\newtcolorbox{sommaire}{enhanced,colback=white,colframe=popink,boxrule=2.2pt,arc=10pt,
  drop shadow=popink,left=18pt,right=18pt,top=13pt,bottom=13pt,before skip=6pt,after skip=20pt,
  before upper={\poppill{Sommaire}{poppurple}{white}\par\vspace{9pt}\popsemi\fontsize{10.6}{14.5}\selectfont\color{popink}}}

% Signature de fin de cours — poussée tout en bas de la dernière page
\newcommand{\findecours}{%
  \par\vfill\nopagebreak
  \begin{center}\popsquiggle{poporange}\end{center}\vspace{4pt}
  \signatureonbuch
}
\AtBeginDocument{\def\llbracket{\mathopen{[\mkern-3.5mu[}}\def\rrbracket{\mathclose{]\mkern-3.5mu]}}} % intervalles d'entiers (glyphe ⟦ absent de la police)
% Glyphes absents de Fira Math : redéfinis à partir de glyphes disponibles
\AtBeginDocument{%
  \def\setminus{\mathbin{\backslash}}\let\smallsetminus\setminus
  \def\vdots{\mathord{\vbox{\baselineskip3.2pt\lineskiplimit0pt\kern4pt\hbox{.}\hbox{.}\hbox{.}}}}%
  \def\ddots{\mathinner{\mkern1mu\raise6.5pt\hbox{.}\mkern2mu\raise3.6pt\hbox{.}\mkern2mu\raise0.7pt\hbox{.}\mkern1mu}}%
  \def\triangle{\mathord{\scalebox{0.9}{$\symup{\Delta}$}}}%
  \def\nabla{\mathord{\raisebox{\depth}{\scalebox{1}[-1]{$\symup{\Delta}$}}}}%
  \def\checkmark{\popcheck{popdark}}%
  \def\star{\mathbin{\ast}}\def\bigstar{\mathord{\ast}}%
  \def\diamond{\mathbin{\scalebox{0.6}{$\Diamond$}}}%
  \def\bigcup{\mathop{\mathchoice{\raisebox{-0.3em}{\scalebox{1.7}{$\cup$}}}{\scalebox{1.25}{$\cup$}}{\cup}{\cup}}\displaylimits}%
  \def\bigcap{\mathop{\mathchoice{\raisebox{-0.3em}{\scalebox{1.7}{$\cap$}}}{\scalebox{1.25}{$\cap$}}{\cap}{\cap}}\displaylimits}%
  \def\lesseqqgtr{\mathrel{\lessgtr}}\def\bigoplus{\mathop{\oplus}\displaylimits}%
}
\providecommand{\ssi}{si et seulement si} % abréviation fréquente
\AtBeginDocument{\providecommand{\wideparen}{\overparen}} % arc de cercle

%% FIGFIX-BEGIN v5 — correctifs globaux des figures (docs/onbuch-generation/tools/figfix)
\usepackage{adjustbox}\usepackage{etoolbox}\tcbuselibrary{raster}
% toute figure TikZ/pgfplots plus large que la ligne est réduite à la largeur disponible
\BeforeBeginEnvironment{tikzpicture}{\begin{adjustbox}{max width=\linewidth}}
\AfterEndEnvironment{tikzpicture}{\end{adjustbox}}
% légendes pgfplots lisibles par-dessus les courbes
\pgfplotsset{every axis/.append style={legend style={fill=white,fill opacity=0.92,text opacity=1,draw=black!15,inner sep=3pt}}}
% étiquettes d'arêtes (syntaxe edge["..."]) avec fond blanc
\tikzset{every edge quotes/.append style={fill=white,inner sep=1.2pt}}
%% FIGFIX-END
\begin{document}

\pagegarde{Transformation et recyclage des déchets}{SVT}{Tle D}{Module IV : La Biotechnologie — Valorisation des déchets de l'environnement}{N20}{6 h}{Cette leçon présente les techniques de transformation et de recyclage des déchets papiers et plastiques : production d'énergie à partir du papier, recyclage en serviettes jetables, papier hygiénique et papier-cadeau, recyclage des plastiques en bouteilles et transformation en pavés. Elle montre enfin les avantages de la valorisation des déchets pour l'environnement et pour l'économie camerounaise.
\vspace{4pt}
\begin{multicols}{2}\small
\begin{itemize}
\item À la fin de cette leçon, je sais définir les notions de transformation, de recyclage et de valorisation des déchets.
\item À la fin de cette leçon, je sais expliquer la technique de production de l'énergie à partir du papier.
\item À la fin de cette leçon, je sais expliquer les étapes du recyclage du papier en serviettes jetables, papier hygiénique et papier-cadeau.
\item À la fin de cette leçon, je sais expliquer la technique de recyclage des déchets plastiques en bouteilles recyclées.
\item À la fin de cette leçon, je sais expliquer la transformation des déchets plastiques en pavés de construction.
\item À la fin de cette leçon, je sais interpréter des documents et des résultats expérimentaux portant sur la valorisation des déchets.
\end{itemize}\end{multicols}}

\begin{sommaire}
\begin{multicols}{2}
\begin{enumerate}
\item Généralités : déchets, transformation, recyclage et valorisation
\item Valorisation énergétique du papier : production d'énergie
\item Recyclage du papier : serviettes jetables, papier hygiénique, papier-cadeau
\item Recyclage des déchets plastiques en bouteilles recyclées
\item Transformation des déchets plastiques en pavés
\item Avantages de la valorisation des déchets et sensibilisation de la communauté
\item Activité d'intégration
\item Méthodes et exercices résolus
\item Exercices
\item Corrigés des exercices
\item Fiche bilan
\end{enumerate}
\end{multicols}
\end{sommaire}

\begin{objectifs}
\begin{itemize}
\item À la fin de cette leçon, je sais définir les notions de transformation, de recyclage et de valorisation des déchets.
\item À la fin de cette leçon, je sais expliquer la technique de production de l'énergie à partir du papier.
\item À la fin de cette leçon, je sais expliquer les étapes du recyclage du papier en serviettes jetables, papier hygiénique et papier-cadeau.
\item À la fin de cette leçon, je sais expliquer la technique de recyclage des déchets plastiques en bouteilles recyclées.
\item À la fin de cette leçon, je sais expliquer la transformation des déchets plastiques en pavés de construction.
\item À la fin de cette leçon, je sais interpréter des documents et des résultats expérimentaux portant sur la valorisation des déchets.
\item À la fin de cette leçon, je sais citer les avantages de la valorisation des déchets pour l'environnement et pour l'économie.
\item À la fin de cette leçon, je sais informer et sensibiliser ma communauté sur la nécessité de transformer et de recycler les déchets.
\end{itemize}
\end{objectifs}

\begin{prerequis}
\begin{itemize}
\item Notion de pollution de l'environnement et types de déchets (biodégradables / non biodégradables) vues en classes antérieures
\item Notion de fermentation et de métabolisme microbien (module Biotechnologie, leçons précédentes)
\item Notion de combustion et de dégagement d'énergie (chimie organique, Première)
\item Notion de polymères et de matières plastiques (notions de base)
\item Lecture et interprétation de tableaux de données et de courbes
\end{itemize}
\end{prerequis}

\newpage

\coursec{Généralités : déchets, transformation, recyclage et valorisation}

\textbf{Situation d'accroche.} À Yaoundé comme à Douala, les marchés et les quartiers débordent chaque jour de sachets plastiques, de bouteilles vides et de papiers usagés qui bouchent les caniveaux et provoquent des inondations. Pourtant, dans certains quartiers, des jeunes transforment ces mêmes déchets en pavés de construction, en papier recyclé ou en source d'énergie. Comment des déchets qui polluent notre environnement peuvent-ils devenir des richesses ? Quelles techniques permettent de valoriser le papier et le plastique ? C'est ce que nous allons découvrir dans cette leçon.

\textbf{Problématique.} Face à l'accumulation croissante des déchets dans nos villes, comment peut-on les transformer et les recycler pour en faire des produits utiles, tout en protégeant l'environnement et en créant de la richesse ?

\courssub{Qu'est-ce qu'un déchet ?}

\textbf{Observation.} Après un repas de famille, on jette les épluchures d'igname, les os, les feuilles de banane et les sachets plastiques. Un an plus tard, dans la fosse à ordures, les épluchures et les feuilles ont disparu, mais les sachets plastiques sont toujours là, presque intacts. Cette simple observation de la vie quotidienne soulève deux questions : qu'est-ce qu'un déchet, et pourquoi tous les déchets ne disparaissent-ils pas au même rythme ?

\begin{definition}[Déchet]
Un \cle{déchet} est tout résidu d'un processus de production, de transformation ou d'utilisation, toute substance, matière ou produit dont le détenteur se débarrasse ou dont il a l'intention ou l'obligation de se débarrasser, car il n'a plus d'utilité immédiate pour lui.
\end{definition}

Un déchet n'est donc pas défini par sa nature chimique, mais par le fait qu'on s'en débarrasse. Le même objet peut être un bien pour une personne et un déchet pour une autre : c'est toute la logique de la valorisation, qui consiste à redonner de la valeur à ce qui a été jeté.

\courssub{Déchets biodégradables et déchets non biodégradables}

\textbf{Démarche d'investigation.} Pour comprendre la différence de comportement des déchets dans la nature, réalisons une expérience simple réalisable dans la cour du lycée.

\begin{experience}[Dégradation comparée d'un papier et d'un sachet plastique]
\textbf{Matériel :} deux carrés de terre nue de même surface, un morceau de papier journal, un morceau de sachet plastique, de l'eau, un marqueur.

\textbf{Protocole :} on enterre à la même profondeur (environ \SI{10}{cm}) le papier dans le premier carré et le plastique dans le second. On arrose régulièrement les deux carrés de la même quantité d'eau et on observe chaque mois pendant six mois.

\textbf{Observations :} au bout de quelques semaines, le papier se ramollit, se fragmente puis disparaît progressivement ; il est attaqué par les micro-organismes du sol (bactéries, champignons). Au bout de six mois, le sachet plastique est toujours intact : ni sa forme ni sa masse n'ont sensiblement changé.

\textbf{Interprétation :} le papier, constitué de cellulose (une molécule organique d'origine végétale), sert de nourriture aux décomposeurs du sol qui le transforment en matière minérale et humus. Le plastique, polymère de synthèse, n'est reconnu par aucune enzyme microbienne : il persiste dans l'environnement.

\textbf{Conclusion :} on distingue les déchets \cle{biodégradables}, décomposés par les micro-organismes en un temps relativement court, et les déchets \cle{non biodégradables}, qui persistent très longtemps dans la nature.
\end{experience}

\begin{definition}[Déchets biodégradables et non biodégradables]
Un déchet \cle{biodégradable} est un déchet qui peut être décomposé par les micro-organismes du milieu (bactéries, champignons) en substances simples (eau, gaz carbonique, matière minérale, humus). Exemples : papiers, restes alimentaires, feuilles mortes, fientes.

Un déchet \cle{non biodégradable} est un déchet qui ne peut pas être décomposé par les micro-organismes, ou seulement sur des durées extrêmement longues. Exemples : plastiques, verre, métaux.
\end{definition}

\begin{center}
\rowcolors{2}{popcream}{white}
\begin{tabularx}{\linewidth}{|l|X|X|}
\hline
\rowcolor{popblueL} \textbf{Critère} & \textbf{Déchets biodégradables} & \textbf{Déchets non biodégradables} \\
\hline
Exemples & Papiers, restes alimentaires, épluchures, feuilles & Sachets et bouteilles plastiques, verre, canettes \\
\hline
Origine & Matière organique d'origine végétale ou animale & Matériaux de synthèse ou minéraux \\
\hline
Décomposition & Par les micro-organismes (bactéries, champignons) & Aucune décomposition microbienne significative \\
\hline
Durée de dégradation & Quelques semaines à quelques mois (papier : 2 à 6 mois) & Plusieurs siècles (sachet plastique : plus de 400 ans) \\
\hline
Devenir naturel & Humus, compost, matière minérale & Accumulation et pollution durable \\
\hline
\end{tabularx}
\end{center}

\begin{savaistu}[Des durées de dégradation très contrastées]
Un sachet plastique jeté dans la nature peut mettre plus de \textbf{400 ans} à se dégrader, alors qu'un papier se dégrade en \textbf{quelques mois} seulement (2 à 6 mois selon l'humidité du sol). Une bouteille en plastique peut persister 100 à 1000 ans. C'est pourquoi l'accumulation des plastiques constitue un problème environnemental majeur : chaque sachet jeté aujourd'hui à Yaoundé existera encore au temps de tes arrière-arrière-petits-enfants !
\end{savaistu}

\courssub{Problématique locale : les déchets dans les villes camerounaises}

\begin{exemplebox}[La production de déchets à Douala]
La ville de Douala produit plusieurs centaines de tonnes de déchets par jour (on estime la production à plus de \SI{1000}{tonnes/jour} pour l'agglomération). Or, une faible part seulement de ces déchets est collectée, et une part encore plus faible est réellement valorisée : l'essentiel finit dans des décharges sauvages, les caniveaux ou la nature.
\end{exemplebox}

L'accumulation des déchets non maîtrisée a des conséquences graves et directement observables dans nos villes :

\begin{itemize}
\item \textbf{Inondations :} les sachets et bouteilles plastiques bouchent les caniveaux et les évacuations d'eau ; à la saison des pluies, l'eau ne s'écoule plus et envahit les quartiers (cas fréquents à Douala, ville basse et très pluvieuse).
\item \textbf{Maladies :} les déchets stagnants abritent les moustiques (vecteurs du paludisme) et contaminent l'eau de boisson, favorisant le choléra, la typhoïde et les diarrhées, surtout chez les enfants.
\item \textbf{Pollution des sols et des eaux :} les plastiques fragmentés (microplastiques) et les substances chimiques des décharges s'infiltrent dans le sol et les nappes phréatiques, appauvrissent les terres agricoles et contaminent les puits et les cours d'eau.
\item \textbf{Pollution de l'air :} la combustion sauvage des ordures dégage des fumées toxiques.
\end{itemize}

Face à ce constat, une question s'impose : au lieu de subir les déchets, ne peut-on pas en faire une ressource ? C'est l'objet de la transformation, du recyclage et, plus largement, de la valorisation.

\courssub{Transformation, recyclage et valorisation : trois notions à ne pas confondre}

\textbf{Intuition.} Une bouteille plastique vide peut connaître deux destins différents : être fondue pour fabriquer une \emph{nouvelle bouteille} (on refabrique un produit similaire), ou être broyée puis mélangée au sable pour fabriquer un \emph{pavé de construction} (on obtient un produit de nature totalement différente). Dans le premier cas, on parle de recyclage ; dans le second, de transformation. Dans les deux cas, le déchet a gagné une nouvelle valeur : c'est la valorisation.

\begin{definition}[Transformation]
La \cle{transformation} d'un déchet est l'ensemble des opérations qui modifient la \textbf{nature} ou la \textbf{forme} de ce déchet pour obtenir un produit nouveau, de nature différente du produit d'origine. Exemple : la transformation des déchets plastiques en pavés de construction.
\end{definition}

\begin{definition}[Recyclage]
Le \cle{recyclage} est la réintroduction d'un déchet dans un nouveau cycle de production, en remplacement total ou partiel de la matière première, pour fabriquer un produit similaire ou équivalent. Exemples : le recyclage du papier usagé en papier hygiénique ou en papier-cadeau ; le recyclage des bouteilles plastiques en nouvelles bouteilles.
\end{definition}

\begin{definition}[Valorisation]
La \cle{valorisation} est l'ensemble des opérations qui donnent une nouvelle valeur (matière, énergie, usage) à un déchet. Elle comprend le recyclage (valorisation matière), la transformation en nouveaux produits et la valorisation énergétique (production d'énergie par combustion ou méthanisation).
\end{definition}

\begin{attention}[Ne pas confondre recycler et transformer]
Erreur fréquente à l'examen : employer « recycler » pour n'importe quelle opération sur un déchet. Retiens la distinction :
\begin{itemize}
\item \textbf{Recycler} = refabriquer un produit \emph{similaire} à partir du déchet (papier usagé $\rightarrow$ papier recyclé ; bouteille plastique $\rightarrow$ nouvelle bouteille).
\item \textbf{Transformer} = obtenir un produit de \emph{nature différente} (plastique $\rightarrow$ pavé ; papier $\rightarrow$ énergie).
\end{itemize}
Dans les deux cas, il y a \emph{valorisation}, terme générique qui englobe les deux. Une rédaction qui confond ces termes est sanctionnée.
\end{attention}

\courssub{Principe général de la valorisation des déchets}

Quel que soit le déchet considéré, la valorisation suit toujours la même logique en quatre étapes : le \cle{tri} (séparer les déchets selon leur nature), la \cle{collecte} (rassembler les déchets triés), le \cle{traitement} (opérations techniques : lavage, broyage, fonte, combustion, etc.) et l'obtention d'un \cle{produit valorisé} (matière première recyclée, nouveau produit ou énergie).

\begin{popfigure}
\begin{center}
\begin{tikzpicture}[
box/.style={rectangle, rounded corners=3pt, draw=popblue, fill=popblueL, thick, minimum height=1cm, minimum width=2.2cm, align=center, font=\small\bfseries},
arr/.style={-{Stealth[length=3mm]}, very thick, poporange},
node distance=0.9cm]
\node[box, fill=popgoldL, draw=popgold] (d) {Déchet};
\node[box, right=of d] (t) {Tri};
\node[box, right=of t] (c) {Collecte};
\node[box, right=of c] (tr) {Traitement};
\node[box, fill=popgreenL, draw=popgreen, right=of tr, align=center] (p){Produit\\valorisé};
\draw[arr] (d) -- (t);
\draw[arr] (t) -- (c);
\draw[arr] (c) -- (tr);
\draw[arr] (tr) -- (p);
\end{tikzpicture}
\end{center}
\legende{Chaîne générale de la valorisation d'un déchet. Le tri sépare les déchets par nature (papier, plastique, organique) ; la collecte les achemine vers l'unité de traitement ; le traitement (lavage, broyage, fonte, combustion) aboutit à un produit utile : matière recyclée, objet transformé ou énergie.}
\end{popfigure}

\begin{methode}[Analyser une filière de valorisation]
Pour décrire correctement une filière de valorisation (au Bac comme dans la vie), procède toujours en quatre étapes :
\begin{enumerate}
\item Identifier le \textbf{déchet de départ} et sa nature (biodégradable ou non).
\item Décrire le \textbf{tri et la collecte} : comment le déchet est séparé et rassemblé.
\item Expliquer le \textbf{traitement} : les opérations techniques subies par le déchet.
\item Nommer le \textbf{produit obtenu} et préciser s'il s'agit d'un recyclage (produit similaire) ou d'une transformation (produit différent).
\end{enumerate}
\end{methode}

\courssub{Les deux grandes voies étudiées dans cette leçon}

Dans la suite de cette leçon, nous appliquerons ce schéma général aux deux familles de déchets les plus abondantes dans nos villes :

\begin{itemize}
\item \textbf{La valorisation des déchets papiers} : d'une part la production d'énergie (valorisation énergétique), d'autre part le recyclage en serviettes jetables, en papier hygiénique et en papier-cadeau (valorisation matière).
\item \textbf{La valorisation des déchets plastiques} : le recyclage des bouteilles plastiques en bouteilles recyclées et leur transformation en pavés de construction.
\end{itemize}

\begin{exoresolu}[Identifier les opérations de valorisation]
Dans un quartier de Yaoundé, une association de jeunes réalise les activités suivantes : (1) elle collecte les vieux cahiers et journaux pour fabriquer du papier-cadeau ; (2) elle broie les bouteilles plastiques et les mélange au sable pour fabriquer des pavés ; (3) elle brûle les résidus de papier non recyclables pour chauffer un four à briques. Pour chaque activité, indique s'il s'agit d'un recyclage, d'une transformation ou d'une valorisation énergétique, en justifiant.
\tcblower
\textbf{Solution.}
\begin{enumerate}
\item Vieux papiers $\rightarrow$ papier-cadeau : il s'agit d'un \textbf{recyclage}, car le déchet (papier) est réintroduit dans un nouveau cycle de production pour fabriquer un produit de même nature (du papier).
\item Bouteilles plastiques $\rightarrow$ pavés : il s'agit d'une \textbf{transformation}, car on obtient un produit de nature différente du produit d'origine (un matériau de construction à partir d'un emballage).
\item Résidus de papier $\rightarrow$ chaleur pour le four : il s'agit d'une \textbf{valorisation énergétique}, car le déchet sert à produire de l'énergie.
\end{enumerate}
Les trois activités relèvent de la \textbf{valorisation}, terme générique qui désigne toute opération redonnant de la valeur à un déchet.
\end{exoresolu}

\begin{aretenir}[L'essentiel de la section]
\begin{itemize}
\item Un \cle{déchet} est toute substance ou objet dont le détenteur se débarrasse.
\item Les déchets \cle{biodégradables} (papiers, restes alimentaires) sont décomposés par les micro-organismes en quelques mois ; les déchets \cle{non biodégradables} (plastiques) persistent des siècles dans la nature.
\item L'accumulation des déchets provoque inondations, maladies et pollution des sols et des eaux.
\item \cle{Recycler} = refabriquer un produit similaire ; \cle{transformer} = obtenir un produit de nature différente ; \cle{valoriser} = redonner une valeur (matière ou énergie) au déchet.
\item Toute valorisation suit la chaîne : Déchet $\rightarrow$ Tri $\rightarrow$ Collecte $\rightarrow$ Traitement $\rightarrow$ Produit valorisé.
\end{itemize}
\end{aretenir}

\begin{exercice}[Pour t'entraîner]
\begin{enumerate}
\item Définis : déchet, biodégradable, recyclage, valorisation.
\item Classe les déchets suivants en biodégradables et non biodégradables : épluchures de manioc, sachet plastique, journal usagé, bouteille en verre, feuilles mortes, canette.
\item Explique pourquoi les sachets plastiques jetés dans les caniveaux aggravent les inondations à Douala.
\end{enumerate}
\end{exercice}

\begin{corrige}[Exercice « Pour t'entraîner »]
\begin{enumerate}
\item Voir les définitions du cours : déchet (substance ou objet dont on se débarrasse), biodégradable (décomposable par les micro-organismes), recyclage (réintroduction du déchet dans un nouveau cycle de production d'un produit similaire), valorisation (ensemble des opérations redonnant une valeur au déchet).
\item Biodégradables : épluchures de manioc, journal usagé, feuilles mortes. Non biodégradables : sachet plastique, bouteille en verre, canette.
\item Les sachets plastiques, non biodégradables, s'accumulent dans les caniveaux et bouchent les évacuations d'eau. À la saison des pluies, l'eau ne peut plus s'écouler : elle déborde et inonde les quartiers, d'autant que Douala est une ville basse et très arrosée.
\end{enumerate}
\end{corrige}

\coursec{Valorisation énergétique du papier : production d'énergie}

Dans la section précédente, nous avons défini les notions de \cle{déchet}, de \cle{transformation}, de \cle{recyclage} et de \cle{valorisation}. Nous avons vu qu'un déchet n'est pas une fatalité : c'est une \emph{ressource mal placée}. Nous allons maintenant étudier la première voie concrète de valorisation des déchets papiers : leur utilisation comme \textbf{source d'énergie}. Autour de toi, à Yaoundé, Douala ou Bafoussam, tu vois chaque jour des papiers usagés (cahiers, journaux, cartons d'emballage) abandonnés ou brûlés à l'air libre. Or ce papier contient de l'énergie chimique que l'on peut récupérer de manière contrôlée. Comment ? C'est l'objet de cette section.

\courssub{Constat expérimental : le papier est un combustible}

\begin{experience}[Le papier brûle en dégageant de la chaleur]
\textbf{Matériel} : une feuille de papier propre et sèche, une pince métallique, une allumette, une coupelle métallique, un thermomètre (ou simplement la main placée à distance de sécurité).

\textbf{Protocole} :
\begin{enumerate}
\item Tenir la feuille de papier avec la pince au-dessus de la coupelle.
\item Enflammer un coin de la feuille et observer la combustion jusqu'à son terme.
\item Approcher prudemment la main (ou le thermomètre) au-dessus de la flamme et noter la sensation ou la température.
\item Observer les résidus : cendres grises et fumées.
\end{enumerate}

\textbf{Observations} : le papier s'enflamme facilement ; la combustion dégage une \textbf{chaleur importante} et de la lumière ; il reste des cendres et il se dégage des fumées.

\textbf{Interprétation} : le papier est constitué essentiellement de \cle{cellulose}, une matière organique (glucide polymère de glucose, de formule brute $(\ce{C6H10O5})_n$) fabriquée par les végétaux grâce à la photosynthèse. La combustion de la cellulose est une réaction chimique \textbf{exothermique} (qui libère de la chaleur) avec le dioxygène de l'air :
\[ \ce{C6H10O5 + 6O2 -> 6CO2 + 5H2O} + \text{énergie (chaleur)} \]

\textbf{Conclusion} : le papier est un \cle{combustible}. L'énergie chimique stockée dans la matière organique (la cellulose) est libérée sous forme de chaleur lors de la combustion. Un déchet papier peut donc servir de \textbf{source d'énergie}.
\end{experience}

\begin{attention}[Combustion contrôlée ou non : toute la différence]
Brûler les déchets à l'air libre dans les quartiers, comme on le voit trop souvent dans nos villes, \textbf{n'est pas une valorisation} : la combustion y est incomplète, elle dégage des fumées toxiques (monoxyde de carbone \ce{CO}, particules fines, composés chlorés si des plastiques sont mélangés) qui polluent l'air et provoquent des maladies respiratoires. La valorisation énergétique exige une \textbf{combustion contrôlée} dans des installations adaptées, avec récupération de la chaleur et contrôle des fumées.
\end{attention}

\courssub{Principe de la valorisation énergétique du papier}

\begin{definition}[Valorisation énergétique]
La \cle{valorisation énergétique} est l'utilisation d'un déchet comme \textbf{source d'énergie} : au lieu d'être simplement éliminé, le déchet est brûlé ou transformé en combustible afin de récupérer l'énergie chimique contenue dans sa matière organique, sous forme de \textbf{chaleur} (chauffage, cuisson, production de vapeur) ou d'\textbf{électricité}.
\end{definition}

Pour le papier, deux voies principales existent :
\begin{itemize}
\item la \textbf{combustion directe} (incinération contrôlée) : le papier sec est brûlé dans un foyer ou un incinérateur équipé d'un système de récupération de chaleur ;
\item la \textbf{transformation préalable en combustible} : le papier est séché, broyé puis compressé en \cle{briquettes} ou en \cle{granulés} de papier, combustibles solides faciles à stocker, à transporter et à utiliser dans les foyers de cuisson ou de chauffage.
\end{itemize}

\begin{aretenir}[Rôle de la cellulose]
La cellulose est la \textbf{molécule énergétique clé} du papier. Synthétisée par les plantes à partir du \ce{CO2} de l'air et de l'eau grâce à l'énergie solaire (photosynthèse), elle constitue un \emph{réservoir d'énergie chimique}. Sa combustion libère cette énergie sous forme de chaleur. C'est pourquoi le papier, riche en cellulose, possède un \textbf{pouvoir calorifique} intéressant, comparable à celui d'un bois de chauffage de qualité moyenne lorsqu'il est bien sec et compressé.
\end{aretenir}

\courssub{Les étapes de la production d'énergie à partir du papier}

La production d'énergie à partir des déchets papiers suit une \textbf{chaîne d'étapes ordonnées}, qu'il faut savoir décrire et justifier.

\begin{enumerate}
\item \textbf{Collecte et tri du papier.} Les papiers usagés (cahiers, journaux, cartons, papiers de bureau) sont ramassés puis \textbf{triés} : on élimine les matières étrangères (plastiques, métaux agrafes, verre) et surtout les papiers souillés de produits dangereux (huiles, peintures, produits chimiques). Le tri garantit une combustion propre et un combustible de qualité.
\item \textbf{Séchage.} Le papier humide brûle mal : une partie de la chaleur serait gaspillée à évaporer l'eau. On le fait donc sécher (au soleil, opération gratuite sous nos climats, ou en séchoir) jusqu'à obtenir un papier bien sec.
\item \textbf{Broyage.} Le papier sec est déchiqueté en petits fragments. Le broyage augmente la surface de contact avec l'air et facilite l'étape suivante.
\item \textbf{Compression éventuelle en briquettes.} Les fragments de papier, éventuellement humidifiés en pâte puis pressés, sont \textbf{compressés} en briquettes denses (ou en granulés) à l'aide d'une presse manuelle ou mécanique, puis séchés à nouveau. La compression augmente la densité du combustible : les briquettes brûlent plus longtemps et plus régulièrement que le papier en vrac.
\item \textbf{Combustion.} Les briquettes (ou le papier sec broyé) sont brûlées dans un foyer ou un incinérateur contrôlé, avec un apport d'air suffisant pour obtenir une combustion aussi complète que possible.
\item \textbf{Utilisation de la chaleur.} La chaleur dégagée est récupérée et utilisée : \textbf{chauffage} et cuisson (foyers domestiques, fours de briqueteries ou de boulangeries), \textbf{production de vapeur} (chaufferie), et éventuellement \textbf{production d'électricité} (la vapeur fait tourner une turbine couplée à un alternateur) dans les grandes unités d'incinération.
\end{enumerate}

\begin{popfigure}
\begin{center}
\begin{tikzpicture}[
  box/.style={draw=popblue, fill=popblueL, rounded corners=3pt, minimum height=1.1cm, minimum width=2.5cm, align=center, font=\small, text=popdark},
  ebox/.style={draw=poporange, fill=poporangeL, rounded corners=3pt, minimum height=1.1cm, minimum width=2.6cm, align=center, font=\small\bfseries, text=popdark},
  fl/.style={-{Stealth[length=3mm]}, thick, popink}
]
\node[box, align=center] (a) at (0,0){Papier usagé\\(collecte)};
\node[box, align=center] (b) at (3.4,0){Tri\\(élimination des\\corps étrangers)};
\node[box, align=center] (c) at (6.8,0){Séchage\\et broyage};
\node[box, align=center] (d) at (10.2,0){Compression\\en briquettes};
\node[box, align=center] (e) at (3.4,-2.4){Combustion\\contrôlée};
\node[ebox, align=center] (f) at (6.8,-2.4){Chaleur\\(chauffage, vapeur)};
\node[ebox, align=center] (g) at (10.2,-2.4){Électricité\\(turbine + alternateur)};
\draw[fl] (a) -- (b);
\draw[fl] (b) -- (c);
\draw[fl] (c) -- (d);
\draw[fl] (d.south) |- ++(0,-0.6) -| (e.north);
\draw[fl] (e) -- (f);
\draw[fl] (f) -- (g);
\end{tikzpicture}
\end{center}
\legende{Chaîne de production d'énergie à partir des déchets papiers. Le papier usagé est collecté, trié, séché, broyé puis compressé en briquettes ; leur combustion contrôlée libère de la chaleur, utilisable directement (chauffage, cuisson, vapeur) ou convertie en électricité.}
\end{popfigure}

\courssub{Le dispositif de combustion contrôlée : foyer et récupération de chaleur}

Pour que la combustion du papier soit une véritable valorisation, elle doit se dérouler dans une \textbf{installation adaptée} dont on peut schématiser les éléments essentiels : une entrée de combustible, une arrivée d'air, une zone de combustion, un échangeur qui récupère la chaleur, et une cheminée équipée d'un dispositif de contrôle des fumées.

\begin{popfigure}
\begin{center}
\begin{tikzpicture}[font=\small]
% Corps de l'incinérateur
\draw[thick, popink, fill=popcream] (2,0) rectangle (8,5);
% Zone de combustion
\draw[thick, poporange, fill=poporangeL] (3,0.4) rectangle (7,2.2);
\node[popdark] at (5,1.3) {\textbf{Zone de combustion}};
% Briquettes
\draw[fill=popgold!60, draw=popgold] (3.4,0.5) rectangle (4.1,0.9);
\draw[fill=popgold!60, draw=popgold] (4.3,0.5) rectangle (5.0,0.9);
\draw[fill=popgold!60, draw=popgold] (5.2,0.5) rectangle (5.9,0.9);
% Flamme stylisée
\draw[poporange, thick, decorate, decoration={coil, amplitude=2pt, segment length=6pt}] (4,2.2) -- (4,3.0);
\draw[poporange, thick, decorate, decoration={coil, amplitude=2pt, segment length=6pt}] (5,2.2) -- (5,3.2);
\draw[poporange, thick, decorate, decoration={coil, amplitude=2pt, segment length=6pt}] (6,2.2) -- (6,3.0);
% Serpentin échangeur
\draw[popblue, very thick] (2.6,3.4) -- (7.4,3.4) (2.6,3.8) -- (7.4,3.8);
\draw[popblue, very thick] (2.6,3.4) .. controls (2.3,3.6) .. (2.6,3.8);
\draw[popblue, very thick] (7.4,3.8) .. controls (7.7,3.6) .. (7.4,3.4);
\node[popblue, align=center] at (5,4.3) {Échangeur : eau chauffée $\rightarrow$ vapeur};
% Entrée combustible
\draw[thick, popink, fill=popgreenL] (0.4,1.2) -- (2,1.2) -- (2,2.0) -- (0.4,2.0) -- cycle;
\draw[-{Stealth[length=3mm]}, thick, popgreen] (0.2,1.6) -- (1.2,1.6);
\node[align=center, popdark] at (1.0,2.5) {Entrée du\\combustible};
% Arrivée d'air
\draw[-{Stealth[length=3mm]}, thick, poppurple] (0.2,0.5) -- (2.6,0.5);
\node[popdark] at (1.2,-0.35) {Arrivée d'air (comburant)};
% Cheminée
\draw[thick, popink, fill=popmuted!20] (6.6,5) -- (6.6,6.6) -- (7.4,6.6) -- (7.4,5);
\draw[-{Stealth[length=3mm]}, thick, popmuted] (7,6.6) -- (7,7.3);
\node[align=center, popdark] at (7,7.8) {Fumées contrôlées\\(cheminée + filtre)};
% Sortie chaleur
\draw[-{Stealth[length=3mm]}, very thick, poporange] (7.4,3.6) -- (9.2,3.6);
\node[align=center, popdark] at (10.6,3.6) {Chaleur récupérée\\(chauffage, vapeur,\\électricité)};
% Cendres
\draw[-{Stealth[length=3mm]}, thick, popmuted] (5,0) -- (5,-0.8);
\node[popdark] at (5,-1.2) {Cendres (résidus)};
\end{tikzpicture}
\end{center}
\legende{Schéma fonctionnel d'un foyer d'incinération contrôlée. Le combustible (briquettes de papier) et l'air entrent dans la zone de combustion ; la chaleur dégagée chauffe l'eau d'un échangeur (vapeur pour le chauffage ou une turbine) ; les fumées sont contrôlées avant rejet et les cendres récupérées.}
\end{popfigure}

\courssub{Conditions d'une bonne valorisation énergétique}

\begin{propriete}[Conditions de qualité de la combustion]
Une valorisation énergétique du papier est efficace et propre si :
\begin{itemize}
\item le papier est \textbf{sec} (l'humidité gaspille la chaleur et favorise une combustion incomplète) ;
\item le papier est \textbf{trié}, exempt de matières dangereuses (plastiques chlorés, papiers souillés de produits chimiques, encres spéciales) ;
\item l'apport d'air est \textbf{suffisant} pour une combustion complète (sinon production de monoxyde de carbone \ce{CO}, gaz toxique) ;
\item les \textbf{fumées sont contrôlées} (cheminée, filtration) pour limiter le rejet de particules et de gaz polluants.
\end{itemize}
\end{propriete}

\begin{attention}[Limites et précautions]
La valorisation énergétique du papier présente des \textbf{limites} qu'il ne faut pas occulter dans une copie :
\begin{itemize}
\item la combustion émet des \textbf{fumées et des gaz} (\ce{CO2}, et \ce{CO} si la combustion est incomplète) ainsi que des particules : des installations adaptées et un contrôle des rejets sont indispensables ;
\item le papier mélangé à des plastiques ou souillé de produits chimiques dégage des composés \textbf{toxiques} : d'où l'importance absolue du tri ;
\item la combustion \emph{détruit} la matière : lorsque le papier est propre et recyclable, le \textbf{recyclage en nouveau papier} (section suivante) est souvent préférable, car il préserve la ressource cellulose.
\end{itemize}
\end{attention}

\courssub{Intérêt concret : les briquettes de papier, un substitut au bois de chauffe}

\begin{exemplebox}[Les briquettes de papier contre la déforestation]
Dans de nombreuses villes camerounaises, le bois de chauffe et le charbon de bois restent les principales sources d'énergie domestique, ce qui contribue à la \textbf{déforestation} des savanes et des lisières forestières. Des briquettes fabriquées à partir de papier et carton usagés compressés peuvent \textbf{remplacer partiellement} le bois de chauffe dans les foyers de cuisson. Chaque tonne de papier valorisée en briquettes épargne une quantité équivalente de bois prélevé dans la nature : la valorisation énergétique du papier devient ainsi un \textbf{outil de protection des forêts}, tout en réduisant le volume des déchets urbains et en créant des revenus pour les jeunes et les femmes organisés en coopératives de collecte et de fabrication.
\end{exemplebox}

\begin{methode}[Expliquer une technique de valorisation au Baccalauréat]
Pour expliquer correctement une technique de valorisation d'un déchet, présente toujours, \textbf{dans l'ordre} :
\begin{enumerate}
\item la \textbf{matière première} (nature du déchet, propriété exploitée — ici : papier riche en cellulose, combustible) ;
\item les \textbf{étapes successives} du procédé, chacune justifiée (collecte et tri $\rightarrow$ séchage et broyage $\rightarrow$ compression en briquettes $\rightarrow$ combustion $\rightarrow$ récupération de la chaleur) ;
\item le \textbf{produit obtenu} et son utilisation (chaleur pour le chauffage ou la vapeur, électricité) ;
\item les \textbf{conditions et précautions} (papier sec et trié, combustion contrôlée, contrôle des fumées).
\end{enumerate}
Une réponse qui se contente d'écrire « on brûle le papier pour avoir de l'énergie » est \textbf{insuffisante} : elle ne décrit ni la technique ni ses conditions.
\end{methode}

\begin{exoresolu}[Expliquer la production d'énergie à partir du papier]
\textbf{Énoncé.} Une coopérative de jeunes de la ville de Bafoussam collecte chaque semaine les papiers et cartons usagés des établissements scolaires et des marchés, les transforme en briquettes combustibles vendues aux ménages pour la cuisson.

1. Définir la valorisation énergétique.

2. Expliquer, en décrivant les étapes, la technique de production d'énergie à partir du papier mise en œuvre par cette coopérative.

3. Citer deux conditions nécessaires à une combustion propre et deux avantages de cette initiative pour l'environnement.

\tcblower
\textbf{Solution détaillée.}

\textbf{1.} La valorisation énergétique est l'utilisation d'un déchet comme source d'énergie : l'énergie chimique contenue dans sa matière organique est récupérée, par combustion contrôlée, sous forme de chaleur ou d'électricité.

\textbf{2.} La matière première est le papier usagé, riche en \textbf{cellulose}, matière organique combustible. Les étapes sont :
\begin{itemize}
\item \emph{collecte et tri} : les papiers et cartons sont ramassés et débarrassés des matières étrangères (plastiques, métaux) et des papiers souillés de produits dangereux ;
\item \emph{séchage} : le papier est séché, car l'humidité gaspillerait la chaleur de combustion ;
\item \emph{broyage} : le papier sec est déchiqueté pour faciliter la compression ;
\item \emph{compression en briquettes} : la pâte de papier est pressée puis séchée, ce qui donne un combustible dense qui brûle longtemps et régulièrement ;
\item \emph{combustion contrôlée} : les briquettes sont brûlées dans les foyers de cuisson ; la combustion de la cellulose, $\ce{C6H10O5 + 6O2 -> 6CO2 + 5H2O}$, libère la chaleur utilisée pour la cuisson.
\end{itemize}
Le produit obtenu est donc de la \textbf{chaleur}, directement utilisée par les ménages.

\textbf{3.} Conditions d'une combustion propre (deux parmi) : utiliser un papier \textbf{sec} ; utiliser un papier \textbf{trié}, exempt de plastiques et de produits chimiques ; assurer un \textbf{apport d'air suffisant} (combustion complète, sans dégagement de monoxyde de carbone) ; brûler dans un \textbf{foyer adapté} et non à l'air libre.

Avantages pour l'environnement (deux parmi) : \textbf{réduction du volume des déchets} urbains ; \textbf{réduction de la déforestation} (les briquettes remplacent partiellement le bois de chauffe) ; limitation des \textbf{brûlis sauvages} et de la pollution de l'air dans les quartiers ; création de revenus favorisant l'adhésion de la communauté à la gestion des déchets.
\end{exoresolu}

\begin{savaistu}[Des déchets à l'électricité]
Dans plusieurs pays, de grandes \textbf{unités de valorisation énergétique} incinèrent les déchets urbains (dont le papier) pour produire de la vapeur qui entraîne des turbines : une partie de l'électricité des villes vient ainsi des poubelles ! Au Cameroun, les initiatives restent artisanales (briquettes de papier et de sciure, foyers améliorés), mais elles s'inscrivent dans la même logique : transformer un problème d'assainissement en ressource énergétique locale, tout en ménageant nos forêts, parmi les plus riches d'Afrique centrale.
\end{savaistu}

\begin{exercice}[À toi de jouer]
Un élève affirme : « Brûler les papiers dans la cour de l'école, c'est déjà de la valorisation énergétique. »

1. Expliquer pourquoi cette affirmation est incorrecte.

2. Décrire les transformations que devrait subir ce papier pour être réellement valorisé en combustible domestique.

3. Donner un argument à présenter à ta communauté pour l'inciter à adopter les briquettes de papier.
\end{exercice}

\begin{corrige}[Exercice : À toi de jouer]
\textbf{1.} Brûler les papiers à l'air libre est une combustion \textbf{non contrôlée} : la chaleur n'est pas récupérée ni utilisée, et la combustion incomplète dégage des fumées toxiques (\ce{CO}, particules) qui polluent l'air. Il n'y a ni récupération d'énergie ni protection de la santé : ce n'est donc pas une valorisation énergétique, mais une simple destruction polluante.

\textbf{2.} Le papier devrait être \textbf{trié} (élimination des plastiques, métaux et papiers souillés), \textbf{séché}, \textbf{broyé}, puis \textbf{compressé en briquettes} et séché à nouveau ; ces briquettes seraient ensuite brûlées dans un \textbf{foyer adapté} avec un apport d'air suffisant, la chaleur étant utilisée pour la cuisson ou le chauffage.

\textbf{3.} Argument (au choix) : les briquettes de papier remplacent partiellement le bois de chauffe, ce qui \textbf{protège nos forêts} et réduit les dépenses des ménages ; elles réduisent en outre le volume des déchets qui encombrent et polluent nos quartiers.
\end{corrige}

En résumé, le papier usagé, grâce à sa richesse en cellulose, est un combustible dont la combustion contrôlée — après tri, séchage, broyage et compression en briquettes — fournit de la chaleur utilisable pour le chauffage, la vapeur voire l'électricité. Mais brûler n'est pas la seule destinée du papier : lorsqu'il est propre, il peut aussi retrouver une seconde vie sous forme de \emph{nouveau papier}. C'est le recyclage du papier en serviettes jetables, papier hygiénique et papier-cadeau que nous étudierons dans la section suivante.

\coursec{Recyclage du papier : serviettes jetables, papier hygiénique, papier-cadeau}

Dans la section précédente, nous avons vu que le papier usagé peut être \emph{brûlé} pour produire de l'énergie : c'est la valorisation énergétique. Mais détruire le papier, c'est aussi détruire ses fibres de cellulose, qui ont demandé des années de croissance à un arbre. Ne pourrait-on pas plutôt \emph{réutiliser} ces fibres pour fabriquer un nouveau papier ? C'est exactement le principe du \cle{recyclage du papier}, que nous allons étudier maintenant. Tu découvriras comment de vieux cahiers, des journaux ou des cartons récupérés dans les rues de Yaoundé ou de Douala peuvent renaître sous la forme de serviettes jetables, de papier hygiénique ou de papier-cadeau.

\courssub{Principe du recyclage du papier : des fibres de cellulose réutilisables}

\textbf{Observation.} Déchire une feuille de papier et observe le bord de la déchirure à la loupe (ou simplement de très près) : tu vois de petits filaments qui dépassent. Ce sont les \cle{fibres de cellulose}, les « briques » du papier.

\begin{experience}[Observer la structure fibreuse du papier]
\textbf{Matériel :} une feuille de papier journal, une feuille de papier propre, une loupe, un verre d'eau.\par
\textbf{Protocole :} (1) Déchirer les deux feuilles et observer les bords à la loupe. (2) Laisser tremper un morceau de papier journal dans l'eau pendant une nuit, puis frotter entre les doigts.\par
\textbf{Observations :} les bords déchirés montrent des filaments enchevêtrés ; après trempage, le papier se désagrège en une bouillie grisâtre de fibres.\par
\textbf{Interprétation :} le papier n'est pas une matière continue comme une feuille de plastique : c'est un \emph{feutre} de fibres de cellulose simplement emmêlées et soudées entre elles lors du séchage. L'eau défait ces liaisons (liaisons hydrogène) et libère les fibres, qui peuvent alors être réassemblées en une nouvelle feuille.\par
\textbf{Conclusion :} le papier usagé est une \emph{matière première} : ses fibres de cellulose peuvent être réutilisées pour fabriquer un nouveau papier. C'est le principe du recyclage.
\end{experience}

\begin{definition}[Pâte à papier]
La \cle{pâte à papier} est une \textbf{suspension de fibres de cellulose dans l'eau}, obtenue à partir de papier usagé désagrégé (pâte recyclée) ou de bois broyé (pâte vierge). C'est le produit intermédiaire de toute fabrication de papier : étalée en couche mince, égouttée, pressée puis séchée, elle donne une feuille de papier.
\end{definition}

\begin{aretenir}[L'idée directrice]
Recycler le papier consiste à \textbf{défaire} la feuille usagée pour retrouver les fibres (mise en pâte), à \textbf{nettoyer} ces fibres (désencrage), puis à \textbf{refabriquer} une feuille (égouttage, pressage, séchage) que l'on façonne selon le produit désiré. Contrairement à l'incinération, le recyclage conserve la matière : la cellulose effectue un nouveau cycle de vie.
\end{aretenir}

\courssub{Les étapes du recyclage du papier}

La transformation du papier usagé en nouveau papier suit une chaîne d'opérations bien ordonnées. Retenons cinq grandes étapes.

\textbf{Étape 1 : collecte et tri des papiers usagés.} Les papiers sont rassemblés (poubelles de tri, ramassage chez les revendeurs, collecte dans les écoles et les administrations) puis triés à la main. On élimine tout ce qui n'est pas de la cellulose : \textbf{agrafes}, trombones, \textbf{plastiques} (fenêtres d'enveloppes, films d'emballage), rubans adhésifs, ainsi que les \textbf{papiers souillés} (taches de graisse, restes alimentaires) qui contamineraient la pâte. Le tri sépare aussi les catégories : journaux, cartons, papiers blancs de bureau, car ils ne donnent pas la même qualité de pâte.

\textbf{Étape 2 : mise en pâte.} Le papier trié est \textbf{trempé dans l'eau} puis \textbf{désagrégé par action mécanique} : brassage vigoureux dans une cuve (le \emph{pulpeur} dans l'industrie, un simple mixeur ou un pilage en atelier artisanal). Les fibres se séparent et l'on obtient une \textbf{pâte à papier}, bouillie grise contenant environ \num{95} à \SI{99}{\percent} d'eau.

\textbf{Étape 3 : nettoyage et désencrage.} La pâte contient encore des \textbf{encres} d'imprimerie et des impuretés (poussières, restes de colle). On la nettoie par lavage et tamisage ; le \cle{désencrage} consiste à décoller les particules d'encre des fibres (souvent à l'aide d'un produit détergent et de bulles d'air qui entraînent l'encre à la surface, où elle est écumée : c'est la \emph{flottation}). Un \textbf{blanchiment éventuel} est réalisé selon le produit visé : indispensable pour un papier bien blanc, inutile pour un papier-cadeau que l'on colorera de toute façon.

\textbf{Étape 4 : formation de la feuille.} La pâte propre est \textbf{étalée en couche mince et régulière sur un tamis} : l'eau s'écoule à travers les mailles (\textbf{égouttage}) tandis que les fibres restent à la surface et s'enchevêtrent. On \textbf{presse} ensuite la feuille humide (entre des chiffons, sous une planche ou entre des rouleaux) pour chasser l'eau restante et souder les fibres entre elles, puis on laisse \textbf{sécher} (à l'air libre, au soleil, ou sur des cylindres chauffants en usine). On obtient la \textbf{feuille de papier recyclé}.

\textbf{Étape 5 : façonnage selon le produit.} La feuille est enfin transformée selon sa destination :
\begin{itemize}
\item \textbf{serviettes jetables} : découpe en carrés ou rectangles, pliage, éventuel gaufrage (motifs en relief) ;
\item \textbf{papier hygiénique} : pâte très fine et bien désencrée, feuilles minces, douces et souples, enroulées en \textbf{rouleaux} et prédécoupées ;
\item \textbf{papier-cadeau} : feuilles plus épaisses, \textbf{décorées ou colorées} (teinture de la pâte, motifs imprimés ou peints à la main).
\end{itemize}

\begin{popfigure}
\begin{center}
\begin{tikzpicture}[font=\small,
etape/.style={rectangle, rounded corners=3pt, draw=popblue, fill=popblueL, text width=3.1cm, minimum height=1cm, align=center},
prod/.style={rectangle, rounded corners=3pt, draw=poporange, fill=poporangeL, text width=2.9cm, minimum height=0.9cm, align=center},
fleche/.style={-{Stealth[length=2.5mm]}, thick, popdark}]
\node[etape, align=center] (e1) at (0,0){\textbf{1. Collecte et tri}\\ (agrafes, plastiques, papiers souillés éliminés)};
\node[etape, align=center] (e2) at (4.4,0){\textbf{2. Mise en pâte}\\ (trempage + désagrégation mécanique dans l'eau)};
\node[etape, align=center] (e3) at (8.8,0){\textbf{3. Nettoyage et désencrage}\\ (encres, impuretés ; blanchiment éventuel)};
\node[etape, align=center] (e4) at (4.4,-2.6){\textbf{4. Égouttage, étalement, pressage, séchage}\\ $\rightarrow$ feuille de papier recyclé};
\node[prod, align=center] (p1) at (0.6,-5.2){\textbf{Serviettes jetables}\\ (découpe, pliage)};
\node[prod, align=center] (p2) at (4.4,-5.2){\textbf{Papier hygiénique}\\ (feuilles fines, douces, en rouleaux)};
\node[prod, align=center] (p3) at (8.2,-5.2){\textbf{Papier-cadeau}\\ (feuilles décorées ou colorées)};
\draw[fleche] (e1) -- (e2);
\draw[fleche] (e2) -- (e3);
\draw[fleche] (e3.south) |- ++(0,-0.6) -| (e4.east);
\draw[fleche] (e4.south) -- ++(0,-0.5) -| (p1.north);
\draw[fleche] (e4.south) -- ++(0,-1.0) |- (p2.north);
\draw[fleche] (e4.south) -- ++(0,-0.5) -| (p3.north);
\end{tikzpicture}
\end{center}
\legende{Organigramme des étapes du recyclage du papier. En bleu : les quatre étapes de fabrication de la feuille recyclée ; en orange : les trois produits obtenus au façonnage.}
\end{popfigure}

\begin{attention}[Le papier ne se recycle pas indéfiniment]
Erreur fréquente : croire que \emph{tout} papier est recyclable \emph{à l'infini}. En réalité, à chaque recyclage, les traitements mécaniques et chimiques \textbf{raccourcissent et fragilisent les fibres de cellulose}. Au bout d'environ \textbf{5 à 7 cycles}, les fibres deviennent trop courtes pour s'assembler en une feuille solide : elles finissent en boues, valorisées alors par incinération (production d'énergie) ou compostage. De plus, certains papiers ne sont \textbf{pas recyclables} : papiers souillés (gras, alimentaire), papiers plastifiés ou cirés, papiers photo, tickets thermiques. C'est pourquoi l'industrie ajoute toujours une part de fibres vierges, et pourquoi le tri (étape 1) est décisif.
\end{attention}

\courssub{Atelier scolaire : fabriquer une feuille de papier recyclé}

La fabrication artisanale du papier est un excellent travail pratique : elle reprend, à petite échelle, toutes les étapes industrielles.

\begin{methode}[Réaliser du papier recyclé en atelier scolaire]
\begin{enumerate}
\item \textbf{Tremper} : déchirer en petits morceaux du papier usagé propre (cahiers, journaux) et le faire tremper dans l'eau pendant plusieurs heures, idéalement une nuit.
\item \textbf{Mélanger au mixeur} : mixer les morceaux trempés avec de l'eau jusqu'à obtenir une pâte homogène (la pâte à papier). Ajouter éventuellement un colorant ou des pétales de fleurs pour le décor.
\item \textbf{Étaler sur le tamis} : plonger un cadre muni d'un tamis fin (moustiquaire tendue) dans la bassine de pâte, le relever à l'horizontale et laisser égoutter : une couche régulière de fibres se dépose sur le tamis.
\item \textbf{Presser} : retourner délicatement la couche de fibres sur un tissu absorbant, couvrir d'un second tissu et presser (éponge, rouleau ou planche lestée) pour extraire l'eau.
\item \textbf{Sécher à l'air} : laisser sécher la feuille à plat, à l'air libre ou au soleil, pendant 24 à \SI{48}{h}, puis la décoller doucement : la feuille de papier recyclé est prête.
\end{enumerate}
\end{methode}

\begin{popfigure}
\begin{center}
\begin{tikzpicture}[font=\small, scale=1]
% coupe du cadre avec tamis
\draw[very thick, popdark] (-3.2,0) -- (-3.2,1.6) -- (3.2,1.6) -- (3.2,0);
\draw[very thick, popdark] (-3.2,0) -- (3.2,0);
% tamis
\draw[popblue, thick] (-3.2,0.5) -- (3.2,0.5);
\foreach \x in {-3.0,-2.6,...,3.0} \draw[popblue!50] (\x,0.42) -- (\x,0.58);
\foreach \y in {0.44,0.5,0.56} \draw[popblue!50] (-3.2,\y) -- (3.2,\y);
% pâte étalée
\fill[popgoldL] (-3.0,0.5) rectangle (3.0,1.1);
\draw[popgold, thick] (-3.0,1.1) -- (3.0,1.1);
% gouttes d'eau qui s'égouttent
\foreach \x in {-2.2,-0.8,0.6,2.0} {\draw[popblue, -{Stealth[length=1.8mm]}] (\x,0.35) -- (\x,-0.55);}
\node[popblue] at (0,-0.95) {eau égouttée à travers le tamis};
% légendes fléchées
\draw[-{Stealth[length=2mm]}, popdark] (3.4,1.35) -- (3.2,1.35);
\node[right, popdark] at (3.4,1.35) {1. cadre en bois};
\draw[-{Stealth[length=2mm]}, popdark] (3.4,0.85) -- (3.0,0.85);
\node[right, popdark] at (3.4,0.85) {2. pâte étalée en couche mince};
\draw[-{Stealth[length=2mm]}, popdark] (-3.4,0.5) -- (-3.2,0.5);
\node[left, popdark] at (-3.4,0.5) {3. tamis (moustiquaire)};
% feuille récupérée
\draw[fill=popcream, draw=popgold, thick] (5.6,-0.6) rectangle (8.4,1.4);
\node[popdark] at (7.0,0.4) {4. feuille};
\node[popdark] at (7.0,-0.05) {récupérée};
\draw[-{Stealth[length=2.2mm]}, thick, poporange] (3.6,-1.5) .. controls (5,-2) .. (6.6,-0.8);
\node[poporange] at (5.1,-2.15) {après pressage et séchage};
\end{tikzpicture}
\end{center}
\legende{Coupe simplifiée d'une forme à papier artisanale. 1 : cadre rigide ; 2 : pâte à papier étalée en couche mince ; 3 : tamis qui retient les fibres et laisse passer l'eau ; 4 : feuille de papier recyclé obtenue après égouttage, pressage et séchage.}
\end{popfigure}

\courssub{Des produits différents pour des usages différents}

Les trois produits du programme — serviettes jetables, papier hygiénique, papier-cadeau — sortent de la même pâte de base, mais diffèrent par la \textbf{qualité de la pâte}, le \textbf{degré de blanchiment}, l'\textbf{épaisseur} et le \textbf{traitement de surface}.

\begin{center}
\begin{tabularx}{\linewidth}{|l|X|X|X|}
\hline
\rowcolor{popblueL} \textbf{Critère} & \textbf{Serviettes jetables} & \textbf{Papier hygiénique} & \textbf{Papier-cadeau} \\
\hline
Qualité de la pâte & moyenne, fibres courtes acceptées & pâte fine, très bien désencrée et lavée & pâte ordinaire, encres résiduelles tolérées \\
\hline
Blanchiment & moyen à fort (aspect propre) & fort (hygiène, blancheur) & faible ou nul (on colore ensuite) \\
\hline
Épaisseur & mince mais absorbante & très mince, souple et douce & plus épaisse et résistante \\
\hline
Traitement de surface & gaufrage, pliage & adoucissement, rouleaux prédécoupés & coloration, motifs décoratifs, parfois vernissage \\
\hline
\end{tabularx}
\end{center}

Ainsi, le papier hygiénique exige la pâte la plus propre (contact avec le corps : désencrage et blanchiment poussés), tandis que le papier-cadeau accepte une pâte plus grossière puisqu'il sera recouvert de couleurs et de motifs : c'est le produit le plus facile à fabriquer artisanalement, ce qui explique son succès dans les petites unités camerounaises.

\begin{exemplebox}[Papier-cadeau artisanal au Cameroun]
Au Cameroun, des petites unités artisanales et des clubs scientifiques de lycées fabriquent du \textbf{papier-cadeau} et des \textbf{emballages} (sachets, pochettes, cartons d'emballage) à partir de papiers récupérés : vieux cahiers, journaux, cartons collectés dans les quartiers. La pâte, teintée avec des colorants, est étalée sur des cadres à tamis, séchée au soleil, puis décorée de motifs. Ces produits sont vendus sur les marchés locaux, notamment en période de fêtes. Cette activité crée des revenus, occupe des jeunes et assainit les rues en réduisant la quantité de papiers abandonnés dans la nature.
\end{exemplebox}

\begin{savaistu}[Une tonne de papier recyclé, combien d'arbres sauvés ?]
Recycler \textbf{1 tonne} de papier permet de sauver environ \textbf{17 arbres}, d'économiser plusieurs \textbf{milliers de litres d'eau} (de l'ordre de \num{25000} à \SI{50000}{L} par rapport à la fabrication de papier neuf) et une quantité importante d'\textbf{énergie} (la fabrication à partir de pâte recyclée consomme nettement moins que la production de pâte vierge à partir du bois). À l'échelle d'une ville comme Douala, où des centaines de tonnes de papiers sont jetées chaque année, le recyclage représente donc un enjeu écologique considérable : moins d'arbres abattus dans nos forêts, moins de déchets dans les rues et les caniveaux, moins de dépenses d'enfouissement.
\end{savaistu}

\begin{exoresolu}[Organiser la chaîne de recyclage]
\textbf{Énoncé.} Un club environnemental d'un lycée de Bafoussam veut fabriquer du papier-cadeau à partir des vieux cahiers de fin d'année. (1) Citer, dans l'ordre, les cinq grandes étapes du procédé. (2) Expliquer pourquoi le club peut se passer de l'étape de blanchiment. (3) Expliquer pourquoi, après plusieurs recyclages successifs, la feuille obtenue devient fragile et se déchire facilement.
\tcblower
\textbf{Solution détaillée.}\par
(1) Les cinq étapes sont : \textbf{collecte et tri} des papiers usagés (élimination des agrafes, plastiques et papiers souillés) ; \textbf{mise en pâte} (trempage et désagrégation mécanique dans l'eau) ; \textbf{nettoyage et désencrage} de la pâte ; \textbf{égouttage, étalement en feuille mince, pressage et séchage} ; \textbf{façonnage} (ici, coloration et décoration des feuilles pour en faire du papier-cadeau).\par
(2) Le blanchiment est inutile car le papier-cadeau sera de toute façon \textbf{coloré et décoré} : la teinte grise de la pâte recyclée sera masquée par les colorants. Supprimer cette étape économise des produits chimiques et réduit la pollution de l'atelier.\par
(3) À chaque recyclage, les traitements mécaniques (mixage, pressage) \textbf{raccourcissent les fibres de cellulose}. Or la solidité d'une feuille de papier vient de l'enchevêtrement de fibres longues. Lorsque les fibres deviennent trop courtes (après environ 5 à 7 cycles), elles s'accrochent mal les unes aux autres : la feuille devient fragile et se déchire facilement. Il faut alors incorporer des fibres vierges ou orienter la pâte vers la valorisation énergétique.
\end{exoresolu}

\begin{exercice}[À toi de jouer]
(1) Définir la pâte à papier. (2) Pour chacun des trois produits (serviettes jetables, papier hygiénique, papier-cadeau), indiquer le degré de blanchiment requis et justifier. (3) Proposer deux arguments que tu utiliserais pour convaincre le proviseur de ton lycée d'installer des bacs de tri du papier.
\end{exercice}

\begin{corrige}[Exercice]
(1) La pâte à papier est une suspension de fibres de cellulose dans l'eau, obtenue à partir de papier usagé désagrégé (ou de bois broyé). (2) Serviettes jetables : blanchiment moyen à fort, car le produit doit paraître propre ; papier hygiénique : blanchiment fort, car il est en contact avec le corps (hygiène) ; papier-cadeau : blanchiment faible ou nul, car la feuille sera colorée et décorée. (3) Arguments possibles : le tri permet de revendre ou de transformer les papiers (revenus pour le lycée, fabrication de papier-cadeau par le club scientifique) ; il réduit la quantité de déchets brûlés ou abandonnés dans la cour, donc la pollution ; il éduque les élèves à la protection de l'environnement et à la préservation des forêts (chaque tonne recyclée sauve environ 17 arbres).
\end{corrige}

En résumé, le recyclage du papier repose sur une propriété remarquable de la cellulose : ses fibres peuvent être séparées puis réassemblées en une nouvelle feuille, plusieurs fois de suite. Maîtrisée industriellement ou artisanalement, cette chaîne de transformation donne des produits utiles au quotidien — serviettes jetables, papier hygiénique, papier-cadeau — tout en épargnant nos forêts. Mais le papier n'est pas le seul déchet valorisable : les \textbf{déchets plastiques}, bien plus persistants dans la nature, peuvent eux aussi être recyclés, comme nous le verrons dans la section suivante.

\coursec{Recyclage des déchets plastiques en bouteilles recyclées}

Après avoir étudié le recyclage du papier en serviettes jetables, en papier hygiénique et en papier-cadeau, nous abordons maintenant le second grand type de déchet valorisable : le \cle{plastique}. À Douala comme à Yaoundé, les bouteilles en plastique jonchent les rues, bouchent les caniveaux et provoquent des inondations en saison des pluies. Pourtant, ces bouteilles usagées ne sont pas des déchets sans valeur : ce sont des matières premières qui peuvent renaître sous forme de \cle{bouteilles recyclées}. Comprendre comment, c'est comprendre l'une des voies les plus prometteuses de la gestion des déchets au Cameroun.

\courssub{Nature des plastiques : des polymères non biodégradables}

\begin{definition}[Plastique]
Un \cle{plastique} est un matériau synthétique constitué de \cle{polymères}, c'est-à-dire de très longues molécules formées par la répétition de petites unités appelées monomères. Les plastiques sont \cle{non biodégradables} : les micro-organismes du sol ne possèdent pas les enzymes capables de couper leurs longues chaînes carbonées. Une bouteille plastique abandonnée dans la nature peut persister pendant plusieurs centaines d'années (on estime souvent \num{100} à \num{1000} ans selon le type).
\end{definition}

Cette non-biodégradabilité explique l'accumulation dramatique des plastiques dans l'environnement. Mais elle cache aussi une bonne nouvelle : puisque le plastique ne se dégrade pas, il conserve ses propriétés et peut être \textbf{retravaillé} presque indéfiniment, à condition de connaître son comportement à la chaleur.

\begin{experience}[Observation : le comportement du plastique à la chaleur]
\textbf{Matériel :} un fragment de bouteille en plastique (PET), un fragment de sachet plastique (PE), un couvercle de vieille marmite en bakélite (plastique dur), une source de chaleur (bec Bunsen ou plaque chauffante), une pince métallique, une hotte aspirante ou lieu aéré.

\textbf{Protocole :} on approche prudemment chaque fragment de la source de chaleur (sans le mettre dans la flamme directe pour éviter la combustion) et on observe les transformations physiques.

\textbf{Observations :}
\begin{itemize}
\item le fragment de bouteille (PET) et le fragment de sachet (PE) \textbf{se ramollissent}, se déforment, puis fondent en une pâte visqueuse ;
\item le couvercle en bakélite \textbf{noircit, dégage une odeur âcre et se carbonise} sans passer par un état ramolli.
\end{itemize}

\textbf{Interprétation :} il existe deux grandes familles de plastiques aux comportements thermiques opposés. Les uns se ramollissent à la chaleur et durcissent en refroidissant, de façon réversible : ce sont les \cle{thermoplastiques}. Les autres, une fois formés par une réaction chimique irréversible (réticulation), ne peuvent plus se refondre : ce sont les \cle{thermodurcissables} ; chauffés, ils se décomposent (pyrolyse) au lieu de fondre.

\textbf{Conclusion :} seuls les thermoplastiques peuvent être recyclés par fusion. C'est cette propriété physique réversible qui est à la base de tout le recyclage des bouteilles plastiques.
\end{experience}

\begin{definition}[Thermoplastique (savoir admis)]
Un \cle{thermoplastique} est un plastique dont les chaînes polymères sont linéaires ou faiblement ramifiées, sans liaisons covalentes interchaînes (ponts). Il \textbf{se ramollit à la chaleur et durcit en refroidissant}, de façon réversible. Cette propriété permet sa \textbf{refonte} : on peut le fondre et lui donner une nouvelle forme par moulage ou extrusion. Exemples : le \cle{PET} (polyéthylène téréphtalate) des bouteilles d'eau et de boissons gazeuses, le PE (polyéthylène) des sachets et films, le PP (polypropylène) des bouchons et cuvettes.
\end{definition}

\begin{attention}[Ne pas confondre thermoplastiques et thermodurcissables]
Les \textbf{thermoplastiques} (bouteilles PET, sachets PE, cuvettes PP, jouets) se ramollissent à la chaleur et sont \textbf{recyclables par fusion}. Les \textbf{thermodurcissables} (bakélite des poignées de marmites, résines époxy, mélamine des assiettes incassables, certaines prises électriques) ont subi une réticulation irréversible : ils \textbf{ne se refondent pas} et se décomposent (brûlent) si on les chauffe fortement. Au Baccalauréat, ne réponds jamais que « tous les plastiques sont recyclables par fusion » : \textbf{seuls les thermoplastiques le sont}.
\end{attention}

\courssub{Les étapes du recyclage des bouteilles plastiques}

Le recyclage des déchets plastiques en bouteilles recyclées suit une chaîne d'opérations rigoureusement ordonnée. Chaque étape prépare la suivante ; en sauter une ou les inverser compromet la qualité du produit final.

\textbf{Étape 1 : la collecte et le tri.} Les bouteilles usagées sont ramassées (par les ménages, les écoles, les collecteurs informels) puis \textbf{triées}. Le tri se fait selon deux critères indispensables :
\begin{itemize}
\item \textbf{par type (nature chimique)} : on sépare les bouteilles (PET), les sachets (PE) et les objets rigides (seaux, chaises en PP), car chaque polymère a sa propre température de fusion et sa rhologie propre ;
\item \textbf{par couleur} : on sépare les plastiques transparents (incolores), bleus, verts et colorés, afin d'obtenir un produit final d'aspect homogène sans avoir à ajouter de pigments coûteux.
\end{itemize}

\textbf{Étape 2 : le nettoyage et le lavage.} Les bouteilles triées sont lavées à l'eau chaude (souvent additionnée de détergent ou de soude diluée) dans des laveurs industriels. Cette étape élimine les \textbf{saletés} (poussière, sable), les \textbf{étiquettes} en papier et les \textbf{résidus} de boissons (sucres, acides) ou de colles. Un plastique mal lavé donnerait une bouteille recyclée tachée, odorante, fragile et impropre au contact alimentaire.

\textbf{Étape 3 : le broyage.} Les bouteilles propres et égouttées passent dans un \textbf{broyeur} muni de lames rotatives qui les découpent en petits fragments réguliers appelés \cle{paillettes} (ou \cle{flocons}), de taille généralement comprise entre \SI{5}{\milli\metre} et \SI{10}{\milli\metre}. Le broyage augmente considérablement la surface de contact avec la chaleur et permet une fusion rapide, homogène et économe en énergie à l'étape suivante. \emph{Note : dans les filières industrielles modernes, les paillettes sont souvent ensuite extrudées et granulées pour former des \cle{granulés} (pellets) cylindriques, plus faciles à doser et à transporter vers les unités de transformation.}

\textbf{Étape 4 : la fusion contrôlée et le moulage par soufflage.} Les paillettes (ou granulés) sont introduites dans une \textbf{extrudeuse} (vis sans fin chauffée) où elles sont \textbf{chauffées de façon contrôlée} jusqu'à leur ramollissement complet. Pour le PET, la température de fusion est maîtrisée autour de \SI{260}{\celsius} (entre \SI{250}{\celsius} et \SI{280}{\celsius}) pour éviter la dégradation thermique (jaunissement, perte de viscosité, formation d'acétaldéhyde). La pâte plastique homogène ainsi obtenue est ensuite façonnée par \textbf{moulage par soufflage} (technique standard pour les corps creux) :
\begin{enumerate}
\item \textbf{Injection de la paraison :} la matière fondue est injectée dans un moule de préformage pour créer une ébauche tubulaire épaisse, fermée à un bout, appelée \cle{paraison} (ou préforme).
\item \textbf{Soufflage :} la paraison, encore chaude et molle, est transférée dans un moule en forme de bouteille (moule de soufflage). On y insuffle de l'\textbf{air comprimé} (pression de \SI{10}{\bar} à \SI{25}{\bar}) qui plaque le plastique ramolli contre les parois froides du moule.
\end{enumerate}

\textbf{Étape 5 : le refroidissement, le contrôle de qualité et le conditionnement.} La bouteille formée dans le moule est \textbf{refroidie} (par circulation d'eau dans les parois du moule) : le thermoplastique durcit en gardant exactement la forme du moule. Chaque bouteille est ensuite soumise à un \textbf{contrôle de qualité} automatisé (vérification de l'épaisseur des parois par ultrasons, absence de fissures ou de bulles, test d'étanchéité à la pression, contrôle du poids et des dimensions). Les bouteilles conformes sont \textbf{conditionnées} (regroupées en fagots, filmées, palettisées) pour être livrées aux embouteilleurs.

\begin{methode}[Expliquer le recyclage du plastique en bouteilles recyclées]
Pour expliquer correctement cette technique au Baccalauréat ou à ta communauté, retiens la chaîne d'opérations dans l'ordre chronologique et justifie chaque étape :
\begin{enumerate}
\item \textbf{Collecte et tri} (par type de polymère et par couleur) : garantit l'homogénéité du lot.
\item \textbf{Lavage} (eau + détergent) : élimine impuretés, étiquettes, colles et résidus organiques.
\item \textbf{Broyage} en paillettes (flocons) : augmente la surface spécifique pour une fusion homogène.
\item \textbf{Fusion contrôlée} (extrusion vers \SI{260}{\celsius} pour le PET) : ramollissement réversible du thermoplastique.
\item \textbf{Moulage par soufflage} : formation de la paraison puis gonflage à l'air comprimé dans le moule final.
\item \textbf{Refroidissement}, contrôle qualité (étanchéité, dimensions, aspect) et conditionnement.
\end{enumerate}
Justifie toujours la possibilité de cette technique par la propriété fondamentale des thermoplastiques : \textbf{ramollissement à la chaleur et durcissement au refroidissement, de façon réversible}.
\end{methode}

\begin{popfigure}
\begin{center}
\begin{tikzpicture}[
  font=\small,
  etape/.style={draw=popblue, fill=popblueL, rounded corners=3pt, text width=2.8cm, align=center, minimum height=1.2cm, text=popdark, inner sep=2pt},
  fleche/.style={-{Stealth[length=3mm]}, thick, poporange}
]
\node[etape, align=center] (a) at (0,0){\textbf{1. Collecte \& Tri}\\(par type \& couleur)};
\node[etape, align=center] (b) at (3.8,0){\textbf{2. Lavage}\\(eau, détergent, séchage)};
\node[etape, align=center] (c) at (7.6,0){\textbf{3. Broyage}\\(paillettes / flocons)};
\node[etape, align=center] (d) at (11.4,0){\textbf{4. Fusion / Extrusion}\\($\sim$260°C pour PET)};
\node[etape, align=center] (e) at (11.4,-2.8){\textbf{5. Moulage par soufflage}\\(paraison + air comprimé)};
\node[etape, align=center] (f) at (7.6,-2.8){\textbf{6. Refroidissement}\\(durcissement dans moule)};
\node[etape, align=center] (g) at (3.8,-2.8){\textbf{7. Contrôle qualité}\\(étanchéité, dimensions)};
\node[etape, fill=popgreenL, draw=popgreen, align=center] (h) at (0,-2.8){\textbf{8. Bouteilles recyclées}\\(conditionnées)};
\draw[fleche] (a) -- (b);
\draw[fleche] (b) -- (c);
\draw[fleche] (c) -- (d);
\draw[fleche] (d) -- (e);
\draw[fleche] (e) -- (f);
\draw[fleche] (f) -- (g);
\draw[fleche] (g) -- (h);
\end{tikzpicture}
\end{center}
\legende{Organigramme du recyclage des déchets plastiques en bouteilles recyclées (filière PET). La chaîne repose sur la réversibilité thermique des thermoplastiques : fusion $\rightarrow$ mise en forme $\rightarrow$ solidification.}
\end{popfigure}

\begin{popfigure}
\begin{center}
\begin{tikzpicture}[font=\small, scale=0.9]
% 1. Paraison
\begin{scope}[xshift=0cm]
\draw[thick, popblue, fill=popblueL] (-0.45,0.3) -- (-0.45,-1.8) arc[start angle=180, end angle=360, radius=0.45] -- (0.45,0.3);
\draw[thick, popblue, fill=popblueL] (-0.6,0.3) rectangle (0.6,0.8); % goulot
\node[popdark, align=center] at (0,-2.5) {1. Paraison\\chauffée (molle)};
\draw[-{Stealth[length=3mm]}, poporange, thick] (0,1.7) -- (0,0.95);
\node[poporange, right] at (0.7,1.4) {Air comprimé};
\end{scope}
% 2. Moule fermé + soufflage
\begin{scope}[xshift=5.0cm]
% Moule (deux demi-coques)
\draw[thick, poppurple, fill=poppurpleL, opacity=0.3] (-1.2,0.8) -- (-1.2,-2.3) -- (-0.6,-2.3) -- (-0.6,0.8) -- cycle;
\draw[thick, poppurple, fill=poppurpleL, opacity=0.3] (1.2,0.8) -- (1.2,-2.3) -- (0.6,-2.3) -- (0.6,0.8) -- cycle;
% Plastique soufflé
\draw[thick, popblue, fill=popblueL] (-0.55,0.3) -- (-0.55,-1.6) arc[start angle=180, end angle=360, radius=0.55] -- (0.55,0.3);
\draw[thick, popblue, fill=popblueL] (-0.65,0.3) rectangle (0.65,0.8);
\draw[-{Stealth[length=3mm]}, poporange, thick] (0,1.8) -- (0,0.95);
\node[poporange, right] at (0.7,1.5) {Air comprimé};
\node[poppurple, left, align=center] at (-1.3,-0.7){Moule\\froid};
\node[popdark, align=center] at (0,-2.8) {2. Soufflage\\dans le moule};
\end{scope}
% 3. Bouteille formée
\begin{scope}[xshift=10.0cm]
\draw[thick, popgreen, fill=popgreenL] (-0.3,1.0) -- (-0.3,0.5) -- (-0.7,-0.2) -- (-0.7,-2.0) arc[start angle=180, end angle=360, radius=0.7] -- (0.7,-0.2) -- (0.3,0.5) -- (0.3,1.0) -- cycle;
\draw[thick, popgreen, fill=popgreenL] (-0.35,1.0) rectangle (0.35,1.3); % bague
\node[popdark, align=center] at (0,-2.8) {3. Bouteille formée\\après refroidissement};
\end{scope}
\draw[-{Stealth[length=3mm]}, very thick, poporange] (2.0,-0.7) -- (2.9,-0.7);
\draw[-{Stealth[length=3mm]}, very thick, poporange] (7.0,-0.7) -- (7.9,-0.7);
\end{tikzpicture}
\end{center}
\legende{Principes du moulage par soufflage (extrusion-soufflage ou injection-soufflage). 1 : La paraison (tube de PET fondu) est maintenue chaude. 2 : Enfermée dans le moule froid, elle est gonflée par air comprimé qui la plaque contre les parois. 3 : Après refroidissement rapide, la bouteille durcit et est éjectée.}
\end{popfigure}

\courssub{L'importance capitale du tri}

\begin{propriete}[Conséquence du mélange de plastiques différents]
Chaque type de plastique possède sa propre \textbf{température de fusion} ($T_f$), sa propre \textbf{viscosité à l'état fondu} et sa propre structure chimique. \textbf{Mélanger des plastiques de natures différentes} (par exemple du PET $T_f \approx \SI{260}{\celsius}$ avec du PE $T_f \approx \SI{110}{\celsius}$ ou du PP $T_f \approx \SI{160}{\celsius}$) \textbf{dégrade irréversiblement la qualité du produit final} : le polymère à $T_f$ basse brûle/dégrade avant que l'autre ne fonde, ou bien le mélange fondu est hétérogène (phases non miscibles). La bouteille recyclée obtenue est alors hétérogène, mécaniquement fragile (rupture aux interfaces), mal colorée et impropre à l'usage alimentaire. Le tri rigoureux par type et par couleur est donc l'étape qui conditionne la réussite de tout le recyclage.
\end{propriete}

C'est pourquoi, dans les filières bien organisées, le tri commence dès la source : dans les ménages (bacs jaunes), les écoles et les marchés. Sensibiliser sa communauté à trier les bouteilles (PET) séparément des sachets (PE) et des bouchons (PP), c'est déjà œuvrer pour la qualité des bouteilles recyclées et la rentabilité de la filière.

\courssub{Autres produits obtenus à partir des granulés plastiques}

Les paillettes et granulés issus du broyage (et de l'extrusion) ne servent pas qu'à fabriquer de nouvelles bouteilles (recyclage \emph{boucle fermée} ou \emph{bottle-to-bottle}). Par moulage par injection, extrusion ou soufflage, on obtient de nombreux \textbf{objets ménagers et industriels} utiles à la vie quotidienne : \textbf{arrosoirs}, \textbf{cuvettes}, \textbf{seaux}, \textbf{tuyaux} d'arrosage, cintres, bols, bassines, palettes, bancs publics, poteaux de clôture. Cette diversité de débouchés (recyclage \emph{boucle ouverte}) renforce l'intérêt économique de la collecte des plastiques et assure un marché même si la demande en bouteilles recyclées fluctue.

\begin{savaistu}[Le PET, champion mondial du recyclage]
Le \cle{PET} (polyéthylène téréphtalate) des bouteilles est l'un des plastiques \textbf{les plus recyclés au monde} (taux de collecte > 50\% en Europe, en progression au Cameroun). Au-delà des nouvelles bouteilles (qualité alimentaire), le PET recyclé (r-PET) est massivement transformé en \textbf{fibres textiles} : les paillettes fondues sont extrudées en filaments fins, étirés et texturés pour produire du \textbf{polyester recyclé}. Ce tissu sert à fabriquer des polaires, des tapis, de la ouate pour couettes et vestes, des sacs, et même des chaussures de sport. \textbf{Environ 10 à 15 bouteilles de 1,5 L recyclées suffisent à produire un tee-shirt en polyester} ou une polaire légère. C'est un exemple parfait d'économie circulaire.
\end{savaistu}

\begin{exemplebox}[L'exemple camerounais : de la rue à l'usine]
Au Cameroun, la filière plastique structure une économie circulaire naissante. Des \textbf{entreprises} (ex. \emph{Redplast}, \emph{Cameroon Plastic Recycling}, \emph{EcoCollect} à Douala/Yaoundé) et de nombreux \textbf{jeunes entrepreneurs} (collecteurs, broyeurs, artisans) ont fait de la bouteille plastique une ressource. Des collecteurs parcourent les quartiers pour ramasser les bouteilles usagées, qu'ils \textbf{revendent} au kg (environ 50--100 FCFA/kg selon les cours) aux unités de transformation ou \textbf{transforment} eux-mêmes en paillettes/granulés, en objets ménagers (cuvettes, seaux) ou en matériaux de construction (pavés, tuiles -- voir section suivante). Cette activité crée des emplois directs et indirects pour des milliers de jeunes, génère des revenus pour les ménages collecteurs et assainit visiblement les quartiers : la preuve concrète que la protection de l'environnement et le développement économique peuvent avancer ensemble.
\end{exemplebox}

\begin{exoresolu}[Ordonner et justifier les étapes du recyclage]
\textbf{Énoncé.} Un élève présente les opérations du recyclage des bouteilles plastiques dans l'ordre suivant : broyage, fusion, collecte et tri, moulage, lavage, refroidissement. (a) Remets ces opérations dans le bon ordre. (b) Explique pourquoi le lavage doit obligatoirement précéder le broyage et la fusion. (c) Indique la propriété physico-chimique des plastiques qui rend ce recyclage possible.
\tcblower
\textbf{Solution détaillée.}

(a) Le bon ordre chronologique est : 
\textbf{1. Collecte et tri} $\rightarrow$ \textbf{2. Lavage} $\rightarrow$ \textbf{3. Broyage} $\rightarrow$ \textbf{4. Fusion} $\rightarrow$ \textbf{5. Moulage (soufflage)} $\rightarrow$ \textbf{6. Refroidissement} (+ contrôle qualité et conditionnement).

(b) Le lavage élimine les saletés minérales (sable), les étiquettes (papier, colle) et les résidus organiques (sucres, acides, microbes). S'il était effectué \emph{après} le broyage, les impuretés seraient intimement mélangées aux paillettes et impossibles à séparer efficacement ; elles seraient incorporées dans la pâte plastique lors de la fusion, créant des défauts (bulles, taches, points de rupture, odeurs) et dégradant les propriétés mécaniques de la bouteille. Le lavage \emph{avant} broyage garantit une matière première propre et saine.

(c) Ce recyclage par refonte est possible parce que le plastique des bouteilles (le PET) est un \textbf{thermoplastique} : ses chaînes polymères ne sont pas réticulées. Il \textbf{se ramollit à la chaleur (fusion réversible) et durcit en refroidissant (solidification)}, sans modification chimique de sa structure. On peut donc le fondre et lui donner une nouvelle forme indéfiniment (dans la limite de la dégradation thermique cumulative).
\end{exoresolu}

\begin{attention}[Pièges fréquents à l'examen (Bac D)]
\begin{itemize}
\item \textbf{Vocabulaire :} Ne dis pas que le plastique est « brûlé » ou « incinéré » pour être recyclé : il est \textbf{fondu} (ramolli par chauffage contrôlé sans combustion), ce qui est radicalement différent de la combustion (oxydation) qui détruit la matière et pollue l'air (fumées toxiques, dioxines).
\item \textbf{Ordre des étapes :} N'oublie jamais le \textbf{tri} comme toute première étape (après collecte) : c'est lui qui garantit l'homogénéité du lot et la qualité du produit final. Un tri absent ou mal fait = recyclage raté.
\item \textbf{Terminologie technique :} Le terme exact de l'ébauche tubulaire gonflée à l'air comprimé est la \textbf{paraison} (ou préforme) ; le procédé industriel est le \textbf{moulage par soufflage} (et non « moulage » tout court qui suggère l'injection dans un moule plein).
\item \textbf{Température :} Pour le PET, cite une température de fusion réaliste (\textbf{environ \SI{260}{\celsius}}, entre \SI{250}{\celsius} et \SI{280}{\celsius}), pas \SI{100}{\celsius} (ébullition de l'eau) ni \SI{1000}{\celsius}.
\end{itemize}
\end{attention}

\begin{aretenir}[L'essentiel de la section : Recyclage du plastique en bouteilles]
\begin{itemize}
\item Les plastiques sont des \textbf{polymères non biodégradables} ; les \textbf{thermoplastiques} (PET, PE, PP) se ramollissent à la chaleur et durcissent en refroidissant (transition de phase réversible), ce qui permet leur refonte.
\item Le recyclage en bouteilles recyclées suit la chaîne industrielle : \textbf{Tri} (type + couleur) $\rightarrow$ \textbf{Lavage} $\rightarrow$ \textbf{Broyage} (paillettes) $\rightarrow$ \textbf{Fusion contrôlée} (extrusion $\sim$\SI{260}{\celsius} pour PET) $\rightarrow$ \textbf{Moulage par soufflage} (paraison + air comprimé) $\rightarrow$ \textbf{Refroidissement}, contrôle qualité, conditionnement.
\item \textbf{Mélanger des plastiques de natures différentes dégrade la qualité} du produit final (incompatibilité polymères) : le tri est l'étape clé, à faire dès la source.
\item Les granulés plastiques servent aussi à fabriquer arrosoirs, cuvettes, tuyaux, textiles (fibres r-PET), matériaux de construction (pavés), créant une économie circulaire locale.
\end{itemize}
\end{aretenir}

\coursec{Transformation des déchets plastiques en pavés}

Dans la section précédente, nous avons vu comment les bouteilles plastiques usagées peuvent être lavées, broyées, fondues puis soufflées pour donner naissance à de nouvelles bouteilles : c'est le recyclage « en boucle fermée ». Mais tous les déchets plastiques ne peuvent pas suivre cette voie : les sachets souillés, les emballages fins, les mélanges de plastiques de natures différentes sont difficiles à retransformer en emballages. Faut-il pour autant les abandonner dans la nature, où ils mettront des siècles à se dégrader ? Non. Une solution ingénieuse, développée notamment par de jeunes entrepreneurs africains, consiste à les transformer en \cle{pavés} de construction. C'est l'objet de cette section : comprendre le principe, décrire les étapes de la fabrication et évaluer les qualités du matériau obtenu.

\courssub{Principe de la transformation : le plastique fondu comme liant}

\textbf{Observation de départ.} Dans une cour d'école de Yaoundé, on compare deux types de pavés : des pavés classiques gris, fabriqués avec du ciment et du sable, et des pavés plus sombres, légèrement luisants, fabriqués à partir de déchets plastiques. Après plusieurs années d'usage, les deux résistent encore au passage des élèves et des motos. Pourtant, les seconds ne contiennent pas un gramme de ciment. Qu'est-ce qui tient les grains de sable ensemble ?

\textbf{Hypothèse.} Le plastique, une fois fondu puis refroidi, joue le rôle de \cle{liant}, exactement comme le ciment dans le béton.

\textbf{Vérification.} On réalise deux éprouvettes de sable : l'une mélangée à de l'eau seule, l'autre mélangée à du plastique fondu, puis laissée refroidir. La première s'effrite dès qu'elle sèche ; la seconde forme un bloc dur, cohérent, qui résiste à la compression et ne laisse pas passer l'eau. L'interprétation est claire : à chaud, le plastique fondu enrobe chaque grain de sable ; au refroidissement, il se solidifie et soude les grains les uns aux autres, formant un matériau composite.

\begin{definition}[Pavé plastique]
Un \cle{pavé plastique} (ou pavé écologique) est un bloc de construction obtenu par mélange de \textbf{plastique fondu}, qui joue le rôle de \cle{liant} (à la place du ciment), et de \textbf{sable}, qui joue le rôle de \cle{granulat} (squelette solide du matériau). Le mélange chaud est coulé dans un moule, compacté, puis refroidi et démoulé. Le plastique solidifié enrobe les grains de sable et assure la cohésion, l'imperméabilité et la résistance du bloc.
\end{definition}

\begin{aretenir}[L'idée directrice]
Dans un pavé classique : \textbf{ciment (liant) + sable (granulat) + eau}. Dans un pavé plastique : \textbf{plastique fondu (liant) + sable (granulat)}, sans eau ni ciment. La chaleur remplace la prise hydraulique du ciment : c'est le refroidissement du plastique qui fait « prendre » le matériau.
\end{aretenir}

\courssub{Les étapes de la fabrication}

\begin{methode}[Décrire la transformation des déchets plastiques en pavés]
Pour décrire correctement cette technique (au Baccalauréat comme dans un exposé de sensibilisation), il faut suivre l'ordre chronologique des opérations :
\begin{enumerate}
\item \textbf{Matière première} : collecte, tri et nettoyage des déchets plastiques ;
\item \textbf{Préparation} : découpe ou broyage en morceaux ;
\item \textbf{Fusion} : chauffage avec brassage jusqu'à obtention d'une pâte homogène ;
\item \textbf{Mélange} : incorporation du sable en proportions contrôlées et malaxage ;
\item \textbf{Moulage} : coulée dans des moules, compactage ;
\item \textbf{Refroidissement} puis \textbf{démoulage} : obtention du pavé fini.
\end{enumerate}
Retiens la logique : \emph{nettoyer → réduire → fondre → mélanger → mouler → refroidir}.
\end{methode}

\textbf{Étape 1 : collecte, tri et nettoyage.} Les déchets plastiques (sachets, bouteilles, emballages divers) sont ramassés dans les quartiers, les marchés ou les décharges. On les trie pour éliminer les corps étrangers (métaux, pierres, matières organiques) et on les lave à l'eau, souvent additionnée de savon, afin d'enlever les souillures. Un plastique sale donne un pavé hétérogène et fragile : les impuretés créent des zones de faiblesse dans le matériau. On laisse ensuite sécher les plastiques, car l'eau résiduelle provoquerait des projections et des bulles lors de la fusion.

\textbf{Étape 2 : découpe ou broyage.} Les plastiques propres et secs sont découpés au couteau ou broyés en petits morceaux de quelques centimètres. Cette étape augmente la surface de contact avec la source de chaleur : plus les morceaux sont petits, plus la fusion est rapide et homogène, et moins on consomme de combustible.

\textbf{Étape 3 : fusion.} Les morceaux sont placés dans un récipient métallique chauffé (sur un foyer au bois, au gaz ou au charbon). Sous l'effet de la chaleur, le plastique ramollit puis fond : on parle de \cle{fusion} (passage de l'état solide à l'état pâteux ou liquide). Un brassage régulier, à l'aide d'une tige métallique, assure une température uniforme et aboutit à une \textbf{pâte homogène}, sans grumeaux. Le chauffage doit rester contrôlé : une température excessive brûle le plastique, dégrade ses propriétés et dégage davantage de fumées.

\textbf{Étape 4 : incorporation du sable et malaxage.} Le sable, propre et sec, est ajouté progressivement à la pâte de plastique fondu, en proportions contrôlées. Un dosage courant est de l'ordre de \textbf{2 à 3 volumes de sable pour 1 volume de plastique fondu}, ajusté selon la qualité recherchée : plus de plastique donne un pavé plus imperméable et plus souple, plus de sable donne un pavé plus rigide mais plus cassant si le liant devient insuffisant. Le tout est malaxé énergiquement jusqu'à ce que chaque grain de sable soit enrobé de plastique : le mélange prend l'aspect d'une pâte sombre et granuleuse.

\textbf{Étape 5 : moulage, compactage, refroidissement et démoulage.} Le mélange chaud est rapidement versé dans des \cle{moules} métalliques (souvent des caissons carrés ou rectangulaires, par exemple de \SI{20}{cm} $\times$ \SI{10}{cm} $\times$ \SI{6}{cm}). On le compacte à l'aide d'un presseur ou d'une presse afin de chasser les bulles d'air et d'augmenter la densité du matériau. Le moule est ensuite laissé au repos : en refroidissant, le plastique se solidifie et le bloc durcit. Après refroidissement complet, on procède au \cle{démoulage} : le pavé est extrait du moule, prêt à l'emploi après contrôle de qualité (absence de fissures, régularité des faces).

\begin{popfigure}
\begin{center}
\begin{tikzpicture}[
font=\small,
etape/.style={rectangle, rounded corners=3pt, draw=popblue, fill=popblueL, text width=2.6cm, minimum height=1.1cm, align=center, inner sep=3pt},
fleche/.style={-{Stealth[length=2.5mm]}, thick, popdark}
]
\node[etape, align=center] (a) at (0,0){\textbf{Plastiques usagés}\\ sachets, bouteilles, emballages};
\node[etape, align=center] (b) at (3.6,0){\textbf{Tri et lavage}\\ élimination des impuretés};
\node[etape, align=center] (c) at (7.2,0){\textbf{Découpe / broyage}\\ petits morceaux};
\node[etape, align=center] (d) at (10.8,0){\textbf{Fusion}\\ chauffage + brassage};
\node[etape, align=center] (e) at (10.8,-2.2){\textbf{Mélange avec le sable}\\ malaxage (pâte homogène)};
\node[etape, align=center] (f) at (7.2,-2.2){\textbf{Moulage}\\ coulée + compactage};
\node[etape, align=center] (g) at (3.6,-2.2){\textbf{Refroidissement}\\ solidification du liant};
\node[etape, draw=popgreen, fill=popgreenL, align=center] (h) at (0,-2.2){\textbf{Pavé}\\ démoulage et contrôle};
\draw[fleche] (a) -- (b);
\draw[fleche] (b) -- (c);
\draw[fleche] (c) -- (d);
\draw[fleche] (d) -- (e);
\draw[fleche] (e) -- (f);
\draw[fleche] (f) -- (g);
\draw[fleche] (g) -- (h);
\node[draw=poporange, fill=poporangeL, rounded corners=2pt, inner sep=2pt, font=\footnotesize] (sable) at (12.9,-2.2) {Sable propre et sec};
\draw[fleche, poporange] (sable) -- (e);
\end{tikzpicture}
\end{center}
\legende{Organigramme de la transformation des déchets plastiques en pavés. Les huit étapes se succèdent dans l'ordre indiqué par les flèches ; le sable (encadré orange) est incorporé au moment du malaxage, après la fusion du plastique.}
\end{popfigure}

\begin{popfigure}
\begin{center}
\begin{tikzpicture}[font=\small, scale=1.05]
% Moule (coupe)
\draw[very thick, popdark, fill=popline!40] (0,0) -- (0,2.2) -- (0.35,2.2) -- (0.35,0.35) -- (3.65,0.35) -- (3.65,2.2) -- (4,2.2) -- (4,0) -- cycle;
% Mélange chaud dans le moule
\fill[poporange!55] (0.35,0.35) rectangle (3.65,1.5);
\foreach \x in {0.6,1.1,1.6,2.1,2.6,3.1,3.5} {\fill[popdark!60] (\x,0.7) circle (0.05); \fill[popdark!60] (\x,1.15) circle (0.05);}
% Presseur
\draw[thick, popblue, fill=popblueL] (1.2,2.6) rectangle (2.8,3.0);
\draw[thick, popblue] (2,3.0) -- (2,3.6);
\draw[-{Stealth[length=2.5mm]}, very thick, popblue] (2,3.55) -- (2,3.1);
% Pavé démoulé
\draw[very thick, popdark, fill=poporange!35] (5.6,0) rectangle (8.6,1.3);
\foreach \x in {5.9,6.5,7.1,7.7,8.3} {\fill[popdark!50] (\x,0.45) circle (0.05); \fill[popdark!50] (\x,0.9) circle (0.05);}
\draw[-{Stealth[length=2.5mm]}, thick, popgreen] (4.3,0.7) -- (5.4,0.7) node[midway, above, font=\footnotesize]{démoulage};
% Légendes fléchées
\draw[-{Stealth[length=2mm]}, popmuted] (2.0,1.0) -- (2.0,-0.5) node[below, align=center, font=\footnotesize]{1. Mélange chaud\\ (plastique fondu + sable)};
\draw[-{Stealth[length=2mm]}, popmuted] (0.18,1.3) -- (-0.9,1.3) node[left, align=right, font=\footnotesize]{2. Moule\\ métallique};
\draw[-{Stealth[length=2mm]}, popmuted] (2.8,2.8) -- (4.3,2.8) node[right, align=left, font=\footnotesize]{3. Presseur\\ (compactage)};
\draw[-{Stealth[length=2mm]}, popmuted] (7.1,0.65) -- (7.1,-0.5) node[below, align=center, font=\footnotesize]{4. Pavé démoulé\\ après refroidissement};
\end{tikzpicture}
\end{center}
\legende{Schéma du moulage d'un pavé plastique. Le mélange chaud (1) est coulé dans le moule métallique (2), compacté à l'aide du presseur (3) pour chasser l'air, puis, après refroidissement et solidification du liant plastique, démoulé pour donner le pavé fini (4).}
\end{popfigure}

\courssub{Propriétés des pavés obtenus et usages}

Les pavés plastique-sable présentent des propriétés remarquables, qui expliquent leur succès croissant :

\begin{itemize}
\item \textbf{Résistance mécanique} : le compactage et la cohésion assurée par le plastique donnent des blocs qui supportent bien la compression ; ils résistent aussi mieux aux chocs que les pavés en ciment, qui sont plus cassants.
\item \textbf{Imperméabilité} : le plastique est insoluble dans l'eau et ne se laisse pas traverser par elle ; le pavé ne se dégrade donc pas sous l'effet des pluies, contrairement aux pavés de terre stabilisée.
\item \textbf{Légèreté} : à dimensions égales, un pavé plastique est souvent plus léger qu'un pavé en béton, ce qui facilite le transport et la pose.
\item \textbf{Durabilité} : le plastique ne pourrit pas, n'est pas attaqué par les champignons ni par les insectes xylophages ; le pavé a une longue durée de vie.
\end{itemize}

Ces qualités en font un matériau adapté au revêtement des \textbf{cours d'écoles et d'habitations}, des \textbf{trottoirs}, des \textbf{parkings}, des allées piétonnes et des places de marché.

\begin{exemplebox}[Un matériau performant et économique]
Des essais de compression menés sur des pavés plastique-sable bien dosés montrent des résistances \textbf{comparables, voire supérieures} à celles des pavés classiques en ciment, alors que leur coût de fabrication est \textbf{inférieur} : la matière première (déchets plastiques) est quasiment gratuite, et le sable est abondant et bon marché au Cameroun. À qualité égale, le pavé écologique coûte donc moins cher tout en assainissant l'environnement : c'est un double gain, économique et écologique.
\end{exemplebox}

\begin{experience}[Comparer la cohésion de deux mélanges sableux (démonstration encadrée)]
\textbf{Matériel} : sable sec, eau, morceaux de sachets plastiques propres, récipient métallique, source de chaleur, tige de brassage, gants et masque, deux petits moules identiques.\par
\textbf{Protocole} : (1) dans le moule A, mélanger le sable avec un peu d'eau et tasser ; (2) dans le moule B, verser du sable mélangé à du plastique fondu (opération réalisée \emph{par l'enseignant}, en espace ventilé) et compacter ; (3) laisser sécher le moule A et refroidir le moule B, puis démouler et comparer la solidité des deux blocs (résistance à la pression des doigts, à la chute de faible hauteur, au ruissellement d'eau).\par
\textbf{Observations} : le bloc A s'effrite facilement et se désagrège sous l'eau ; le bloc B est dur, cohérent et imperméable.\par
\textbf{Interprétation} : l'eau n'est pas un liant durable ; en revanche, le plastique fondu enrobe les grains de sable et, en se solidifiant au refroidissement, les soude fermement les uns aux autres. C'est bien le plastique qui joue le rôle de liant, à la place du ciment.
\end{experience}

\courssub{Précautions et sécurité}

\begin{attention}[Fusion du plastique : une opération à encadrer strictement]
La fusion du plastique dégage des \textbf{fumées} potentiellement irritantes ou toxiques, et le mélange chaud peut provoquer des \textbf{brûlures graves}. Par conséquent :
\begin{itemize}
\item l'opération doit se faire \textbf{en espace ventilé} (à l'extérieur ou dans un local largement ouvert) ;
\item l'opérateur porte des \textbf{équipements de protection} : gants épais, masque, lunettes, vêtements couvrants ;
\item le chauffage doit rester \textbf{contrôlé} : ne jamais surchauffer le plastique jusqu'à le brûler ;
\item cette manipulation \textbf{ne doit jamais être réalisée par des élèves sans encadrement} : seule une démonstration par l'enseignant ou un technicien est envisageable.
\end{itemize}
Piège de rédaction fréquent : écrire que le plastique « bout » ou « s'évapore ». Non : il \emph{fond} (changement d'état solide → pâteux), puis \emph{se solidifie} au refroidissement. Emploie le vocabulaire exact des changements d'état.
\end{attention}

\begin{exoresolu}[Doser le mélange pour une production de pavés]
Une coopérative de jeunes de Douala fabrique des pavés avec un dosage de \num{2,5} volumes de sable pour 1 volume de plastique fondu. Pour une production, elle dispose de \SI{40}{L} de plastique broyé (dont le volume est conservé après fusion, en première approximation). Quel volume de sable doit-elle prévoir ? Quel volume total de mélange obtient-elle (en supposant les volumes additifs) ?
\tcblower
\textbf{Solution.} Le dosage impose : $V_{\text{sable}} = \num{2,5} \times V_{\text{plastique}}$.
\begin{align*}
V_{\text{sable}} &= \num{2,5} \times 40 = \SI{100}{L}\\
V_{\text{total}} &= V_{\text{plastique}} + V_{\text{sable}} = 40 + 100 = \SI{140}{L}
\end{align*}
La coopérative doit prévoir \SI{100}{L} de sable propre et sec, et obtient environ \SI{140}{L} de mélange. Avec des moules de \SI{2,4}{L} (\SI{20}{cm} $\times$ \SI{10}{cm} $\times$ \SI{12}{cm} remplis à \SI{6}{cm} de hauteur utile, soit \SI{1,2}{L} par pavé), elle peut produire environ $\dfrac{140}{\num{1,2}} \approx 116$ pavés. Remarque : dans la pratique, le compactage réduit légèrement le volume apparent, ce qui améliore encore la densité et la résistance des pavés.
\end{exoresolu}

\begin{savaistu}[Des pavés qui créent des emplois en Afrique]
En Afrique, la transformation des déchets plastiques en pavés est devenue une véritable filière économique. Au \textbf{Cameroun}, des jeunes entrepreneurs, à Douala comme à Yaoundé, collectent les sachets et emballages usagés pour fabriquer des pavés destinés aux cours, trottoirs et parkings. En \textbf{Côte d'Ivoire}, une entreprise fondée par une jeune ingénieure transforme chaque jour des tonnes de plastique en pavés pour construire des salles de classe. Au \textbf{Kenya}, une entrepreneure a développé des pavés plastique-sable réputés plus résistants que le béton, en recyclant les déchets des industries locales. Ces initiatives montrent que la biotechnologie environnementale n'est pas qu'une affaire de laboratoire : elle peut naître d'un atelier de quartier, assainir la ville et créer des emplois pour la jeunesse. Et si ton prochain projet de club scientifique était de paver la cour de ton lycée ?
\end{savaistu}

\begin{aretenir}[L'essentiel de la section]
\begin{itemize}
\item Le pavé plastique est un matériau composite : le \textbf{plastique fondu} sert de liant (à la place du ciment) et le \textbf{sable} de granulat.
\item La fabrication comporte cinq grandes étapes : \textbf{tri/lavage → découpe/broyage → fusion avec brassage → mélange avec le sable et malaxage → moulage, compactage, refroidissement et démoulage}.
\item Les pavés obtenus sont \textbf{résistants, imperméables, légers et durables} ; ils conviennent aux cours, trottoirs et parkings, souvent à moindre coût que les pavés en ciment.
\item La fusion dégage des fumées : \textbf{ventilation, équipements de protection et encadrement} sont indispensables.
\end{itemize}
\end{aretenir}

\coursec{Avantages de la valorisation des déchets et sensibilisation de la communauté}

Dans la section précédente, nous avons vu comment les déchets plastiques, broyés puis mélangés à du sable, pouvaient être transformés en pavés solides et vendables. Tu as donc découvert, au fil de cette leçon, que le papier usagé peut devenir source d'énergie ou nouveau papier, et que le plastique peut renaître en bouteilles ou en pavés. Mais à quoi servent exactement toutes ces techniques ? Pourquoi une commune, une école ou une famille devrait-elle s'y intéresser ? C'est la question à laquelle cette dernière section répond : nous allons évaluer les \cle{avantages} de la valorisation des déchets pour l'\cle{environnement}, pour l'\cle{économie} et pour la \cle{société}, puis apprendre à \cle{sensibiliser} notre communauté, car aucune technique ne fonctionne sans l'adhésion des populations.

\courssub{Observation de départ : deux quartiers, deux réalités}

Imaginons deux quartiers d'une même ville camerounaise. Dans le premier, les papiers et les plastiques sont jetés dans les caniveaux : à la saison des pluies, les caniveaux se bouchent, l'eau stagne dans les rues, les moustiques pullulent et les cas de paludisme et de choléra augmentent. Dans le second quartier, une association de jeunes collecte les papiers et les plastiques, les trie et les vend à des recycleurs : les caniveaux restent propres, les rues ne sont pas inondées, et les jeunes gagnent un revenu.

\begin{experience}[Comparer pour comprendre : démarche d'investigation]
\textbf{Problème :} la valorisation des déchets améliore-t-elle réellement les conditions de vie ?

\textbf{Hypothèse :} un quartier où les déchets sont triés et valorisés présente moins de nuisances (inondations, maladies) qu'un quartier où ils sont abandonnés.

\textbf{Protocole :} on choisit deux quartiers comparables (même population, même relief). Dans le quartier A, on organise pendant 6 mois le tri et la collecte des papiers et plastiques ; dans le quartier B, on ne change rien. On mesure avant et après : le nombre de points d'eau stagnante après une pluie, le volume de déchets dans les caniveaux (en kg), et le nombre de cas de paludisme déclarés au centre de santé.

\textbf{Résultats attendus :} dans le quartier A, le volume de déchets dans les caniveaux diminue fortement, l'eau s'écoule mieux après les pluies, et les cas de paludisme reculent ; dans le quartier B, rien ne change.

\textbf{Conclusion :} la valorisation des déchets réduit les nuisances environnementales et sanitaires. L'hypothèse est vérifiée.
\end{experience}

Cette démarche montre que les avantages de la valorisation ne sont pas des slogans : ce sont des effets \emph{mesurables} et \emph{démontrables}.

\courssub{Les avantages pour l'environnement}

\begin{propriete}[Avantages de la valorisation des déchets (savoir admis)]
La valorisation des déchets \textbf{réduit la pollution} (sols, eau, air), \textbf{préserve les ressources naturelles} (arbres, pétrole) et \textbf{génère des revenus}. Elle transforme un problème (le déchet) en solution (matière première et énergie).
\end{propriete}

Détaillons ces avantages environnementaux :

\begin{itemize}
\item \textbf{Réduction des déchets dans la nature et les caniveaux.} Chaque kilogramme de papier ou de plastique collecté est un kilogramme qui ne finit pas dans la rue, la brousse ou le caniveau. Or un sac plastique abandonné peut mettre plusieurs centaines d'années à se dégrader.
\item \textbf{Diminution des inondations.} Les caniveaux débarrassés des plastiques laissent s'écouler les eaux de pluie : les inondations urbaines, fréquentes à Douala ou à Yaoundé pendant la saison des pluies, s'atténuent.
\item \textbf{Diminution des maladies.} Moins d'eau stagnante signifie moins de moustiques (vecteurs du paludisme) et moins de contamination des eaux (choléra, typhoïde).
\item \textbf{Moindre pollution des sols, de l'eau et de l'air.} Le plastique qui se décompose lentement libère des substances toxiques dans le sol ; brûlé à l'air libre, il pollue l'air. Le recyclage évite ces deux pollutions.
\item \textbf{Économie des ressources naturelles.} Recycler 1 tonne de papier évite d'abattre environ 17 arbres ; recycler le plastique économise le pétrole, matière première de sa fabrication. La valorisation ménage donc les forêts et les ressources fossiles.
\end{itemize}

\courssub{Les avantages économiques et sociaux}

\textbf{Avantages économiques.} La valorisation crée toute une chaîne d'activités rémunératrices :

\begin{itemize}
\item \textbf{Création d'emplois et de revenus :} collecteurs, trieurs, recycleurs, artisans fabricants de pavés ou de papier recyclé gagnent leur vie grâce aux déchets ;
\item \textbf{Production de biens vendables :} pavés, papier recyclé, bouteilles recyclées, briquettes de papier se vendent sur les marchés locaux ;
\item \textbf{Réduction des coûts d'assainissement :} les communes dépensent moins pour ramasser et enfouir les ordures lorsqu'une partie est valorisée ;
\item \textbf{Économie de matières premières importées :} utiliser du papier et du plastique recyclés réduit les importations de matières premières, ce qui épargne des devises au pays.
\end{itemize}

\textbf{Avantages sociaux.} Les quartiers deviennent plus propres et plus sains, ce qui améliore la qualité de vie de tous. Les activités de collecte et de recyclage sont aussi des occasions d'\cle{éducation environnementale} et d'\cle{entrepreneuriat des jeunes} : un jeune qui fabrique des pavés à partir de plastique est à la fois un protecteur de l'environnement et un entrepreneur.

\begin{exemplebox}[Un club environnemental scolaire en action]
Le club environnemental d'un lycée peut organiser une \textbf{journée de collecte} de papiers et de plastiques dans les salles de classe et la cour. Les papiers récoltés servent à fabriquer du \textbf{papier recyclé} (trempage, broyage en pâte, désencrage si nécessaire, étalement sur tamis, séchage, pressage), qui peut devenir papier-cadeau ou serviettes pour la fête de fin d'année. Les plastiques sont revendus à un artisan fabricant de pavés. Le club gagne de l'argent pour ses activités, et tout le lycée apprend le geste du tri.
\end{exemplebox}

\courssub{Le cercle vertueux de la valorisation}

Tous ces avantages s'enchaînent logiquement : le tri permet le recyclage, le recyclage produit des biens, les biens vendus rapportent des revenus, et un environnement sain motive à continuer le tri. C'est un \cle{cercle vertueux}.

\begin{popfigure}
\begin{tikzpicture}[scale=1, every node/.style={align=center}]
\foreach \i/\txt in {0/{Tri des\\déchets}, 72/{Recyclage et\\transformation}, 144/{Produits\\(pavés, papier,\\bouteilles)}, 216/{Revenus et\\emplois}, 288/{Environnement\\sain}} {
  \node[draw=popgreen, fill=popgreenL, rounded corners=3pt, minimum width=2.6cm, minimum height=1.1cm, font=\small] (n\i) at (\i:3.2cm) {\txt};
}
\foreach \a/\b in {0/72, 72/144, 144/216, 216/288, 288/0} {
  \draw[-{Stealth[length=3mm]}, popdark, thick] (n\a) to[bend left=10] (n\b);
}
\node[font=\small\itshape, text=popmuted, align=center] at (0,0){Cercle vertueux\\de la valorisation};
\end{tikzpicture}
\legende{Le cercle vertueux de la valorisation des déchets : chaque étape entraîne la suivante. Le tri initial déclenche toute la chaîne ; les revenus et l'environnement sain obtenus motivent la population à poursuivre le tri.}
\end{popfigure}

\begin{savaistu}[Le tri à la source : le premier maillon de la chaîne]
Le \textbf{tri à la source} (à la maison, à l'école, au marché) est le premier maillon \emph{indispensable} de toute la chaîne de valorisation. Un papier souillé de restes de nourriture ou un plastique mélangé à des ordures ménagères devient très difficile à recycler. Séparer dès le départ les papiers, les plastiques et les ordures dans des récipients différents rend tout le reste possible. Sans tri à la source, pas de recyclage efficace !
\end{savaistu}

\courssub{Informer et sensibiliser sa communauté}

Les techniques de valorisation ne servent à rien si la population continue de jeter ses déchets n'importe où. Il faut donc \cle{informer} et \cle{sensibiliser} la communauté sur trois messages essentiels :

\begin{enumerate}
\item \textbf{Trier les déchets à la source} (maison, école, marché) : séparer papiers, plastiques et ordures ;
\item \textbf{Ne jamais jeter les déchets dans les caniveaux}, les rivières ou la brousse ;
\item \textbf{Soutenir les filières de recyclage} : donner ou vendre ses déchets triés aux collecteurs et recycleurs locaux.
\end{enumerate}

\begin{methode}[Comment sensibiliser sa communauté en quatre étapes]
\begin{enumerate}
\item \textbf{Identifier le problème local :} observer son quartier ou son école (caniveaux bouchés, tas d'ordures, plastiques brûlés) et recueillir quelques données simples (photos, volumes, témoignages) ;
\item \textbf{Expliquer simplement les techniques :} montrer, avec des mots simples et des exemples locaux, comment le papier et le plastique peuvent être recyclés (papier recyclé, briquettes, pavés, bouteilles) ;
\item \textbf{Montrer les avantages :} insister sur ce qui touche directement les gens — moins d'inondations, moins de maladies, des revenus possibles pour les jeunes ;
\item \textbf{Proposer des gestes concrets :} tri à la maison, point de collecte à l'école ou au marché, participation à une journée de nettoyage.
\end{enumerate}
\end{methode}

Les \textbf{moyens de sensibilisation} sont nombreux et accessibles : les \emph{causeries éducatives} (réunions de quartier, réunions de parents d'élèves), les \emph{affiches} placardées à l'école et au marché, les émissions des \emph{radios communautaires} en langues locales, les \emph{clubs environnementaux scolaires}, et surtout les \emph{démonstrations pratiques} — rien ne convainc mieux qu'un pavé fabriqué sous les yeux du public à partir de sacs plastiques !

\begin{attention}[La valorisation n'est pas une baguette magique]
La valorisation \textbf{ne supprime pas toutes les pollutions}. L'incinération du papier pour produire de l'énergie dégage des fumées ; la fusion du plastique peut rejeter des gaz toxiques ; le recyclage du papier consomme de l'eau et des produits de désencrage. Ces activités doivent donc être \textbf{encadrées et contrôlées} (ventilation, traitement des rejets, respect des normes). Dire au Bac que « le recyclage ne pollue pas du tout » est une erreur : il pollue \emph{moins} que l'abandon des déchets, à condition d'être bien conduit.
\end{attention}

\courssub{Bilan de la leçon : tableau récapitulatif}

\begin{popfigure}
\begin{center}
\begin{tabularx}{\linewidth}{|l|X|X|X|}
\hline
\rowcolor{popblueL} \textbf{Déchet} & \textbf{Technique} & \textbf{Produit obtenu} & \textbf{Avantage} \\
\hline
Papiers & Incinération contrôlée (combustion) & Chaleur, énergie (briquettes de papier) & Économie de bois et de charbon ; moins de déchets \\
\hline
Papiers & Recyclage (trempage, broyage en pâte, désencrage, étalement, séchage, pressage) & Serviettes jetables, papier hygiénique, papier-cadeau & Économie d'arbres ; revenus des artisans \\
\hline
Plastiques & Tri, lavage, broyage, fusion, moulage & Bouteilles recyclées & Économie de pétrole ; moins de plastique dans la nature \\
\hline
Plastiques & Broyage, mélange au sable, chauffage, moulage & Pavés & Rues et cours pavées à moindre coût ; caniveaux débarrassés \\
\hline
\end{tabularx}
\end{center}
\legende{Tableau récapitulatif des techniques de valorisation du papier et du plastique étudiées dans cette leçon : pour chaque déchet, une technique, un produit vendable et un avantage pour l'environnement et l'économie.}
\end{popfigure}

\begin{exoresolu}[Exercice résolu : argumenter pour convaincre]
\textbf{Énoncé.} Le maire de ta commune hésite à soutenir un projet de recyclage des papiers et plastiques porté par des jeunes. Rédige, en cinq arguments scientifiquement fondés, une plaidoirie pour le convaincre.

\tcblower
\textbf{Solution détaillée.}
\begin{enumerate}
\item \textbf{Argument sanitaire :} les plastiques bouchent les caniveaux ; l'eau stagnante favorise les moustiques et le paludisme. Recycler réduit les eaux stagnantes et donc les maladies.
\item \textbf{Argument environnemental :} le plastique met des centaines d'années à se dégrader et pollue sols et eaux ; le recyclage du papier épargne environ 17 arbres par tonne recyclée.
\item \textbf{Argument économique :} le projet crée des emplois (collecte, tri, fabrication) et des produits vendables (pavés, papier recyclé), donc des revenus pour les jeunes.
\item \textbf{Argument financier pour la commune :} moins de déchets à ramasser et à enfouir signifie une réduction des coûts d'assainissement.
\item \textbf{Argument social et éducatif :} le projet sensibilise la population au tri à la source et à la propreté, ce qui améliore durablement le cadre de vie.
\textbf{Conclusion :} en soutenant ce projet, la commune gagne sur tous les plans : santé, environnement, économie et cohésion sociale.
\end{enumerate}
\end{exoresolu}

\begin{aretenir}[L'essentiel de la section]
\begin{itemize}
\item La valorisation des déchets \textbf{réduit la pollution} (sols, eau, air), \textbf{préserve les ressources naturelles} (arbres, pétrole) et \textbf{génère des revenus} ;
\item Elle diminue les \textbf{inondations} et les \textbf{maladies} (paludisme, choléra) en débarrassant les caniveaux des plastiques ;
\item Elle crée des \textbf{emplois} (collecteurs, recycleurs, artisans) et des \textbf{produits vendables} (pavés, papier recyclé, bouteilles, briquettes) ;
\item Le \textbf{tri à la source} est le premier maillon indispensable de toute la chaîne ;
\item Sensibiliser sa communauté se fait en quatre étapes : identifier le problème, expliquer les techniques, montrer les avantages, proposer des gestes concrets ;
\item La valorisation doit être \textbf{encadrée}, car elle ne supprime pas toutes les pollutions.
\end{itemize}
\end{aretenir}

Tu es maintenant capable non seulement d'expliquer les techniques de transformation et de recyclage du papier et du plastique, mais aussi de défendre leur intérêt devant ta communauté. C'est là tout l'esprit de la biotechnologie au service du développement durable : transformer les problèmes de l'environnement en opportunités pour tous.

\newpage

\coursec{Activité d'intégration}

\coursec{Transformation et recyclage des déchets}

\courssub{Objectifs de l'activité}

À la fin de cette activité d'intégration, tu seras capable de :

\begin{itemize}
\item identifier et classer les déchets valorisables produits dans un milieu scolaire (démarche du tri) ;
\item décrire, sous forme d'organigrammes, les techniques de valorisation du papier (production d'énergie, recyclage en serviettes jetables, papier hygiénique, papier-cadeau) et des plastiques (bouteilles recyclées, pavés) ;
\item exploiter des données chiffrées (masses, rendements, coûts, prix de vente) pour estimer une production mensuelle et un revenu potentiel ;
\item comparer les avantages et les limites de plusieurs filières de valorisation ;
\item rédiger un message de sensibilisation destiné à ta communauté sur la nécessité de trier, transformer et recycler les déchets.
\end{itemize}

\courssub{Situation}

Le lycée bilingue de Nkol-Eton, dans la banlieue de Yaoundé, accueille \num{1200} élèves. Chaque semaine, la cour et les salles de classe accumulent des montagnes de déchets : feuilles de cahiers usagées, vieux journaux, cartons, sachets plastiques, bouteilles d'eau minérale et de boissons gazeuses. Faute de tri, ces déchets sont brûlés à l'air libre derrière le terrain de sport, ce qui provoque des fumées âcres dont se plaignent les riverains. Le club environnement du lycée, en partenariat avec une coopérative de jeunes du quartier qui fabrique des pavés en plastique-sable et un petit atelier de papier recyclé, décide de lancer le projet « Zéro déchet ». Tu es membre de ce club.

\textbf{Document 1 : Quantités de déchets produits par semaine dans le lycée}

\begin{center}
\begin{tabularx}{\linewidth}{|l|X|X|}\hline
\rowcolor{popblueL} \textbf{Type de déchet} & \textbf{Masse par semaine (kg)} & \textbf{Exemples} \\ \hline
Papiers et cartons propres & 60 & feuilles de cahiers, journaux, cartons d'emballage \\ \hline
Plastiques durs (bouteilles PET) & 25 & bouteilles d'eau et de boissons gazeuses \\ \hline
Plastiques souples (sachets) & 15 & sachets d'emballage, emballages de beignets \\ \hline
Déchets organiques et autres & 40 & épluchures, restes de repas \\ \hline
\end{tabularx}
\end{center}

\textbf{Document 2 : Techniques artisanales et rendements}

\begin{itemize}
\item \emph{Atelier de papier recyclé} : le papier trié est déchiqueté, trempé dans l'eau pendant 24 h, réduit en pâte (défibrage), éventuellement désencré et blanchi, puis étalé en feuilles sur des tamis (couchage), pressé et séché au soleil. Rendement : 1 kg de papiers-cartons propres donne environ 0,8 kg de pâte, soit \textbf{40 feuilles} de papier recyclé (format A4) utilisables comme papier-cadeau ou serviettes jetables après découpe.
\item \emph{Unité de pavés plastique-sable} : les plastiques (sachets et bouteilles broyées) sont lavés, séchés, fondus dans une marmite chauffée, mélangés à du sable (proportion 1 volume de plastique fondu pour 3 volumes de sable), puis versés dans des moules et refroidis. Rendement : 1 kg de plastique permet de fabriquer \textbf{5 pavés} de dimensions $20 \times 10 \times 6$ cm.
\item \emph{Valorisation énergétique} : le papier non recyclable (souillé) peut être brûlé dans un incinérateur équipé d'un système de récupération de chaleur ; la chaleur dégagée produit de la vapeur qui peut chauffer de l'eau ou faire tourner un alternateur. Le pouvoir calorifique du papier sec est d'environ 15 MJ/kg.
\end{itemize}

\textbf{Document 3 : Coûts et prix de vente}

\begin{center}
\begin{tabularx}{\linewidth}{|l|X|X|}\hline
\rowcolor{popblueL} \textbf{Produit valorisé} & \textbf{Coût de production unitaire (FCFA)} & \textbf{Prix de vente unitaire (FCFA)} \\ \hline
Pavé plastique-sable & 50 & 100 \\ \hline
Feuille de papier recyclé (papier-cadeau) & 15 & 25 \\ \hline
\end{tabularx}
\end{center}

Le mois scolaire compte 4 semaines. Le club dispose d'un petit budget de démarrage et souhaite présenter au proviseur et aux parents d'élèves un plan d'action chiffré, accompagné d'une campagne de sensibilisation.

\courssub{Consignes}

Par groupes de quatre élèves, réponds aux questions suivantes :

\begin{enumerate}
\item À partir du document 1, identifie les déchets valorisables du lycée et classe-les en deux grandes catégories (déchets papiers, déchets plastiques). Précise le devenir possible de chaque catégorie.
\item Construis deux organigrammes : l'un décrivant les étapes du recyclage du papier en serviettes jetables ou papier-cadeau, l'autre décrivant la transformation des plastiques en pavés. Indique à chaque étape l'opération réalisée.
\item Calcule la masse mensuelle de papiers et de plastiques valorisables. En déduis le nombre de feuilles de papier recyclé et le nombre de pavés produisibles par mois.
\item Calcule le revenu mensuel brut puis le bénéfice mensuel (recette moins coûts de production) pour chaque filière. Quelle filière est la plus rentable ?
\item Le proviseur propose de brûler le papier souillé (10 kg par semaine) dans un incinérateur pour chauffer l'eau de la cantine. Calcule l'énergie thermique récupérable par mois. Cite une précaution indispensable.
\item Rédige une affiche ou un message de sensibilisation (10 lignes maximum) destiné aux élèves et aux habitants du quartier, expliquant pourquoi il faut trier, transformer et recycler les déchets.
\item Discussion critique : cite deux limites ou deux risques de ces techniques artisanales (fumées, hygiène, équipements) et propose une solution pour chacun.
\end{enumerate}

\courssub{Ressources à mobiliser}

\begin{aretenir}[Ce qu'il faut avoir en tête]
\begin{itemize}
\item Le \cle{recyclage} réintroduit un déchet dans le cycle de production d'un produit de même nature (papier $\rightarrow$ papier recyclé, bouteilles PET $\rightarrow$ nouvelles bouteilles) ; la \cle{transformation} donne un produit nouveau (plastiques $\rightarrow$ pavés plastique-sable) ; la \cle{valorisation énergétique} convertit le déchet en énergie (chaleur, électricité).
\item Chaîne du papier recyclé : tri $\rightarrow$ déchiquetage $\rightarrow$ trempage $\rightarrow$ défibrage (pâte) $\rightarrow$ désencrage/blanchiment $\rightarrow$ couchage sur tamis $\rightarrow$ pressage $\rightarrow$ séchage $\rightarrow$ découpe.
\item Chaîne des pavés : tri $\rightarrow$ lavage $\rightarrow$ séchage $\rightarrow$ broyage $\rightarrow$ fusion $\rightarrow$ mélange plastique-sable (1 volume plastique / 3 volumes sable) $\rightarrow$ moulage $\rightarrow$ refroidissement $\rightarrow$ démoulage.
\item Calculs utiles : production = masse $\times$ rendement ; recette = quantité $\times$ prix de vente ; bénéfice = recette $-$ coûts ; énergie = masse $\times$ pouvoir calorifique ($E = m \times PC$).
\end{itemize}
\end{aretenir}

\begin{savaistu}[Contexte camerounais : initiatives locales]
Au Cameroun, la gestion des déchets urbains est un enjeu majeur. À Yaoundé et Douala, des coopératives de jeunes (ex. \emph{Green Youth}, \emph{Plastic Man}) collectent les sachets et bouteilles pour fabriquer des pavés écologiques qui revêtent les cours d'écoles et les allées de quartiers. HYSACAM expérimente la valorisation énergétique du papier/carton dans certains centres de transit. Le recyclage du papier en serviettes ou papier hygiénique reste artisanal mais se développe dans les ateliers de formation professionnelle (ex. CFP de Mbalmayo). Ces filières créent des emplois verts et réduisent la pollution plastique qui bouche les caniveaux et aggrave les inondations.
\end{savaistu}

\courssub{Démarche et corrigé commenté}

\begin{exoresolu}[Projet « Zéro déchet » : résolution complète]
Énoncé : répondre aux sept consignes à partir des documents 1, 2 et 3.
\tcblower
\textbf{1. Identification et classement des déchets valorisables.}

Déchets papiers (60 kg/semaine) : feuilles, journaux, cartons propres $\rightarrow$ recyclage en papier-cadeau, serviettes jetables, papier hygiénique ; fraction souillée (10 kg/semaine, consigne 5) $\rightarrow$ valorisation énergétique. Déchets plastiques ($25 + 15 = 40$ kg/semaine) : bouteilles PET (plastiques durs) $\rightarrow$ recyclage industriel en nouvelles bouteilles (non traité ici) ou transformation artisanale en pavés ; sachets (plastiques souples) $\rightarrow$ transformation en pavés. Les déchets organiques (40 kg/semaine) relèvent du compostage (hors programme de cette leçon).

\textbf{2. Organigrammes des filières de valorisation.}

\begin{popfigure}
\begin{tikzpicture}[
    node distance=1.2cm and 0.5cm,
    box/.style={rectangle, rounded corners, draw=popblue, thick, fill=popblueL, text width=2.8cm, align=center, minimum height=1.2cm, font=\small},
    arrow/.style={-Stealth, thick, popdark},
    start/.style={ellipse, draw=popgreen, thick, fill=popgreenL, text width=2.5cm, align=center, minimum height=1.2cm, font=\small\bfseries},
    end/.style={ellipse, draw=poporange, thick, fill=poporangeL, text width=2.5cm, align=center, minimum height=1.2cm, font=\small\bfseries}
]
% Paper chain
\node[start, align=center] (startP){Papiers/cartons\\propres triés};
\node[box, below=of startP] (d1) {Déchiquetage};
\node[box, below=of d1] (d2) {Trempage (24 h)};
\node[box, below=of d2] (d3) {Défibrage $\rightarrow$ Pâte};
\node[box, below=of d3] (d4) {Désencrage / Blanchiment};
\node[box, below=of d4] (d5) {Couchage sur tamis};
\node[box, below=of d5] (d6) {Pressage};
\node[box, below=of d6] (d7) {Séchage au soleil};
\node[end, below=of d7, align=center] (endP){Papier recyclé\\(serviettes, cadeau)};

\draw[arrow] (startP) -- (d1);
\draw[arrow] (d1) -- (d2);
\draw[arrow] (d2) -- (d3);
\draw[arrow] (d3) -- (d4);
\draw[arrow] (d4) -- (d5);
\draw[arrow] (d5) -- (d6);
\draw[arrow] (d6) -- (d7);
\draw[arrow] (d7) -- (endP);

% Plastic chain (shifted right)
\begin{scope}[xshift=9cm]
\node[start, align=center] (startPl){Plastiques triés\\(bouteilles, sachets)};
\node[box, below=of startPl] (p1) {Lavage};
\node[box, below=of p1] (p2) {Séchage};
\node[box, below=of p2] (p3) {Broyage};
\node[box, below=of p3] (p4) {Fusion};
\node[box, below=of p4, align=center] (p5){Mélange + Sable\\(1 vol. pl. / 3 vol. sab.)};
\node[box, below=of p5] (p6) {Moulage};
\node[box, below=of p6] (p7) {Refroidissement};
\node[end, below=of p7, align=center] (endPl){Pavés\\plastique-sable};

\draw[arrow] (startPl) -- (p1);
\draw[arrow] (p1) -- (p2);
\draw[arrow] (p2) -- (p3);
\draw[arrow] (p3) -- (p4);
\draw[arrow] (p4) -- (p5);
\draw[arrow] (p5) -- (p6);
\draw[arrow] (p6) -- (p7);
\draw[arrow] (p7) -- (endPl);
\end{scope}
\end{tikzpicture}
\legende{Organigrammes des deux filières de valorisation artisanale : recyclage du papier (à gauche) et transformation des plastiques en pavés (à droite).}
\end{popfigure}

\textbf{3. Production mensuelle.}

Masse mensuelle de papiers propres : $60 \times 4 = 240$ kg.
Masse mensuelle de plastiques (durs + souples) : $(25 + 15) \times 4 = 160$ kg.

Feuilles de papier recyclé : $240 \text{ kg} \times 40 \text{ feuilles/kg} = \num{9600}$ feuilles par mois.
Pavés : $160 \text{ kg} \times 5 \text{ pavés/kg} = 800$ pavés par mois.

\textbf{4. Revenus et bénéfices mensuels.}

\begin{align*}
\text{Filière Papier :}\quad
\text{Recette} &= 9600 \times 25 = \num{240000} \text{ FCFA}\\
\text{Coûts} &= 9600 \times 15 = \num{144000} \text{ FCFA}\\
\text{Bénéfice} &= \num{240000} - \num{144000} = \num{96000} \text{ FCFA}\\[1em]
\text{Filière Pavés :}\quad
\text{Recette} &= 800 \times 100 = \num{80000} \text{ FCFA}\\
\text{Coûts} &= 800 \times 50 = \num{40000} \text{ FCFA}\\
\text{Bénéfice} &= \num{80000} - \num{40000} = \num{40000} \text{ FCFA}
\end{align*}

\begin{popfigure}
\begin{tikzpicture}
\begin{axis}[
    ybar,
    bar width=25pt,
    width=0.7\linewidth,
    height=6cm,
    symbolic x coords={Papier recyclé, Pavés plastique-sable},
    xtick=data,
    ylabel={Bénéfice mensuel (kFCFA)},
    ymin=0, ymax=110,
    ytick={0,20,40,60,80,100},
    nodes near coords,
    nodes near coords align={vertical},
    every node near coord/.append style={font=\small\bfseries, popdark},
    tick label style={font=\small},
    label style={font=\small},
    axis lines*=left,
    enlarge x limits=0.3,
    fill=popblue!70,
]
\addplot coordinates {(Papier recyclé,96) (Pavés plastique-sable,40)};
\end{axis}
\end{tikzpicture}
\legende{Comparaison des bénéfices mensuels des deux filières. La filière papier est plus rentable mais traite un flux de déchets différent.}
\end{popfigure}

La filière papier recyclé est la plus rentable (\num{96000} FCFA contre \num{40000} FCFA de bénéfice mensuel), mais les deux filières sont complémentaires car elles valorisent des déchets distincts. Bénéfice total potentiel : \num{136000} FCFA par mois.

\textbf{5. Valorisation énergétique du papier souillé.}

Masse mensuelle de papier souillé : $10 \times 4 = 40$ kg.
Énergie thermique récupérable :
\[ E = m \times PC = 40 \text{ kg} \times 15 \text{ MJ/kg} = 600 \text{ MJ par mois.} \]
\emph{Précaution indispensable :} l'incinération doit se faire dans un appareil fermé (incinérateur), équipé d'un système de traitement des fumées (filtres, lavage des gaz), à distance des habitations. La combustion à l'air libre est interdite car elle émet des polluants toxiques (dioxines, furanes, particules fines, CO).

\textbf{6. Message de sensibilisation (exemple — 9 lignes).}

\begin{center}
\fbox{\begin{minipage}{0.9\linewidth}
\textbf{ENSEMBLE POUR UN LYCÉE PROPRE ET VERT !}

Chers élèves, chers voisins, chaque semaine 100 kg de papiers et plastiques quittent notre établissement. Brûlés, ils empoisonnent l'air ; triés, ils deviennent papier-cadeau, serviettes et pavés pour nos rues, tout en finançant nos activités. \textbf{3 gestes simples :} \textbullet\ Papier $\rightarrow$ poubelle bleue ; \textbullet\ Plastique $\rightarrow$ poubelle jaune ; \textbullet\ Zéro brûlage à l'air libre. Trier, transformer, recycler : protégeons notre santé, embellissons Nkol-Eton, créons de la richesse. Rejoignez le club « Zéro Déchet » !
\end{minipage}}
\end{center}

\textbf{7. Limites, risques et solutions.}

\begin{itemize}
\item \emph{Fusion du plastique (pavés) :} dégage des fumées toxiques (monomères, additifs, éventuel PVC) et risque de brûlures graves. \textbf{Solution :} travailler en zone ventilée ou sous hotte d'extraction, porter masque à cartouches (vapeurs organiques), gants thermiques, tablier ; interdire formellement le PVC (reconnaissable au test du fil de cuivre vert).
\item \emph{Hygiène du papier recyclé :} les fibres issues de papiers souillés ou mélangés à des déchets organiques portent des germes pathogènes ; l'eau de trempage stagne (moustiques). \textbf{Solution :} tri rigoureux à la source (séparer propre/souillé), désinfection de la pâte (eau de Javel diluée, rinçage), séchage rapide au soleil (UV), réserver le papier recyclé aux usages non alimentaires/non corporels (papier-cadeau, emballage) sauf norme sanitaire validée.
\item \emph{Équipements artisanaux :} marmites sur feu nu, moules métalliques chauds $\rightarrow$ risques d'incendie et de brûlures. \textbf{Solution :} zone de feu délimitée, extincteur à poudre ABC à proximité, formation aux gestes de premiers secours, manipulation des moules uniquement avec pinces/gants.
\end{itemize}
\end{exoresolu}

\begin{attention}[Erreurs fréquentes à éviter au Bac]
\begin{itemize}
\item Oublier de convertir les masses hebdomadaires en masses mensuelles (facteur 4) avant d'appliquer les rendements.
\item Confondre \textbf{recette} (chiffre d'affaires) et \textbf{bénéfice} (recette $-$ coûts de production) ; le bénéfice seul mesure la rentabilité.
\item Additionner les masses de papiers et de plastiques pour un calcul global : les rendements (feuilles/kg vs pavés/kg) et les prix unitaires diffèrent, chaque filière se calcule séparément.
\item Présenter l'incinération comme une solution « propre » sans mentionner l'obligation de traitement des fumées et de cendres (métaux lourds).
\item Rédiger un message de sensibilisation trop long, technique ou moralisateur : il doit être court, positif, actionnable (verbes d'action) et adapté au public visé.
\end{itemize}
\end{attention}

\courssub{Bilan de l'activité}

Cette activité t'a permis de mobiliser l'ensemble des savoir-faire de la leçon dans une situation authentique. Tu as d'abord \cle{classé} les déchets d'un milieu réel et associé à chaque catégorie une filière de valorisation adaptée (recyclage, transformation, énergie). Tu as ensuite \cle{décrit} sous forme d'organigrammes les techniques de recyclage du papier et de transformation des plastiques, ce qui t'a obligé à ordonner logiquement les étapes (tri, lavage, défibrage ou fusion, moulage ou couchage, séchage). L'exploitation des données chiffrées t'a montré que la valorisation n'est pas seulement un geste écologique : avec \num{136000} FCFA de bénéfice mensuel potentiel, c'est aussi une \cle{activité économique} créatrice de revenus, comme le prouvent les coopératives de jeunes qui pavent les cours et les rues de nombreux quartiers de Yaoundé et de Douala. Enfin, la rédaction du message de sensibilisation et la discussion critique sur les limites (fumées toxiques, hygiène, sécurité des équipements) t'ont appris qu'informer sa communauté suppose d'être honnête sur les risques et rigoureux sur les solutions. Trier, transformer, recycler : trois verbes qui résument une attitude citoyenne et scientifique face aux problèmes de pollution de l'environnement.

\newpage

\coursec{Méthodes et exercices résolus}

\coursec{Transformation et recyclage des déchets}

\courssub{Généralités : définitions et hiérarchie des déchets}

\begin{definition}[Vocabulaire fondamental]
\begin{itemize}
    \item \textbf{Déchet :} Toute substance ou objet dont le détenteur se défait ou a l'intention ou l'obligation de se défaire (Loi 96/12, Art. 2).
    \item \textbf{Valorisation :} Toute opération dont le résultat principal est que le déchet sert à une fin utile en se substituant à d'autres matières (matière ou énergie).
    \item \textbf{Recyclage (matière) :} Traitement des déchets pour en extraire des matières réutilisables (fibres de papier, polymères plastiques). La matière est conservée.
    \item \textbf{Valorisation énergétique :} Incinération avec récupération d'énergie (chaleur, électricité). La matière est détruite (oxydée).
    \item \textbf{Upcycling (Surcyclage) :} Transformation de déchets en produits de qualité ou valeur supérieure (ex. pavés en plastiques mélangés).
    \item \textbf{CSR (Combustibles Solides de Récupération) :} Déchets non dangereux préparés pour être utilisés comme combustible (papiers souillés, plastiques non recyclables, PCI \SI{\sim 15-20}{MJ/kg}).
\end{itemize}
\end{definition}

\begin{propriete}[Hiérarchie des modes de gestion des déchets (Directive 2008/98/CE, Loi camerounaise 96/12)]
L'ordre de priorité légal et environnemental est :
\begin{enumerate}
    \item \textbf{Prévention} (réduction à la source, éco-conception).
    \item \textbf{Réemploi} (utilisation à nouveau pour la même fin).
    \item \textbf{Recyclage} (valorisation matière : papier $\rightarrow$ papier, PET $\rightarrow$ rPET).
    \item \textbf{Autre valorisation} (valorisation énergétique : CSR $\rightarrow$ chaleur/électricité ; upcycling : plastiques mélangés $\rightarrow$ pavés).
    \item \textbf{Élimination} (incinération sans récupération, enfouissement en centre de stockage).
\end{enumerate}
\textbf{Principle :} On ne doit incinérer (valorisation énergétique) que les déchets qui ne peuvent pas être recyclés (matière) de façon techniquement et économiquement viable.
\end{propriete}

\begin{popfigure}
\begin{tikzpicture}[node distance=1.2cm, every node/.style={align=center, font=\small}, >=Stealth]
\node (prev) [draw=popgreen, fill=popgreenL, rounded corners, minimum width=3.5cm, minimum height=1cm, align=center]{\textbf{1. Prévention}\\(Réduction source)};
\node (reuse) [below of=prev, draw=popblue, fill=popblueL, rounded corners, minimum width=3.5cm, minimum height=1cm, align=center]{\textbf{2. Réemploi}\\(Même usage)};
\node (recycl) [below of=reuse, draw=poporange, fill=poporangeL, rounded corners, minimum width=4cm, minimum height=1cm, align=center]{\textbf{3. Recyclage (Matière)}\\(Fibres, Polymères)};
\node (recov) [below of=recycl, draw=poppink, fill=poppinkL, rounded corners, minimum width=4.5cm, minimum height=1cm, align=center]{\textbf{4. Valorisation Énergétique / Upcycling}\\(CSR $\rightarrow$ Énergie, Plastiques mélangés $\rightarrow$ Pavés)};
\node (elim) [below of=recov, draw=popdark, fill=popmuted, rounded corners, minimum width=3.5cm, minimum height=1cm, align=center]{\textbf{5. Élimination}\\(Enfouissement, Incinération simple)};
\draw[->, thick, popdark] (prev) -- (reuse);
\draw[->, thick, popdark] (reuse) -- (recycl);
\draw[->, thick, popdark] (recycl) -- (recov);
\draw[->, thick, popdark] (recov) -- (elim);
\node[left=0.8cm of prev, font=\small\bfseries, text width=2.5cm, align=right, rotate=90, anchor=south] {Priorité environnementale et légale croissante $\uparrow$};
\end{tikzpicture}
\legende{Hiérarchie des déchets : tout traitement doit viser le niveau le plus haut techniquement et économiquement possible. Le recyclage matière (niveau 3) est prioritaire sur la valorisation énergétique (niveau 4).}
\end{popfigure}

\courssub{Valorisation énergétique du papier : production d'énergie}

\begin{methode}[Production d'énergie à partir des déchets papiers (Valorisation énergétique / CSR)]
\textbf{Objectif :} Expliquer la technique de production d'énergie (chaleur, électricité) par incinération des papiers non recyclables ou en fin de vie (Combustibles Solides de Récupération - CSR).
\begin{enumerate}
    \item \textbf{Tri et préparation :} Séparer les papiers souillés, composites ou à fibres trop courtes. Broyage et compactage en balles ou granulés (CSR) pour homogénéiser le pouvoir calorifique (\SI{\sim 15-18}{MJ/kg}).
    \item \textbf{Combustion contrôlée :} Incinération dans un four à grille ou lit fluidisé à \SI{850-1100}{\celsius} avec excès d'air (\SI{\sim 6-10}{\%} O$_2$ résiduel) pour détruire les polluants (dioxines, furannes). Réaction exothermique : \ce{C6H10O5 (cellulose) + 6O2 -> 6CO2 + 5H2O + \Delta H}.
    \item \textbf{Récupération thermique :} Les fumées chaudes (\SI{> 850}{\celsius}) traversent une chaudière à tubes d'eau $\rightarrow$ production de vapeur surchauffée (\SI{400}{\celsius}, \SI{40}{bar}).
    \item \textbf{Conversion électrique (cogénération) :} La vapeur détend dans une turbine à vapeur couplée à un alternateur (rendement électrique \SI{\sim 20-25}{\%}). La chaleur fatale (condenseur) alimente un réseau de chaleur urbain ou des processus industriels (rendement global \SI{> 80}{\%}).
    \item \textbf{Traitement des fumées et résidus :} Épuration (chaux pour HCl/SO$_2$, charbon actif pour métaux lourds/dioxines, filtre à manches pour poussières). Valorisation des mâchefers (granulats) et traitement des REFIOM (résidus d'épuration) en centre de stockage classe 1.
\end{enumerate}
\cle{CSR}, \cle{Pouvoir Calorifique Inférieur (PCI)}, \cle{Cogénération}, \cle{Épuration des fumées}, \cle{REFIOM}
\end{methode}

\begin{exoresolu}[Bilan énergétique d'une unité de valorisation énergétique de papier]
\textbf{Énoncé :} Une unité de valorisation énergétique (UVE) traite \SI{50 000}{tonnes/an} de CSR papier (PCI = \SI{16,5}{MJ/kg}). La chaudière produit \SI{180 000}{tonnes/an} de vapeur surchauffée (\SI{400}{\celsius}, \SI{40}{bar}) à partir d'eau d'alimentation à \SI{105}{\celsius}. La turbine génère \SI{45 000}{MWh/an} d'électricité brute. L'autoconsommation de l'usine est de \SI{8 000}{MWh/an}. Le condenseur fournit \SI{120 000}{MWh/an} de chaleur vendue au réseau de chaleur.
\begin{enumerate}
    \item Calculer l'énergie thermique brute apportée par le combustible ($E_{th,in}$) en MWh/an.
    \item Déterminer le rendement chaudière $\eta_{chaud}$ (énergie vapeur / énergie combustible). Donnée : enthalpie vapeur $h_v = \SI{3214}{kJ/kg}$, enthalpie eau $h_e = \SI{440}{kJ/kg}$.
    \item Calculer le rendement électrique brut $\eta_{el,brut}$ et net $\eta_{el,net}$.
    \item Calculer le rendement global de cogénération $\eta_{glob}$.
    \item Conclure sur la performance de l'installation par rapport aux seuils réglementaires (Directive 2008/98/CE : $\eta_{el,net} > 16,2\%$ pour installation $> \SI{100}{MW_{th}}$ ou formule R1).
\end{enumerate}
\tcblower
\textbf{Solution détaillée :}
\begin{enumerate}
    \item \textbf{Énergie thermique entrante :}
    $m_{CSR} = 50\,000\,\text{t} = 5,0 \times 10^7\,\text{kg}$.
    $PCI = 16,5\,\text{MJ/kg} = \frac{16,5}{3,6}\,\text{kWh/kg} \approx 4,583\,\text{kWh/kg}$.
    $E_{th,in} = m \times PCI = 5,0 \times 10^7 \times 4,583 = 2,2915 \times 10^8\,\text{kWh} = \boxed{229\,150\,\text{MWh/an}}$.
    
    \item \textbf{Rendement chaudière :}
    Énergie vapeur utile : $Q_{vapeur} = m_{vapeur} \times (h_v - h_e)$.
    $m_{vapeur} = 180\,000\,\text{t} = 1,8 \times 10^8\,\text{kg}$.
    $\Delta h = 3214 - 440 = 2774\,\text{kJ/kg} = \frac{2774}{3600}\,\text{kWh/kg} \approx 0,7706\,\text{kWh/kg}$.
    $Q_{vapeur} = 1,8 \times 10^8 \times 0,7706 = 1,387 \times 10^8\,\text{kWh} = 138\,700\,\text{MWh/an}$.
    $\eta_{chaud} = \frac{Q_{vapeur}}{E_{th,in}} = \frac{138\,700}{229\,150} \approx 0,605 = \boxed{60,5\,\%}$.
    \textit{(Rendement réaliste pour chaudière à déchets, pertes fumées, cendres, rayonnement).}
    
    \item \textbf{Rendements électriques :}
    $\eta_{el,brut} = \frac{W_{el,brut}}{E_{th,in}} = \frac{45\,000}{229\,150} \approx 0,196 = \boxed{19,6\,\%}$.
    $W_{el,net} = 45\,000 - 8\,000 = 37\,000\,\text{MWh/an}$.
    $\eta_{el,net} = \frac{37\,000}{229\,150} \approx 0,1615 = \boxed{16,15\,\%}$.
    
    \item \textbf{Rendement global cogénération :}
    $\eta_{glob} = \frac{W_{el,net} + Q_{chaleur\_vendue}}{E_{th,in}} = \frac{37\,000 + 120\,000}{229\,150} = \frac{157\,000}{229\,150} \approx 0,685 = \boxed{68,5\,\%}$.
    
    \item \textbf{Conclusion :}
    Le rendement électrique net (\SI{16,15}{\%}) est \textbf{légèrement inférieur} au seuil de \SI{16,2}{\%} exigé par la formule R1 de la Directive européenne pour les installations de grande capacité afin d'être classé en "Valorisation Énergétique" (R1) et non en "Élimination" (D10). L'installation est performante thermiquement (\SI{68,5}{\%} global), mais l'optimisation du cycle vapeur (pression/température plus élevées, meilleure isolation) ou la réduction de l'autoconsommation permettrait de franchir le seuil réglementaire.
\end{enumerate}
\end{exoresolu}

\courssub{Recyclage du papier : serviettes jetables, papier hygiénique, papier-cadeau}

\begin{methode}[Recyclage du papier en produits d'hygiène et d'emballage (Boucle fibreuse)]
\textbf{Objectif :} Expliquer les étapes du recyclage matière du papier/carton en pâte désencrée pour fabriquer serviettes, papier hygiénique, papier-cadeau.
\begin{enumerate}
    \item \textbf{Collecte et tri (amont) :} Séparation à la source (bac bleu/jaune). Tri industriel : élimination des corps étrangers (métaux, plastiques, ficelles) par criblage, aimantation, séparation aéraulique. Tri par qualité (papiers graphiques, cartons, journaux).
    \item \textbf{Hydropulpage (Défibrage) :} Mélange papiers + eau (\SI{4-5}{\%} siccité) dans un pulpeur (grosse cuve avec rotor). Cisaillement hydraulique $\rightarrow$ séparation des fibres. Les contaminants lourds (agrafe, sable) tombent au fond (dépulpage "propre" ou "sale" selon tri amont).
    \item \textbf{Épuration grossière et fine :} 
    \begin{itemize}
        \item \textit{Dépurage centrifuge (cleaners) :} Cyclones (D $\sim$ \SI{150-250}{mm}) éliminent les particules denses (sable, verre, agrafes) par force centrifuge.
        \item \textit{Criblage (fente \SI{0,15-0,30}{mm}) :} Élimine les "stickies" (colles, plastiques souples) et gros morceaux de fibres non défibrées.
    \end{itemize}
    \item \textbf{Désencrage (Flottation) :} \textbf{Étape clé pour blancheur.} Ajout de réactifs : savon/collecteur (hydrophobe), chaux (pH \num{\sim 9-10}), dispersant. Injection d'air (bulles \SI{50-100}{\mu m}) $\rightarrow$ particules d'encre hydrophobes se fixent sur bulles $\rightarrow$ remontée en mousse (écrémage). Rendement désencrage \SI{70-90}{\%}.
    \item \textbf{Lavage et épaississement :} Élimination des fines, charges (carbonate de calcium), sels solubles, résidus de savon par presse à vis ou bande pressante. Siccité \SI{\sim 30-40}{\%}.
    \item \textbf{Affinage (Raffinage) :} Traitement mécanique (cône-refineur ou disque) pour développer le potentiel de liaison des fibres (fibrillation, raccourcissement contrôlé). Mesure : degré \cle{Schopper-Riegler} (\cle{°SR}) cible selon produit final (ex. \SI{25-35}{\degree SR} pour tissu/hygiène).
    \item \textbf{Fabrication papier (Machine à papier) :} 
    \begin{itemize}
        \item \textit{Partie humide :} Tête de forme $\rightarrow$ toile formant $\rightarrow$ pressage (siccité \SI{45-50}{\%}).
        \item \textit{Partie sèche :} Sécheurs à cylindres (vapeur) $\rightarrow$ calandrage (lissage/brillance).
        \item \textit{Spécificités :} \textbf{Papier hygiénique/serviettes :} Crêpage (lame Yankee \SI{> 90}{\%} siccité) pour douceur/volume (micro-plis). \textbf{Papier-cadeau :} Enduisage (couche pigmentée/brillante) + calandrage supercalandré.
    \end{itemize}
    \item \textbf{Bobinage et conversion :} Bobines mères $\rightarrow$ découpe, gaufrage, pliage, emballage.
\end{enumerate}
\cle{Hydropulpage}, \cle{Flottation (désencrage)}, \cle{Affinage}, \cle{Crêpage}, \cle{°SR}, \cle{Stickies}
\end{methode}

\begin{popfigure}
\begin{tikzpicture}[node distance=0.8cm and 1.5cm, every node/.style={font=\small, align=center}, >=Stealth]
\node (collecte) [draw, rounded corners, fill=popblueL, align=center]{Collecte \& Tri\\source (bac bleu)};
\node (triind) [right of=collecte, draw, rounded corners, fill=popblueL, align=center]{Tri industriel\\criblage, aimantation};
\node (pulpeur) [right of=triind, draw, rounded corners, fill=poporangeL, align=center]{Hydropulpage\\défibrage 4-5\%};
\node (epur) [below of=pulpeur, draw, rounded corners, fill=poporangeL, align=center]{Épuration\\cleaners, criblage};
\node (desenc) [right of=epur, draw, rounded corners, fill=popgreenL, align=center]{Désencrage\\Flottation pH 9-10};
\node (lavage) [right of=desenc, draw, rounded corners, fill=popgreenL, align=center]{Lavage \& Épaississement\\30-40\%};
\node (affin) [right of=lavage, draw, rounded corners, fill=poppurpleL, align=center]{Affinage\\°SR cible};
\node (machine) [right of=affin, draw, rounded corners, fill=poppinkL, align=center]{Machine à papier\\Yankee (crêpage) / Calandrage};
\node (produits) [below of=machine, draw, rounded corners, fill=popgoldL, align=center]{Produits finis\\Serviettes, Hygiénique, Cadeau};
\draw[->, thick] (collecte) -- (triind);
\draw[->, thick] (triind) -- (pulpeur);
\draw[->, thick] (pulpeur) -- (epur);
\draw[->, thick] (epur) -- (desenc);
\draw[->, thick] (desenc) -- (lavage);
\draw[->, thick] (lavage) -- (affin);
\draw[->, thick] (affin) -- (machine);
\draw[->, thick] (machine) -- (produits);
% Reject loops
\node[below of=triind, yshift=-0.5cm] (rej1) {Rejets CSR/Enfouissement};
\draw[->, dashed, popmuted] (triind) -- (rej1);
\node[below of=epur, yshift=-0.5cm] (rej2) {Rejets (sables, stickies)};
\draw[->, dashed, popmuted] (epur) -- (rej2);
\node[below of=desenc, yshift=-0.5cm] (rej3) {Boues d'encre};
\draw[->, dashed, popmuted] (desenc) -- (rej3);
\end{tikzpicture}
\legende{Schéma simplifié de la boucle de recyclage du papier (fibres graphiques $\rightarrow$ papiers d'hygiène/emballage). Les rejets partent en valorisation énergétique (CSR) ou élimination.}
\end{popfigure}

\begin{exoresolu}[Analyse d'une chaîne de recyclage papier - Identification des étapes critiques]
\textbf{Énoncé :} Une papeterie de recyclage produit du papier hygiénique à partir de magazines récupérés (Old Magazines - OM). Le process comprend : Pulpeur $\rightarrow$ Dépurage $\rightarrow$ Criblage \SI{0,25}{mm} $\rightarrow$ Flottation (2 stades) $\rightarrow$ Lavage $\rightarrow$ Affinage (Disque double) $\rightarrow$ Machine à papier (Yankee).
On observe sur le produit final : (1) Taches sombres résiduelles (maculage), (2) Faible résistance à la traction (papier "court"), (3) Présence de micro-billes plastiques brillantes.
\begin{enumerate}
    \item Identifier l'origine probable de chaque défaut et l'étape du process défaillante ou insuffisante.
    \item Proposer une modification du process (ajout/suppression/modification paramètre) pour corriger chaque défaut.
    \item Justifier pourquoi le crêpage sur Yankee est indispensable pour le papier hygiénique mais pas pour du papier journal recyclé.
\end{enumerate}
\tcblower
\textbf{Solution détaillée :}
\begin{enumerate}
    \item \textbf{Analyse des défauts :}
    \begin{itemize}
        \item \textbf{Défaut 1 : Taches sombres (maculage).} $\rightarrow$ \textit{Origine :} Particules d'encre résiduelles non éliminées. \textit{Étape défaillante :} \textbf{Flottation}. Soit chimie inadaptée (pH trop bas, savon insuffisant), soit temps de rétention trop court, soit 2ème stade inefficace.
        \item \textbf{Défaut 2 : Faible résistance à la traction (papier court).} $\rightarrow$ \textit{Origine :} Fibres trop courtes / sur-affinage / perte de fines lors lavage. \textit{Étape défaillante :} \textbf{Affinage} (énergie spécifique trop élevée, jeu disques trop serré) \textbf{OU} \textbf{Lavage} (perte excessive de fines colloïdales qui font la liaison fibre-fibre).
        \item \textbf{Défaut 3 : Micro-billes plastiques brillantes.} $\rightarrow$ \textit{Origine :} Contaminants "stickies" (colles thermofusibles, fragments de films plastiques des couvertures brillantes des magazines). \textit{Étape défaillante :} \textbf{Criblage} (fente \SI{0,25}{mm} trop large pour micro-billes < \SI{100}{\mu m}) \textbf{ET/OU} \textbf{Flottation} (stickies non hydrophobes ou trop petits pour bulles).
    \end{itemize}
    
    \item \textbf{Propositions correctives :}
    \begin{itemize}
        \item \textit{Maculage :} Optimiser chimie flottation (pH \num{9,5-10}, dose collecteur), augmenter temps rétention ou ajouter 3ème cellule. Contrôler dureté eau (chélatant).
        \item \textit{Résistance :} Réduire énergie spécifique affinage (cible °SR plus bas, ex \SI{28}{\degree SR} au lieu de \SI{35}{\degree SR}). Installer un \textit{save-all} (récupérateur de fines) sur eau blanche machine ou réduire taux de rejet lavage.
        \item \textit{Micro-billes :} Ajouter un \textbf{criblage fin fente \SI{0,10}{mm}} (ou \SI{0,15}{mm}) en amont flottation. Ajouter une étape \textbf{dispersion à chaud} (dispersion kneader \SI{95}{\celsius}, \SI{30}{\%} siccité) avant flottation pour éclater les stickies et les rendre flottables.
    \end{itemize}
    
    \item \textbf{Justification Crêpage (Yankee) :}
    Le papier hygiénique/serviettes requiert \textbf{volume, douceur, absorption rapide}. Le crêpage (décollage lame acier sur cylindre Yankee chauffé \SI{\sim 180}{\celsius}, siccité \SI{> 93}{\%}) crée des \textbf{micropuis} par flambage localisé des fibres. Cela augmente l'épaisseur (volume), la surface spécifique (capillarité $\rightarrow$ absorption) et la douceur (flexibilité). Le papier journal recyclé (standard ISO 287) cherche \textbf{opacité, rigidité, imprimabilité, faible coût} : pas de crêpage, calandrage léger pour lisser surface, grammage plus élevé (\SI{45-48}{g/m^2} vs \SI{15-20}{g/m^2}).
\end{enumerate}
\end{exoresolu}

\courssub{Recyclage des déchets plastiques en bouteilles recyclées (rPET alimentaire)}

\begin{methode}[Recyclage des déchets plastiques en bouteilles recyclées (Boucle fermée rPET)]
\textbf{Objectif :} Expliquer la technique de recyclage mécanique des bouteilles PET (Polyéthylène Téréphtalate) en nouvelles bouteilles (rPET alimentaire).
\begin{enumerate}
    \item \textbf{Collecte et Tri optique :} Bacs jaunes $\rightarrow$ Centre de tri. Séparation par \textbf{proche infrarouge (NIR)} : identification signature spectrale PET. Éjection par jets d'air. Pureté flux entrant > \SI{95}{\%}.
    \item \textbf{Broyage et Lavage à chaud :} Broyage en paillettes (\textit{flakes} \SI{\sim 10-12}{mm}). Lavage à \SI{85-95}{\celsius} avec soude (\ce{NaOH} \SI{1-2}{\%}) + tensioactif $\rightarrow$ décolle étiquettes (colle), élimine salissures, boissons résiduelles. Séparation par densité : PET (\SI{1,38}{g/cm^3}) coule, PP/PE bouchons (\SI{0,9}{g/cm^3}) flottent $\rightarrow$ séparation par flottation.
    \item \textbf{Séchage et Tri couleur :} Sécheur lit fluidisé. Tri optique couleur (caméras) : séparation clair (transparent), bleu, vert, couleur. Le flux "clair" est premium pour bouteille-bouteille.
    \item \textbf{Extrusion et Polycondensation en phase solide (SSP - Solid State Polycondensation) :} \textbf{Étape critique qualité alimentaire.}
    \begin{itemize}
        \item Extrusion double vis : fusion \SI{280}{\celsius}, filtration fine (\SI{30-60}{\mu m}), granulation (pellets). Élimination volatils (acétaldéhyde, éthylène glycol).
        \item SSP : Pellets chauffés \SI{180-220}{\celsius} sous vide/azote flux \SI{10-20}{h}. Diffusion des chaînes polymères $\rightarrow$ augmentation masse molaire (IV \textit{Intrinsic Viscosity} passe \SI{0,65}{dl/g} $\rightarrow$ \SI{0,80-0,85}{dl/g}). Élimination contaminants migrateurs (limite EFSA < \SI{0,1}{\mu g/kg} pour benzène etc.).
    \end{itemize}
    \item \textbf{Injection-Soufflage (ISBM) :} Préchauffage préformes \SI{110-120}{\celsius}. Soufflage biaxial (étirage axial + gonflage) $\rightarrow$ orientation molécules $\rightarrow$ transparence, résistance mécanique, barrière gaz.
    \item \textbf{Contrôle qualité et Traçabilité :} IV, couleur (L*a*b*), acétaldéhyde, métaux lourds, migration globale/spécifique (règlement UE 10/2011, 282/2008). Certificat EFSA pour process "super-clean".
\end{enumerate}
\cle{PET}, \cle{rPET}, \cle{SSP (Polycondensation phase solide)}, \cle{IV (Viscosité intrinsèque)}, \cle{Super-clean process}, \cle{EFSA}, \cle{ISBM}
\end{methode}

\begin{popfigure}
\begin{tikzpicture}[node distance=0.8cm and 1.2cm, every node/.style={font=\small, align=center}, >=Stealth]
\node (collecte) [draw, rounded corners, fill=popblueL, align=center]{Collecte sélective\\Bacs jaunes};
\node (tri) [right of=collecte, draw, rounded corners, fill=popblueL, align=center]{Tri optique NIR\\Pureté > 95\%};
\node (broy) [right of=tri, draw, rounded corners, fill=poporangeL, align=center]{Broyage\\Flakes 10-12 mm};
\node (lavage) [below of=broy, draw, rounded corners, fill=poporangeL, align=center]{Lavage chaud NaOH\\Séparation densité\\PET coule, PP/PE flotte};
\node (sech) [right of=lavage, draw, rounded corners, fill=popgreenL, align=center]{Séchage\\Tri couleur (Clair/Premium)};
\node (extr) [right of=sech, draw, rounded corners, fill=poppurpleL, align=center]{Extrusion double vis\\Filtration 30-60 µm\\Granulés};
\node (ssp) [right of=extr, draw, rounded corners, fill=poppurpleL, align=center]{SSP (Polycond. phase solide)\\180-220°C, Vide/N2\\IV: 0.65 $\rightarrow$ 0.82 dl/g};
\node (isb) [right of=ssp, draw, rounded corners, fill=poppinkL, align=center]{Injection-Soufflage (ISBM)\\Préformes $\rightarrow$ Bouteilles rPET};
\node (qual) [below of=isb, draw, rounded corners, fill=popgoldL, align=center]{Contrôle Qualité\\IV, Migration, EFSA};
\draw[->, thick] (collecte) -- (tri);
\draw[->, thick] (tri) -- (broy);
\draw[->, thick] (broy) -- (lavage);
\draw[->, thick] (lavage) -- (sech);
\draw[->, thick] (sech) -- (extr);
\draw[->, thick] (extr) -- (ssp);
\draw[->, thick] (ssp) -- (isb);
\draw[->, thick] (isb) -- (qual);
% Byproducts
\node[below of=lavage, yshift=-0.5cm] (bp1) {Bouchons PP/PE $\rightarrow$ Recyclage};
\draw[->, dashed, popmuted] (lavage) -- (bp1);
\node[below of=extr, yshift=-0.5cm, align=center] (bp2){Volatils (Acétaldéhyde)\\Épurés};
\draw[->, dashed, popmuted] (extr) -- (bp2);
\end{tikzpicture}
\legende{Processus « Super-clean » de recyclage bouteille-bouteille (rPET alimentaire). L'étape SSP (Solid State Polycondensation) est critique pour remonter la viscosité intrinsèque (IV) et éliminer les contaminants migrants (norme EFSA).}
\end{popfigure}

\begin{exoresolu}[Calcul de taux d'incorporation et bilan masse rPET]
\textbf{Énoncé :} Une usine de recyclage PET traite \SI{25 000}{tonnes/an} de balles de bouteilles triées (teneur PET \SI{92}{\%}, humidité \SI{8}{\%}, étiquettes/bouchons/autres \SI{8}{\%} matière sèche). Le process de lavage/flottation a un rendement masse (paillettes propres / entrée sèche PET) de \SI{88}{\%}. L'extrusion+SSP a un rendement masse (pellets rPET / paillettes) de \SI{96}{\%}. L'injection-soufflage a un rendement (bouteilles finies / pellets) de \SI{99}{\%}.
\begin{enumerate}
    \item Calculer la masse annuelle de pellets rPET produits (tonnes/an).
    \item Calculer le taux d'incorporation maximal de rPET dans une nouvelle bouteille de \SI{30}{g} si la préforme pèse \SI{32}{g} (perte soufflage) et que l'usine de préformes mélange rPET et PET vierge (IV vierge = \SI{0,84}{dl/g}, IV rPET = \SI{0,82}{dl/g}) pour viser IV finale \SI{0,83}{dl/g}. (Loi des mélanges linéaire sur IV).
    \item Déterminer le tonnage annuel de déchets ultimes (rebuts process) générés par l'usine de recyclage (hors bouchons/étiquettes séparés en amont).
\end{enumerate}
\tcblower
\textbf{Solution détaillée :}
\begin{enumerate}
    \item \textbf{Masse pellets rPET :}
    Masse entrée sèche = $25\,000 \times (1 - 0,08) = 23\,000\,\text{t/an}$.
    Masse PET pur entrée = $23\,000 \times 0,92 = 21\,160\,\text{t/an}$.
    Masse paillettes propres = $21\,160 \times 0,88 = 18\,620,8\,\text{t/an}$.
    Masse pellets rPET = $18\,620,8 \times 0,96 = \boxed{17\,876\,\text{t/an}}$ (arrondi).
    
    \item \textbf{Taux d'incorporation rPET (calcul sur IV) :}
    Soit $x$ la fraction massique de rPET dans le mélange préforme.
    $IV_{mélange} = x \cdot IV_{rPET} + (1-x) \cdot IV_{vierge}$.
    $0,83 = x \cdot 0,82 + (1-x) \cdot 0,84$.
    $0,83 = 0,82x + 0,84 - 0,84x$.
    $0,02x = 0,01 \Rightarrow x = 0,50$.
    Taux d'incorporation maximal = \boxed{50\,\% \text{ en masse}}.
    \textit{Remarque :} La perte au soufflage (\SI{2}{g}) n'affecte pas la composition du mélange, seulement le rendement matière final bouteille/préforme.
    
    \item \textbf{Déchets ultimes process (rebuts internes) :}
    Pertes lavage/flottation (PET non récupéré) = $21\,160 \times (1-0,88) = 2\,539,2\,\text{t}$.
    Pertes extrusion/SSP = $18\,620,8 \times (1-0,96) = 744,8\,\text{t}$.
    Pertes injection-soufflage = $17\,876 \times (1-0,99) = 178,8\,\text{t}$.
    Total rebuts process = $2\,539,2 + 744,8 + 178,8 = \boxed{3\,463\,\text{t/an}}$.
    \textit{Note :} Les \SI{2\,000}{t} (humidité + étiquettes/bouchons/autres) sont séparés en amont du process cœur (tri flottation), ils ne sont pas des "rebuts process" mais des "sous-produits tri" (CSR, recyclage PP/PE, boue station épuration).
\end{enumerate}
\end{exoresolu}

\courssub{Transformation des déchets plastiques mélangés en pavés (Upcycling local)}

\begin{methode}[Transformation des déchets plastiques mélangés en pavés (Upcycling / Valorisation matière)]
\textbf{Objectif :} Expliquer la technique de transformation des plastiques souillés, mélangés (PE, PP, PS, PET résidus) ou non recyclables en boucle fermée, en pavés/planches/plots (économie circulaire locale).
\begin{enumerate}
    \item \textbf{Collecte et Prétraitement "Low-cost" :} Pas de tri fin par polymère. Broyage grossier (\SI{20-50}{mm}). Lavage à l'eau froide/ambiante (simple rinçage) ou \textbf{pas de lavage} (tolérance terres/sables \SI{< 5-10}{\%}). Séchage solaire ou air chaud basse T°. \textit{Avantage :} faible CAPEX/OPEX, accepte plastiques sales (films agricoles, sacs, emballages multicouches).
    \item \textbf{Extrusion-Fusion / Moulage par compression (Batch) :} 
    \begin{itemize}
        \item \textit{Extrusion continue (vis unique/conique) :} Fusion \SI{180-220}{\celsius} (T° < dégradation PET/PVC). Vis à fort cisaillement $\rightarrow$ mélange homogène (compatibilisation physique). Filage $\rightarrow$ granulats "éco-granulats" $\rightarrow$ ré-injection/moulage.
        \item \textit{Moulage par compression direct (Batch - très courant au Cameroun) :} Mélange paillettes + additifs (charges minérales sable \SI{20-40}{\%}, colorants, compatibilisants type SEBS/g-MA \SI{1-3}{\%}) dans moule métallique chauffé (\SI{180-200}{\celsius}, \SI{50-100}{bar}) $\rightarrow$ maintien \SI{10-30}{min} $\rightarrow$ refroidissement sous presse $\rightarrow$ démoulage pavé.
    \end{itemize}
    \item \textbf{Gestion de l'hétérogénéité (Compatibilisation) :} Mélange PE/PP (non miscibles) $\rightarrow$ morphologie "île-mer" $\rightarrow$ fragilité interface. Ajout \textbf{compatibilisant} (copolymère bloc SEBS, polyoléfine greffée anhydride maléique) $\rightarrow$ réduit tension interfaciale, améliore adhérence, ductilité.
    \item \textbf{Contrôle qualité pavé :} Dimensions, masse volumique (\SI{1,1-1,4}{g/cm^3} selon charge), résistance compression (\SI{> 20}{MPa} pour voirie piétonne, \SI{> 40}{MPa} carrossable), absorption eau (\SI{< 1}{\%}), dilatation thermique (coeff élevé $\sim \SI{100-150 e-6}{K^{-1}}$ $\rightarrow$ joints de dilatation obligatoires).
    \item \textbf{Applications locales :} Pavés routes/rues (Yaoundé, Douala), dalles latérites stabilisées, plots électricité, bancs publics, latrines.
\end{enumerate}
\cle{Upcycling}, \cle{Compatibilisant}, \cle{Moulage compression}, \cle{Éco-granulats}, \cle{Économie circulaire locale}
\end{methode}

\begin{popfigure}
\begin{tikzpicture}[scale=0.8, every node/.style={font=\small}, >=Stealth]
% Mold cross section
\draw[thick, fill=popmuted!30] (0,0) rectangle (6,4); % Mold cavity
\draw[thick, fill=popdark!50] (-0.5,0) rectangle (0,4); % Left wall
\draw[thick, fill=popdark!50] (6,0) rectangle (6.5,4); % Right wall
\draw[thick, fill=popdark!50] (-0.5,0) rectangle (6.5,-0.5); % Bottom plate
\draw[thick, fill=popdark!50] (-0.5,4) rectangle (6.5,4.5); % Top platen
% Material
\draw[fill=poporange!50] (0.2,0.2) rectangle (5.8,3.8);
% Arrows
\draw[->, thick, popred] (3,4.8) -- (3,4.2) node[midway, right] {Pression 50-100 bar};
\draw[->, thick, popred] (-1,2) -- (-0.2,2) node[midway, above] {Chauffage 180-200°C};
\draw[->, thick, popred] (7,2) -- (6.2,2) node[midway, above] {Chauffage 180-200°C};
% Composition labels
\node[align=left] at (8, 3.5) {Formulation massique :};
\node[align=left] at (8, 3.0) {55\% Plastiques mélangés (PE/PP/PS/PET)};
\node[align=left] at (8, 2.5) {40\% Sable fin (charge minérale)};
\node[align=left] at (8, 2.0) {3\% Compatibilisant (SEBS-g-MA)};
\node[align=left] at (8, 1.5) {2\% Pigments/Additifs};
\node[align=left] at (8, 1.0) {Retrait volumique ~2,5\%};
\node[align=left] at (8, 0.5) {$\rho_{\text{pavé}} \approx 1,28\,\text{g/cm}^3$};
\node[align=left] at (8, 0.0) {$R_c \approx 32\,\text{MPa}$};
\end{tikzpicture}
\legende{Principes du moulage par compression (batch) pour pavés en plastiques mélangés. Le sable assure la rigidité et réduit le coût ; le compatibilisant améliore l'adhérence aux interfaces polymères immiscibles. Joints de dilatation obligatoires (fort coefficient de dilatation thermique).}
\end{popfigure}

\begin{exoresolu}[Formulation et dimensionnement d'un pavé en plastiques mélangés]
\textbf{Énoncé :} Une coopérative camerounaise produit des pavés \SI{20 x 10 x 6}{cm} par moulage compression. Formulation massique : \SI{55}{\%} plastiques mélangés (PE/PP/PS/PET broyés, $\rho \approx \SI{0,95}{g/cm^3}$), \SI{40}{\%} sable fin (\SI{0-2}{mm}, $\rho = \SI{2,65}{g/cm^3}$), \SI{3}{\%} compatibilisant (SEBS-g-MA, $\rho = \SI{0,91}{g/cm^3}$), \SI{2}{\%} pigment/autres. Le moule a un volume utile de \SI{1200}{cm^3}. Le taux de retrait volumique au refroidissement est de \SI{2,5}{\%}.
\begin{enumerate}
    \item Calculer la masse théorique d'un pavé brut (sortie moule) en kg.
    \item Calculer la masse de chaque composant pour produire \SI{1000}{pavés/jour}.
    \item Estimer la résistance à la compression approximative si la matrice polymère pure donne \SI{15}{MPa} et l'ajout de \SI{40}{\%} sable (charge rigide) multiplie la résistance par un facteur 2,5 (règle de mélange simplifiée), mais la présence de \SI{5}{\%} PET (fragile, non compatibilisé localement) réduit le gain de \SI{15}{\%}. Conclure sur l'usage voirie.
\end{enumerate}
\tcblower
\textbf{Solution détaillée :}
\begin{enumerate}
    \item \textbf{Masse pavé brut :}
    Volume moule $V_m = 20 \times 10 \times 6 = 1200\,\text{cm}^3$.
    Volume pavé final $V_p = V_m \times (1 - 0,025) = 1200 \times 0,975 = 1170\,\text{cm}^3$.
    Masse volumique mélange $\rho_{mél}$ (règle des mélanges sur masse/volume) :
    $\frac{1}{\rho_{mél}} = \sum \frac{w_i}{\rho_i} = \frac{0,55}{0,95} + \frac{0,40}{2,65} + \frac{0,03}{0,91} + \frac{0,02}{1,20\,(est.)} \approx 0,5789 + 0,1509 + 0,0330 + 0,0167 = 0,7795\,\text{cm}^3/\text{g}$.
    $\rho_{mél} \approx 1,283\,\text{g/cm}^3$.
    Masse pavé $m = \rho_{mél} \times V_p = 1,283 \times 1170 \approx 1501\,\text{g} = \boxed{1,50\,\text{kg}}$.
    
    \item \textbf{Masses quotidiennes (1000 pavés) :}
    Masse totale = $1000 \times 1,501 = 1501\,\text{kg}$.
    Plastiques : $0,55 \times 1501 = \boxed{825,6\,\text{kg}}$.
    Sable : $0,40 \times 1501 = \boxed{600,4\,\text{kg}}$.
    Compatibilisant : $0,03 \times 1501 = \boxed{45,0\,\text{kg}}$.
    Pigments : $0,02 \times 1501 = \boxed{30,0\,\text{kg}}$.
    
    \item \textbf{Résistance compression estimée :}
    $R_{matrice} = 15\,\text{MPa}$.
    Gain charge sable (40\%) : facteur 2,5 $\rightarrow R_1 = 15 \times 2,5 = 37,5\,\text{MPa}$.
    Pénalité PET 5\% (fragilité interface) : réduction 15\% $\rightarrow R_{final} = 37,5 \times (1 - 0,15) = 37,5 \times 0,85 = \boxed{31,9\,\text{MPa}}$.
    \textbf{Conclusion :} Résistance \SI{\sim 32}{MPa} > \SI{20}{MPa} (voirie piétonne/parking léger) mais < \SI{40}{MPa} (voirie lourde carrossable classe T4/T5). Adapté pour trottoirs, cours, places de marché. Nécessite joints de dilatation (dilatation thermique élevée).
\end{enumerate}
\end{exoresolu}

\courssub{Avantages de la valorisation et sensibilisation de la communauté}

\begin{methode}[Sensibilisation communautaire : Informer sur la nécessité de transformer et recycler]
\textbf{Objectif :} Structurer un message de sensibilisation efficace (oral, affiche, radio, réseaux sociaux) basé sur les arguments scientifiques, économiques et réglementaires.
\begin{enumerate}
    \item \textbf{Diagnostic local (Contexte) :} Identifier les déchets visibles (dépôts sauvages, caniveaux bouchés $\rightarrow$ inondations, brûlage à l'air libre $\rightarrow$ fumées toxiques). Chiffres clés locaux (ex. Douala \SI{\sim 2000}{t/j} déchets, \SI{< 10}{\%} recyclés).
    \item \textbf{Arguments Scientifiques (Savoirs) :}
    \begin{itemize}
        \item \textit{Pollution :} Lixiviats (métaux lourds, COD) $\rightarrow$ nappe phréatique / rivière (Wouri, Mfoundi). Microplastiques $\rightarrow$ chaîne alimentaire (poissons Wouri). Brûlage \ce{PVC} $\rightarrow$ dioxines (cancérigènes CIRC 1).
        \item \textit{Ressources :} Papier = bois/énergie/eau économisés (\SI{1}{t} papier recyclé = \SI{17}{arbres}, \SI{50}{\%} énergie, \SI{80}{\%} eau). Plastique = pétrole économisé (\SI{1}{t} PET recyclé = \SI{\sim 700}{kg} pétrole brut).
        \item \textit{Climat :} Évitement \ce{CH4} (décharge anaérobie, PRG 28) et \ce{CO2} (incinération sans récupération / production vierge).
    \end{itemize}
    \item \textbf{Arguments Économiques \& Sociaux :}
    \begin{itemize}
        \item Filières locales créatrices d'emplois (collecteurs, trieurs, broyeurs, artisans pavés, papeteries). Ex: \textit{Redplast, EcoCollect, Namé Recycling}.
        \item Réduction coûts gestion municipale (transport, enfouissement).
        \item Revenus pour ménages (rachat kg plastique/papier).
    \end{itemize}
    \item \textbf{Appel à l'action concret (Gestes quotidiens) :}
    \begin{enumerate}
        \item Tri à la source : 2 bacs (Organique / Sec recyclable).
        \item Dépôt points de collecte identifiés (écoles, marchés, mairies, ONG).
        \item Refus sacs plastiques interdits (\SI{< 60}{\mu m}), adoption sacs réutilisables.
        \item Signalement dépôts sauvages (numéro vert / appli).
    \end{enumerate}
    \item \textbf{Format du message :} Court, visuel (schéma cycle), langue locale + français, ton positif (gain) pas seulement punitif. "Mon déchet = Ma ressource".
\end{enumerate}
\cle{Tri à la source}, \cle{Lixiviat}, \cle{Microplastiques}, \cle{Économie circulaire}, \cle{Responsabilité Élargie du Producteur (REP)}
\end{methode}

\begin{exoresolu}[Rédaction d'un message de sensibilisation radio (2 minutes)]
\textbf{Énoncé :} Vous êtes chargé de concevoir un spot radio de \SI{120}{secondes} (environ 250 mots) en français et en langue locale (ex. Duala, Bassa, Ewondo, Fulfulde) pour le "Mois de l'Environnement" à Yaoundé. Thème : "Trions nos déchets, valorisons notre quartier". Le spot doit intégrer : 1 fait chiffré local, 1 explication scientifique simple (cycle), 1 avantage économique direct pour l'auditeur, 1 geste précis, 1 contact utile.
\tcblower
\textbf{Proposition de script (Version Française - \SI{\sim 130}{mots} = \SI{\sim 60}{s}, laisser place traduction + musique) :}

\textbf{[Jingle court : Bruit de bouteilles qui s'entrechoquent, puis musique douce type Makossa/Acoustique]}

\textbf{Voix Homme (Calme) :} « Chers habitants de Yaoundé, saviez-vous que notre ville produit \textbf{1 500 tonnes de déchets par jour} ? Et que \textbf{la moitié} pourrait avoir une seconde vie ? »

\textbf{Voix Femme (Dynamique) :} « Vos vieux journaux ? Ils deviennent du \textbf{papier hygiénique} ou des \textbf{serviettes} ici même, à la papeterie de la SOPECAM. Vos bouteilles en plastique ? Elles sont transformées en \textbf{nouvelles bouteilles} ou en \textbf{pavés solides} pour nos routes ! »

\textbf{Voix Homme :} « Brûler ou jeter dans le caniveau, c'est polluer l'eau du robinet et l'air qu'on respire. Trier, c'est \textbf{protéger la santé de nos enfants}. »

\textbf{Voix Femme :} « Et c'est \textbf{gagner de l'argent} ! Apportez vos plastiques et papiers aux points \textbf{EcoCollect} ou chez les \textbf{collecteurs agréés} : on vous les rachète au kilo. »

\textbf{Ensemble (Tonique) :} « \textbf{Un bac pour l'organique, un sac pour le sec recyclable.} Simple, non ? »

\textbf{Voix Homme :} « Points de collecte et infos : \textbf{15 00 00 00} (numéro fictif) ou page Facebook \textit{Yaoundé Propre}. »

\textbf{[Jingle final]} « \textit{Trions, Valorisons, Vivons mieux !} »

\textbf{Analyse pédagogique :} Respecte la contrainte temps. Intègre : Chiffre local (1500 t/j), Cycle (journal $\rightarrow$ hygiène, bouteille $\rightarrow$ pavé/bouteille), Santé (lixiviat/air), Économie (rachat kg), Geste clair (2 bacs), Contact. Adaptable en langue locale (vocabulaire courant : "bouteille", "papier", "argent", "santé", "enfant").
\end{exoresolu}

\begin{methode}[Analyse comparative des filières de valorisation : Avantages Environnementaux \& Économiques]
\textbf{Objectif :} Comparer les scénarios de fin de vie (Enfouissement, Incinération simple, Valorisation Énergétique, Recyclage Matière, Upcycling) pour argumenter le choix de la valorisation.
\begin{enumerate}
    \item \textbf{Indicateurs Environnementaux (ACV simplifiée - Analyse Cycle de Vie) :}
    \begin{itemize}
        \item \textit{GES (kg \ce{CO2} eq/t déchet) :} Enfouissement (\SI{+400}{à +800} \ce{CH4} fuite) > Incinération simple (\SI{+300}{à +500}) > Valorisation Énergétique (\SI{-200}{à +100} crédit énergie) > Recyclage Papier (\SI{-500}{à -900} crédit bois/énergie) > Recyclage Plastique (\SI{-1000}{à -2000} crédit pétrole/énergie).
        \item \textit{Épuisement ressources abiotiques :} Recyclage $\ll$ Enfouissement/Incinération.
        \item \textit{Qualité air/eau/sols :} Valorisation maîtrisée (épuration) < Enfouissement (lixiviat) < Brûlage sauvage.
    \end{itemize}
    \item \textbf{Indicateurs Économiques :}
    \begin{itemize}
        \item \textit{Coût global (Collecte + Traitement - Recettes matière/énergie) :} Enfouissement (\SI{30-50}{kFCFA/t}), Valorisation Énergétique (\SI{40-70}{kFCFA/t} mais recettes énergie), Recyclage Industriel (Coût net souvent négatif = \textbf{bénéfice} si tri amont bon), Upcycling local (Faible CAPEX, main d'œuvre locale, valeur ajoutée locale).
        \item \textit{Emplois (ETP / 10 000 t/an) :} Enfouissement (1-2), Incinération (5-10), Tri/Recyclage (30-100), Upcycling artisanal (50-200).
    \end{itemize}
    \item \textbf{Hiérarchie des déchets (Directive Cadre / Loi Cadre Cameroun 96/12) :} Prévention > Réemploi > Recyclage (Matière) > Valorisation (Énergie) > Élimination. \textbf{Le recyclage matière (papier, PET bouteilles) est prioritaire sur la valorisation énergétique.}
    \item \textbf{Argumentaire de décision :} "Pour \SI{1}{tonne} de papier mélangé : Recycler en hygiène économise \SI{17}{arbres} et \SI{50}{\%} énergie vs vierge. Brûler en CSR récupère \SI{16}{MJ/kg} mais perd la fibre à jamais. Choix optimal : Tri $\rightarrow$ Recyclage matière (fibres longues) $\rightarrow$ Valorisation énergie (fibres courtes/souillées)."
\end{enumerate}
\cle{ACV}, \cle{Hiérarchie des déchets}, \cle{Crédit carbone}, \cle{REP (Responsabilité Élargie du Producteur)}, \cle{Économie circulaire}
\end{methode}

\begin{exoresolu}[Choix de filière pour un déchet papier/plastique mixte - Étude de cas communal]
\textbf{Énoncé :} La commune de Dschang (\SI{\sim 200 000}{hab}) produit \SI{50}{t/j} de déchets ménagers (OM \SI{60}{\%}, Papier/Carton \SI{15}{\%}, Plastique \SI{12}{\%}, Verre/Métaux/Autres \SI{13}{\%}). Trois scénarios sont proposés pour la fraction sèche (Papier + Plastique = \SI{13,5}{t/j}) :
\begin{itemize}
    \item \textbf{Scén. A :} Enfouissement tout-venant (coût \SI{25 000}{FCFA/t}).
    \item \textbf{Scén. B :} Tri mécanique centralisé + Valorisation Énergétique CSR (Papier/Plastique) dans cimenterie locale (coût tri+transport \SI{35 000}{FCFA/t}, revenu énergie \SI{5 000}{FCFA/t}).
    \item \textbf{Scén. C :} Tri à la source (2 bacs) + Collecte sélective + Recyclage matière Papier (Papeterie \SI{50}{km}, revenu \SI{40 000}{FCFA/t} papier trié) + Recyclage matière Plastique (Unité locale pavés, coût net \SI{10 000}{FCFA/t} plastique trié).
\end{itemize}
Hypothèses : Taux de capture tri source = \SI{60}{\%}. Taux tri mécanique = \SI{80}{\%}. Pureté papier tri source = \SI{90}{\%}, tri méca = \SI{70}{\%}. Pureté plastique tri source = \SI{85}{\%}, tri méca = \SI{60}{\%}.
\begin{enumerate}
    \item Calculer le coût net annuel (FCFA/an) pour chaque scénario sur la fraction sèche.
    \item Calculer le tonnage annuel de papier et plastique \textbf{effectivement recyclés/valorisés} (sortie process) pour B et C.
    \item Recommander le scénario optimal en intégrant critères économique, environnemental (hiérarchie), social (emplois).
\end{enumerate}
\tcblower
\textbf{Solution détaillée :}
\textbf{Données de base :} $M_{sec} = 13,5\,\text{t/j} \times 365 = 4\,927,5\,\text{t/an}$. Papier entrant = $0,15/0,27 \times 13,5 = 7,5\,\text{t/j}$. Plastique entrant = $6,0\,\text{t/j}$.

\begin{enumerate}
    \item \textbf{Coûts nets annuels :}
    \textbf{Scén. A (Enfouissement) :} $4\,927,5 \times 25\,000 = \boxed{123\,187\,500\,\text{FCFA/an}}$.
    
    \textbf{Scén. B (Tri méca + VE Cimenterie) :}
    Tonnage traité = $4\,927,5 \times 0,80 = 3\,942\,\text{t/an}$.
    Coût = $3\,942 \times 35\,000 = 137\,970\,000$.
    Revenu = $3\,942 \times 5\,000 = 19\,710\,000$.
    Coût net = $137,97 - 19,71 = \boxed{118\,260\,000\,\text{FCFA/an}}$.
    \textit{(Légèrement moins cher qu'enfouissement, mais perte matière).}
    
    \textbf{Scén. C (Tri source + Recyclage Matière) :}
    \textit{Papier :} Capté = $7,5 \times 365 \times 0,60 = 1\,642,5\,\text{t/an}$. Trié (90\%) = $1\,478\,\text{t/an}$. Revenu = $1\,478 \times 40\,000 = 59\,120\,000$.
    \textit{Plastique :} Capté = $6,0 \times 365 \times 0,60 = 1\,314\,\text{t/an}$. Trié (85\%) = $1\,117\,\text{t/an}$. Coût net = $1\,117 \times 10\,000 = 11\,170\,000$ (charge).
    \textit{Restes (non captés + refus tri) :} $4\,927,5 - (1\,642,5+1\,314) = 1\,971\,\text{t/an}$ enfouis ($1\,971 \times 25\,000 = 49\,275\,000$).
    \textit{Coût collecte sélective (supplément) :} Estimation \SI{15 000}{FCFA/t} capté $\rightarrow (1\,642,5+1\,314) \times 15\,000 = 44\,347\,500$.
    \textbf{Total Scén C :} $49,275 + 44,347 + 11,170 - 59,120 = \boxed{45\,672\,500\,\text{FCFA/an}}$.
    
    \item \textbf{Tonnages valorisés (Sortie process) :}
    \textbf{B :} $3\,942\,\text{t/an}$ valorisées en énergie (perte matière totale).
    \textbf{C :} Papier recyclé : $1\,478\,\text{t/an}$. Plastique recyclé (pavés) : $1\,117\,\text{t/an}$. Total matière = \boxed{2\,595\,\text{t/an}}.
    
    \item \textbf{Recommandation :} \textbf{Scén. C (Tri à la source + Recyclage Matière)}.
    \begin{itemize}
        \item \textit{Économique :} Coût net \textbf{2,6 fois inférieur} au Scén. A/B (\SI{45,7}{M} vs \SI{118-123}{M} FCFA). Revenus matière réels.
        \item \textit{Environnemental :} Respecte hiérarchie (Recyclage > Valorisation Énergie). Évite \ce{CO2} production vierge. Préserve ressources.
        \item \textit{Social :} Crée emplois collecte sélective, tri manuel, unité pavés locale (forte intensité main d'œuvre). Éducation citoyenne (geste tri).
    \end{itemize}
    \textit{Condition de réussite :} Campagne sensibilisation soutenue (Méthode 5) pour atteindre 60\% capture. Investissement initial bacs/collecte.
\end{enumerate}
\end{exoresolu}

\begin{attention}[Les 5 erreurs qui coûtent des points au Bac]
\begin{enumerate}
    \item \textbf{Confondre "Recyclage" et "Valorisation Énergétique" :} Écrire "le papier est recyclé en énergie" est une \textbf{erreur conceptuelle majeure}. Le recyclage matière conserve la matière (fibre, polymère). La valorisation énergétique détruit la matière pour en récupérer l'énergie. \textit{Au Bac : "Valorisation" est le terme générique ; il faut préciser "Recyclage matière" ou "Valorisation énergétique".}
    \item \textbf{Oublier l'étape de "Compatibilisation" pour les pavés en plastiques mélangés :} Dire "on fond tout ensemble et ça fait un pavé solide" ignore l'immiscibilité PE/PP/PET. Sans compatibilisant (ou charge minérale forte), le pavé est friable (mauvaise adhérence interfaces). \textit{Point clé : citer le rôle du compatibilisant (SEBS, g-MA) ou de la charge (sable) comme "pont physique".}
    \item \textbf{Ne pas distinguer les filières PET "Bouteille-Bouteille" (rPET alimentaire) et "Bouteille-Textile/Pavé" :} Le process "Super-clean" (SSP, IV cible \SI{>0,80}{dl/g}, EFSA) est \textbf{exigible} pour le contact alimentaire. Le simple broyage-lavage-extrusion (IV \SI{\sim 0,65}{dl/g}) ne produit que du rPET non alimentaire (fibres, straps, pavés). \textit{Au Bac : Préciser "Process Super-clean / SSP / Agrément EFSA" pour bouteille-bouteille.}
    \item \textbf{Calculs de rendement : Oublier l'humidité ou la siccité :} Les tonnages "balles" ou "entrées usine" incluent l'eau (\SI{8-12}{\%}). Les rendements process s'appliquent sur \textbf{masse sèche}. Ne pas corriger l'humidité fausse tous les bilans de masse. \textit{Réflexe : } $M_{sec} = M_{brut} \times (1 - H)$.
    \item \textbf{Sensibilisation : Message trop technique ou culpabilisant sans solution concrète :} "Vous polluez, triez !" sans dire \textbf{OÙ, COMMENT, QUOI GAGNER}. Perte de points en "Savoir-être / Communication". \textit{Exiger : Geste clair (2 bacs), Point de collecte précis, Avantage direct (argent, santé, cadre de vie), Contact.}
\end{enumerate}
\end{attention}

\newpage

\coursec{Exercices}

\courssub{Je vérifie mes connaissances}

\begin{exercice}[QCM : choisis la (ou les) bonne(s) réponse(s)]
Pour chaque question, une seule réponse est exacte. Recopie la lettre correspondante.

\textbf{1.} La valorisation d'un déchet consiste à :
\begin{itemize}
\item[a)] l'enfouir dans une décharge contrôlée ;
\item[b)] lui donner une nouvelle utilité (matière ou énergie) ;
\item[c)] le brûler à l'air libre ;
\item[d)] l'exporter vers un autre pays.
\end{itemize}

\textbf{2.} Dans la fabrication des pavés à partir de déchets plastiques, le plastique fondu joue le rôle de :
\begin{itemize}
\item[a)] granulat ; \quad b)] liant ; \quad c)] colorant ; \quad d)] combustible.
\end{itemize}

\textbf{3.} L'opération qui consiste à éliminer les encres d'impression au cours du recyclage du papier s'appelle :
\begin{itemize}
\item[a)] la mise en pâte ; \quad b)] le pressage ; \quad c)] le désencrage ; \quad d)] le calandrage.
\end{itemize}

\textbf{4.} Les briquettes de papier compressé sont un produit de :
\begin{itemize}
\item[a)] la valorisation énergétique ; \quad b)] la valorisation matière ; \quad c)] l'enfouissement ; \quad d)] le compostage.
\end{itemize}
\end{exercice}

\begin{exercice}[Vrai ou faux ? Justifie]
Réponds par \textbf{vrai} ou \textbf{faux}, puis justifie ta réponse en une ou deux phrases.

\textbf{1.} Le recyclage du papier permet de réduire le nombre d'arbres abattus pour la fabrication du papier neuf.

\textbf{2.} Tous les types de plastiques peuvent être mélangés sans tri préalable lors du recyclage en bouteilles.

\textbf{3.} La production d'énergie à partir du papier est une transformation chimique du déchet.

\textbf{4.} La fabrication de pavés à base de plastique et de sable ne présente aucun intérêt économique pour les jeunes entrepreneurs.
\end{exercice}

\begin{exercice}[Définitions]
Définis en deux ou trois lignes chacun des termes suivants, en donnant un exemple camerounais pour chacun :
\begin{itemize}
\item déchet ;
\item transformation ;
\item recyclage ;
\item valorisation.
\end{itemize}
\end{exercice}

\begin{exercice}[Calcul direct : briquettes de papier]
Une coopérative de la ville de Bafoussam collecte \SI{500}{kg} de papiers usagés par mois. Après mise en pâte, essorage et séchage, le rendement en briquettes combustibles est de \SI{60}{\%} en masse. Chaque briquette a une masse moyenne de \SI{250}{g}.

\textbf{1.} Calcule la masse totale de briquettes obtenue en un mois.

\textbf{2.} Calcule le nombre de briquettes produites en un mois.

\textbf{3.} Sachant que le pouvoir calorifique d'une briquette est de \SI{15}{MJ/kg}, calcule l'énergie totale (en \si{MJ}) que peut fournir la production mensuelle.
\end{exercice}

\courssub{Je m'entraîne}

\begin{exercice}[Recyclage ou transformation ?]
Voici une liste d'opérations réalisées dans des unités de valorisation des déchets :
\begin{itemize}
\item[a)] fabrication de briquettes combustibles à partir de vieux papiers ;
\item[b)] fabrication de pavés à partir de sachets plastiques fondus et de sable ;
\item[c)] fabrication de nouvelles bouteilles à partir de bouteilles plastiques usagées ;
\item[d)] fabrication de papier-cadeau à partir de papier journal usagé ;
\item[e)] fabrication de papier hygiénique à partir de cartons usagés.
\end{itemize}

\textbf{1.} Rappelle la différence entre recyclage et transformation.

\textbf{2.} Classe chacune de ces opérations dans la catégorie « recyclage » ou « transformation », en justifiant ton choix pour chacune.
\end{exercice}

\begin{exercice}[La chaîne de recyclage du papier]
On donne ci-dessous, dans le désordre, les étapes du recyclage du papier en papier hygiénique :
\begin{center}
pressage -- collecte et tri -- séchage -- mise en pâte -- désencrage -- blanchiment -- bobinage et découpe
\end{center}

\textbf{1.} Ordonne ces étapes de façon chronologique.

\textbf{2.} Définis chacune des opérations suivantes : mise en pâte, désencrage, pressage, séchage.

\textbf{3.} Explique pourquoi l'étape de blanchiment est particulièrement importante pour la fabrication du papier hygiénique.
\end{exercice}

\begin{exercice}[Étude de données : papier recyclé contre bois de chauffe]
Le tableau ci-dessous compare le bois de chauffe et les briquettes de papier utilisées par des ménages de Dschang.

\begin{center}
\begin{tabularx}{\linewidth}{|l|X|X|}
\hline
\rowcolor{popblueL} \textbf{Caractéristique} & \textbf{Bois de chauffe} & \textbf{Briquette de papier} \\
\hline
Pouvoir calorifique (\si{MJ/kg}) & 16 & 15 \\
\hline
Coût moyen (FCFA/kg) & 150 & 80 \\
\hline
Production de fumées & importante & modérée \\
\hline
Impact sur les forêts & déforestation & aucun \\
\hline
\end{tabularx}
\end{center}

\textbf{1.} Calcule l'énergie obtenue avec \SI{10}{kg} de chaque combustible.

\textbf{2.} Calcule le coût de \SI{10}{kg} de chaque combustible, puis le coût par mégajoule (FCFA/MJ) dans les deux cas.

\textbf{3.} À partir de ces résultats et des autres données du tableau, dégage deux avantages des briquettes de papier par rapport au bois de chauffe.
\end{exercice}

\begin{exercice}[Schéma de la fabrication des pavés]
Une association de jeunes de Douala fabrique des pavés en mélangeant des sachets plastiques fondus et du sable.

\textbf{1.} Réalise un organigramme simple (étapes encadrées reliées par des flèches) présentant les quatre grandes étapes de cette fabrication, de la collecte des sachets au pavé fini.

\textbf{2.} Indique sur ton organigramme l'étape qui nécessite le port d'équipements de protection, et précise lesquels.

\textbf{3.} Explique pourquoi le sable doit être sec et tamisé avant d'être mélangé au plastique fondu.
\end{exercice}

\courssub{Je me prépare au Bac}

\begin{exercice}[Type Bac : l'unité de recyclage du papier de Yaoundé]
\textbf{Document 1.} Une PME installée dans le quartier de Nkolbisson (Yaoundé) récupère chaque semaine \SI{2}{tonnes} de papiers et cartons usagés auprès des écoles, administrations et ménages. Le processus de production comprend les étapes suivantes, présentées ici dans le désordre :

\begin{center}
\begin{tabularx}{\linewidth}{|c|X|}
\hline
\rowcolor{popblueL} \textbf{Étape} & \textbf{Description} \\
\hline
A & La feuille humide est débarrassée de l'eau par compression entre des rouleaux. \\
\hline
B & Les papiers sont triés, puis broyés et mélangés à l'eau chaude pour former une suspension de fibres. \\
\hline
C & Les encres d'impression sont séparées des fibres par action de produits chimiques et par flottation. \\
\hline
D & La feuille est chauffée puis enroulée en bobines, avant d'être découpée selon le produit voulu. \\
\hline
E & La pâte est étalée en fine couche sur une toile pour former une feuille continue. \\
\hline
\end{tabularx}
\end{center}

\textbf{Document 2.} L'unité fabrique trois produits : des serviettes jetables, du papier hygiénique et du papier-cadeau. Le rendement global de la fabrication est de \SI{70}{\%} en masse : \SI{1}{kg} de papier collecté fournit \SI{0,70}{kg} de papier recyclé.

\textbf{1.} À l'aide du document 1, ordonne les étapes A à E de façon chronologique et nomme chaque opération (mise en pâte, désencrage, formation de la feuille, pressage, séchage-bobinage).

\textbf{2.} Explique le principe de l'étape C et précise pourquoi elle est indispensable pour obtenir un papier propre.

\textbf{3.} Indique, pour chacun des trois produits du document 2, une caractéristique de fabrication particulière (épaisseur, traitement, coloration).

\textbf{4.} À l'aide du document 2, calcule la masse de papier recyclé produite chaque semaine par l'unité, puis la masse produite en une année de 48 semaines de travail.

\textbf{5.} Sachant qu'il faut abattre en moyenne un arbre pour produire \SI{60}{kg} de papier neuf, estime le nombre d'arbres « économisés » par an grâce à cette unité. Conclus sur l'intérêt environnemental de cette activité.
\end{exercice}

\begin{exercice}[Type Bac : les pavés écologiques de Buea]
\textbf{Document 1.} Une entreprise de la région du Sud-Ouest transforme les sachets plastiques usagés en pavés de voirie. Le procédé est le suivant : les sachets sont collectés, triés, lavés et séchés ; ils sont ensuite fondus dans une cuve chauffée à environ \SI{180}{\celsius} ; le sable sec et tamisé est incorporé au plastique fondu dans un rapport de \SI{70}{\%} de sable pour \SI{30}{\%} de plastique (en masse) ; le mélange pâteux est versé dans des moules, compacté, puis laissé à refroidir. Les pavés obtenus sont plus résistants et moins chers que les pavés en ciment classiques.

\textbf{Document 2.} La fusion des plastiques dégage des vapeurs irritantes ; de plus, la cuve de fusion et le mélange chaud présentent des risques de brûlures.

\textbf{1.} Dégage le principe général de la technique de fabrication des pavés plastiques.

\textbf{2.} Précise le rôle joué par le plastique fondu et celui joué par le sable dans le matériau obtenu. À quel matériau de construction classique peut-on comparer ce pavé ? Justifie l'analogie.

\textbf{3.} À l'aide du document 2, cite deux précautions que doivent prendre les ouvriers lors de la fabrication, en précisant le risque évité dans chaque cas.

\textbf{4.} Calcule la masse de sachets plastiques et la masse de sable nécessaires à la fabrication d'un lot de pavés de masse totale \SI{500}{kg}.

\textbf{5.} Dégage deux avantages de cette technique pour l'environnement et deux avantages pour l'économie locale.
\end{exercice}

\begin{exercice}[Type Bac : recyclage des bouteilles plastiques et sensibilisation]
\textbf{Document 1.} Voici l'organigramme incomplet de la chaîne de recyclage des bouteilles plastiques dans une usine de Douala :

\begin{center}
\begin{tabularx}{\linewidth}{|X|X|X|X|X|}
\hline
\rowcolor{popblueL} Collecte & Étape 1 & Lavage & Étape 2 & Fusion et moulage \\
\hline
\end{tabularx}
\end{center}

\textbf{Document 2.} Les bouteilles plastiques sont fabriquées à partir de différents types de polymères (PET, PEHD, PVC...). Chaque type de plastique possède une température de fusion et des propriétés propres. Un mélange de plastiques incompatibles donne un matériau recyclé fragile et inutilisable.

\textbf{1.} Complète l'organigramme du document 1 en nommant les étapes 1 et 2 manquantes, puis décris brièvement ce qui se passe à chacune des cinq étapes de la chaîne.

\textbf{2.} À l'aide du document 2, explique pourquoi le tri des bouteilles par type de plastique est une étape indispensable avant la fusion.

\textbf{3.} Cite deux produits, autres que les bouteilles recyclées, qui peuvent être obtenus à partir des déchets plastiques.

\textbf{4.} Rédige un texte argumentatif de 15 à 20 lignes destiné à sensibiliser les habitants de ton quartier sur la nécessité de transformer et de recycler les déchets. Ton texte devra :
\begin{itemize}
\item rappeler ce qu'est la valorisation des déchets ;
\item citer au moins trois avantages concrets pour l'environnement et l'économie (assainissement, réduction de la déforestation, création d'emplois et de revenus...) ;
\item proposer deux gestes simples que chaque habitant peut adopter au quotidien.
\end{itemize}
\end{exercice}

\newpage

\coursec{Corrigés des exercices}

\begin{corrige}[Exercice 1 — QCM : choisis la (ou les) bonne(s) réponse(s)]
\textbf{1. Réponse b).} La valorisation d'un déchet consiste à lui donner une \cle{nouvelle utilité}, soit sous forme de \cle{matière} (recyclage, transformation), soit sous forme d'\cle{énergie} (valorisation énergétique). L'enfouissement (a) est un simple stockage, le brûlage à l'air libre (c) est une pollution, l'exportation (d) ne fait que déplacer le problème.

\textbf{2. Réponse b).} Dans les pavés, le plastique fondu enrobe les grains de sable et se solidifie en les agglomérant : il joue le rôle de \cle{liant}, comme le ciment dans le béton. Le sable, lui, joue le rôle de \cle{granulat}.

\textbf{3. Réponse c).} L'élimination des encres d'impression, réalisée par action de produits chimiques et par flottation, s'appelle le \cle{désencrage}. La mise en pâte (a) est la dispersion des fibres dans l'eau ; le pressage (b) et le calandrage (d) interviennent après.

\textbf{4. Réponse a).} Les briquettes de papier compressé sont brûlées comme combustible : on récupère l'\cle{énergie} contenue dans le déchet. C'est donc une \cle{valorisation énergétique}, et non une valorisation matière (qui donnerait un nouvel objet).
\end{corrige}

\begin{corrige}[Exercice 2 — Vrai ou faux ? Justifie]
\textbf{1. Vrai.} Le recyclage du papier réutilise les fibres de cellulose déjà extraites : chaque tonne de papier recyclé remplace une tonne de papier neuf et évite l'abattage d'arbres. On estime qu'une tonne de papier recyclé économise environ 15 à 17 arbres.

\textbf{2. Faux.} Les différents plastiques (PET, PEHD, PVC\ldots) ont des températures de fusion et des propriétés différentes ; un mélange de polymères incompatibles donne un matériau recyclé fragile et inutilisable. Le \cle{tri par type de plastique} est donc une étape préalable indispensable.

\textbf{3. Vrai.} La production d'énergie à partir du papier repose sur la \cle{combustion} (ou la carbonisation), réaction chimique au cours de laquelle la matière organique réagit avec le dioxygène pour libérer de la chaleur, du dioxyde de carbone et de la vapeur d'eau : les espèces chimiques de départ sont transformées.

\textbf{4. Faux.} La fabrication de pavés plastique-sable utilise des matières premières très bon marché (déchets plastiques gratuits ou presque, sable abondant) et produit un matériau vendable, plus résistant et moins cher que le pavé en ciment : c'est une activité génératrice de revenus et d'emplois, particulièrement adaptée aux jeunes entrepreneurs.
\end{corrige}

\begin{corrige}[Exercice 3 — Définitions]
\textbf{Déchet :} tout résidu ou objet dont le détenteur se débarrasse parce qu'il n'a plus d'utilité immédiate pour lui. \emph{Exemple camerounais :} les sachets plastiques jetés dans les rues de Douala après consommation de l'eau en sachet.

\textbf{Transformation :} ensemble des opérations qui modifient la nature ou la forme d'un déchet pour en faire un produit nouveau, différent du produit d'origine. \emph{Exemple :} la fabrication de pavés à partir de sachets plastiques fondus et de sable à Buea.

\textbf{Recyclage :} réintroduction d'un déchet dans le cycle de production d'un produit de même nature que le produit d'origine, après traitement. \emph{Exemple :} la fabrication de nouvelles bouteilles plastiques à partir de bouteilles usagées collectées à Yaoundé.

\textbf{Valorisation :} ensemble des procédés qui redonnent une utilité à un déchet, soit sous forme de matière (recyclage, transformation), soit sous forme d'énergie (valorisation énergétique). \emph{Exemple :} la fabrication de briquettes combustibles à partir de vieux papiers à Bafoussam.
\end{corrige}

\begin{corrige}[Exercice 4 — Calcul direct : briquettes de papier]
\textbf{1. Masse totale de briquettes.} Le rendement massique est de \SI{60}{\percent} :
\[ m = 500 \times \dfrac{60}{100} = \SI{300}{kg}. \]
La coopérative obtient \SI{300}{kg} de briquettes par mois.

\textbf{2. Nombre de briquettes.} Chaque briquette a une masse de \SI{250}{g} = \SI{0,250}{kg} :
\[ N = \dfrac{300}{0,250} = 1200 \text{ briquettes}. \]

\textbf{3. Énergie totale.} Le pouvoir calorifique est de \SI{15}{MJ/kg} :
\[ E = 300 \times 15 = \SI{4500}{MJ}. \]
La production mensuelle peut fournir une énergie de \SI{4500}{MJ} (soit \SI{4,5}{GJ}).

\textbf{Point de vigilance :} convertir la masse d'une briquette en kilogrammes avant la division, et ne pas confondre le rendement en masse (\SI{60}{\percent}) avec le pouvoir calorifique (\SI{15}{MJ/kg}).
\end{corrige}

\begin{corrige}[Exercice 5 — Recyclage ou transformation ?]
\textbf{1. Différence.} Le \cle{recyclage} réintroduit le déchet dans la fabrication d'un produit de \emph{même nature} que le produit d'origine (boucle fermée). La \cle{transformation} convertit le déchet en un produit de \emph{nature différente}, ou en énergie.

\textbf{2. Classement.}
\begin{itemize}
\item \textbf{a) Briquettes combustibles : transformation} (valorisation énergétique) : le papier devient un combustible, produit de nature différente, destiné à être détruit par combustion.
\item \textbf{b) Pavés plastique-sable : transformation} : le plastique devient un matériau de construction, produit de nature différente de l'objet d'origine (sachets).
\item \textbf{c) Nouvelles bouteilles : recyclage} : le déchet (bouteille) redonne un produit identique (bouteille).
\item \textbf{d) Papier-cadeau : recyclage} : le papier journal redevient du papier, produit de même nature (fibres de cellulose reformées en feuille).
\item \textbf{e) Papier hygiénique : recyclage} : le carton (produit papetier) redevient du papier ; la nature de la matière (fibres de cellulose) est conservée.
\end{itemize}

\textbf{Remarque :} pour d) et e), on accepte aussi l'argument de la transformation au sens large (le produit fini diffère du produit initial) ; l'essentiel est de justifier de façon cohérente. Au sens strict du cours, la refabrication de papier à partir de papier relève du recyclage.
\end{corrige}

\begin{corrige}[Exercice 6 — La chaîne de recyclage du papier]
\textbf{1. Ordre chronologique :}
\begin{center}
collecte et tri $\Rightarrow$ mise en pâte $\Rightarrow$ désencrage $\Rightarrow$ blanchiment $\Rightarrow$ pressage $\Rightarrow$ séchage $\Rightarrow$ bobinage et découpe.
\end{center}

\textbf{2. Définitions.}
\begin{itemize}
\item \textbf{Mise en pâte :} les papiers triés sont broyés et mélangés à l'eau chaude dans une cuve (pulpeur) ; les fibres de cellulose se dispersent et forment une suspension appelée pâte à papier.
\item \textbf{Désencrage :} élimination des encres d'impression par action de produits chimiques (détergents, soude) et par flottation : des bulles d'air entraînent les particules d'encre vers la surface, où elles sont écumées.
\item \textbf{Pressage :} la feuille humide passe entre des rouleaux qui la compriment et en expulsent l'eau, tout en serrant les fibres les unes contre les autres.
\item \textbf{Séchage :} la feuille pressée est chauffée (cylindres chauffants) pour évaporer l'eau résiduelle et obtenir une feuille sèche et résistante.
\end{itemize}

\textbf{3. Importance du blanchiment.} Le papier hygiénique est en contact direct avec la peau et les muqueuses : il doit être \emph{propre, blanc et dépourvu de résidus d'encre et de contaminants}. Le blanchiment élimine les colorants restants après désencrage et améliore l'aspect et l'hygiène du produit fini. Il est moins exigeant pour le papier-cadeau, qui sera de toute façon coloré.
\end{corrige}

\begin{corrige}[Exercice 7 — Étude de données : papier recyclé contre bois de chauffe]
\textbf{1. Énergie obtenue avec \SI{10}{kg}.}
\begin{align*}
E_{\text{bois}} &= 10 \times 16 = \SI{160}{MJ} ;\\
E_{\text{briquettes}} &= 10 \times 15 = \SI{150}{MJ}.
\end{align*}

\textbf{2. Coûts et coût par mégajoule.}
\begin{align*}
\text{Coût bois} &= 10 \times 150 = \SI{1500}{FCFA} ; & \text{Coût briquettes} &= 10 \times 80 = \SI{800}{FCFA}.\\[2pt]
\text{Coût/MJ bois} &= \dfrac{1500}{160} = \SI{9,375}{FCFA/MJ} \approx \SI{9,4}{FCFA/MJ} ;\\
\text{Coût/MJ briquettes} &= \dfrac{800}{150} \approx \SI{5,33}{FCFA/MJ}.
\end{align*}

\textbf{3. Deux avantages des briquettes.}
\begin{itemize}
\item \textbf{Avantage économique :} le mégajoule fourni par les briquettes coûte environ \SI{5,3}{FCFA} contre \SI{9,4}{FCFA} pour le bois, soit près de \SI{43}{\percent} d'économie ; le budget « cuisson » des ménages est allégé.
\item \textbf{Avantage environnemental :} les briquettes n'entraînent \emph{aucune déforestation} (elles utilisent des déchets papiers) et produisent des fumées modérées, donc moins de pollution de l'air et moins de maladies respiratoires pour les cuisinières.
\end{itemize}

\textbf{Point de vigilance :} le bois a un pouvoir calorifique légèrement supérieur (16 contre 15 MJ/kg) ; c'est le \emph{rapport coût/énergie}, et non le seul prix au kilogramme, qui doit guider la comparaison.
\end{corrige}

\begin{corrige}[Exercice 8 — Schéma de la fabrication des pavés]
\textbf{1. Organigramme.}

\begin{popfigure}
\begin{center}
\begin{tikzpicture}[node distance=0.55cm,
  etape/.style={rectangle, draw=popblue, fill=popblueL, rounded corners, text width=6.2cm, align=center, minimum height=0.9cm, font=\small}]
\node[etape] (e1) {1. Collecte, tri, lavage et séchage des sachets plastiques};
\node[etape, below=of e1] (e2) {2. Fusion du plastique dans une cuve chauffée (environ \SI{180}{\celsius})};
\node[etape, below=of e2] (e3) {3. Incorporation du sable sec et tamisé au plastique fondu, malaxage};
\node[etape, below=of e3] (e4) {4. Moulage, compactage, refroidissement et démoulage du pavé};
\draw[-{Stealth}, thick, popdark] (e1) -- (e2);
\draw[-{Stealth}, thick, popdark] (e2) -- (e3);
\draw[-{Stealth}, thick, popdark] (e3) -- (e4);
\node[draw=poporange, fill=poporangeL, rounded corners, right=0.4cm of e2, font=\small, align=center, text width=3.4cm] (prot) {Équipements de protection obligatoires};
\draw[-{Stealth}, thick, poporange] (prot) -- (e2);
\end{tikzpicture}
\end{center}
\legende{Organigramme de fabrication des pavés plastique-sable : quatre étapes successives, de la collecte des sachets au pavé fini.}
\end{popfigure}

\textbf{2. Étape à risque.} C'est l'\textbf{étape 2 (fusion)} — et la manipulation du mélange chaud à l'étape 3 — qui exige le port d'équipements de protection : \cle{masque} respiratoire (vapeurs irritantes dégagées par la fusion), \cle{gants} thermiques, \cle{blouse} ou tablier, \cle{lunettes} de protection et travail dans un espace \cle{bien ventilé}.

\textbf{3. Pourquoi du sable sec et tamisé ?} Le tamisage élimine les graviers et débris trop gros, ce qui garantit un mélange \emph{homogène} et un pavé sans défaut de structure. Le séchage est indispensable car l'eau, au contact du plastique fondu à \SI{180}{\celsius}, se vaporiserait brutalement : projections de matière brûlante (danger) et formation de bulles ou de pores qui fragiliseraient le pavé.
\end{corrige}

\begin{corrige}[Exercice 9 — Type Bac : l'unité de recyclage du papier de Yaoundé]
\textbf{1. Ordre chronologique et nom des opérations.}
\begin{center}
B (mise en pâte) $\Rightarrow$ C (désencrage) $\Rightarrow$ E (formation de la feuille) $\Rightarrow$ A (pressage) $\Rightarrow$ D (séchage-bobinage).
\end{center}

\textbf{2. Principe de l'étape C (désencrage).} Des produits chimiques détachent les particules d'encre des fibres de cellulose ; l'injection de bulles d'air (flottation) entraîne ces particules vers la surface de la cuve, où elles sont écumées. Elle est indispensable car, sans elle, les encres resteraient mélangées aux fibres : le papier obtenu serait gris, taché et impropre à la fabrication de produits propres comme le papier hygiénique ou les serviettes.

\textbf{3. Caractéristiques de fabrication des trois produits.}
\begin{itemize}
\item \textbf{Serviettes jetables :} feuille fine, absorbante, parfois gaufrée (motifs en relief) et pliée ; blanchiment soigné.
\item \textbf{Papier hygiénique :} feuille très fine, douce et blanche ; blanchiment et hygiène renforcés ; découpe en rouleaux.
\item \textbf{Papier-cadeau :} feuille plus épaisse et résistante, \cle{colorée} ou imprimée de motifs ; le blanchiment peut être allégé puisque la coloration masque les teintes résiduelles.
\end{itemize}

\textbf{4. Masse de papier recyclé produite.}
\begin{align*}
m_{\text{semaine}} &= 2 \times 0{,}70 = \SI{1,4}{t} \text{ par semaine} ;\\
m_{\text{année}} &= 1{,}4 \times 48 = \SI{67,2}{t} = \SI{67200}{kg} \text{ par an}.
\end{align*}

\textbf{5. Arbres économisés et conclusion.}
\[ N = \dfrac{67200}{60} = 1120 \text{ arbres par an}. \]
Cette seule PME évite l'abattage d'environ \num{1120} arbres chaque année : le recyclage du papier contribue directement à la \cle{lutte contre la déforestation}, à la préservation de la biodiversité des forêts camerounaises, tout en assainissant la ville (moins de papiers dans les rues et les caniveaux) et en créant des emplois.

\textbf{Barème indicatif (sur 10 pts) :} ordre et nom des étapes (2 pts) ; principe du désencrage (2 pts) ; caractéristiques des trois produits (1,5 pt) ; calculs avec unités (2 pts) ; nombre d'arbres et conclusion argumentée (2,5 pts).

\textbf{Points de vigilance :} convertir les tonnes en kilogrammes avant la division finale ; justifier chaque étape par le document 1 (le jury attend que la réponse s'appuie sur les descriptions A--E) ; ne pas oublier l'unité dans chaque résultat.
\end{corrige}

\begin{corrige}[Exercice 10 — Type Bac : les pavés écologiques de Buea]
\textbf{1. Principe général.} La technique consiste à utiliser le plastique fondu comme \cle{liant} : chauffé vers \SI{180}{\celsius}, il devient pâteux, enrobe les grains de sable, puis se solidifie en refroidissant et agglomère l'ensemble en un matériau dur et résistant, moulé sous forme de pavé.

\textbf{2. Rôles des constituants et analogie.} Le \textbf{plastique fondu} joue le rôle de liant (il assure la cohésion du matériau) ; le \textbf{sable} joue le rôle de \cle{granulat} (il constitue l'ossature et apporte la masse et la résistance mécanique). Ce pavé est comparable au \textbf{béton} : dans le béton, le ciment est le liant et le sable (avec le gravier) est le granulat ; dans le pavé écologique, le plastique recyclé remplace le ciment.

\textbf{3. Précautions (document 2).}
\begin{itemize}
\item Porter un \textbf{masque} respiratoire et travailler dans un espace ventilé : évite l'inhalation des vapeurs irritantes dégagées par la fusion (risque d'irritation des voies respiratoires).
\item Porter des \textbf{gants thermiques}, une blouse et des lunettes : évite les brûlures au contact de la cuve de fusion et du mélange chaud (\SI{180}{\celsius}).
\end{itemize}

\textbf{4. Masses pour un lot de \SI{500}{kg}.}
\begin{align*}
m_{\text{plastique}} &= 500 \times \dfrac{30}{100} = \SI{150}{kg} ;\\
m_{\text{sable}} &= 500 \times \dfrac{70}{100} = \SI{350}{kg}.
\end{align*}
Vérification : $150 + 350 = \SI{500}{kg}$. Il faut donc \SI{150}{kg} de sachets plastiques et \SI{350}{kg} de sable.

\textbf{5. Avantages.}
\begin{itemize}
\item \textbf{Pour l'environnement :} (i) assainissement des villes par l'élimination des sachets plastiques qui bouchent les caniveaux et provoquent des inondations ; (ii) réduction de la pollution plastique des sols et des cours d'eau (un plastique met des siècles à se dégrader).
\item \textbf{Pour l'économie locale :} (i) création d'emplois et de revenus (collecte, fabrication, vente) pour les jeunes ; (ii) production d'un matériau de construction moins cher que le pavé en ciment, ce qui réduit le coût des travaux de voirie.
\end{itemize}

\textbf{Barème indicatif (sur 10 pts) :} principe (2 pts) ; rôles et analogie avec le béton (2 pts) ; deux précautions justifiées par le risque évité (2 pts) ; calculs avec vérification (2 pts) ; quatre avantages bien distingués (2 pts).

\textbf{Points de vigilance :} exiger la justification de chaque précaution par le document 2 ; bien distinguer « avantage environnemental » et « avantage économique » (le jury sanctionne les réponses confondues) ; dans le calcul, appliquer les pourcentages à la masse totale du lot.
\end{corrige}

\begin{corrige}[Exercice 11 — Type Bac : recyclage des bouteilles plastiques et sensibilisation]
\textbf{1. Organigramme complété et description des étapes.}
\begin{center}
Collecte $\Rightarrow$ \textbf{Étape 1 : tri par type de plastique} $\Rightarrow$ Lavage $\Rightarrow$ \textbf{Étape 2 : broyage en paillettes} $\Rightarrow$ Fusion et moulage.
\end{center}
\begin{itemize}
\item \textbf{Collecte :} rassemblement des bouteilles usagées auprès des ménages, écoles et points de vente.
\item \textbf{Tri :} séparation des bouteilles selon le type de polymère (PET, PEHD, PVC\ldots) et élimination des corps étrangers.
\item \textbf{Lavage :} élimination des saletés, étiquettes et résidus de boisson pour obtenir un plastique propre.
\item \textbf{Broyage :} les bouteilles propres sont réduites en paillettes, plus faciles à fondre de façon homogène.
\item \textbf{Fusion et moulage :} les paillettes sont fondues puis moulées (soufflage, injection) pour former de nouvelles bouteilles.
\end{itemize}

\textbf{2. Nécessité du tri.} D'après le document 2, chaque type de plastique possède une température de fusion et des propriétés propres. Si l'on fond ensemble des polymères incompatibles, ils ne se mélangent pas correctement : le matériau recyclé obtenu est \emph{fragile et inutilisable}. Le tri garantit donc l'homogénéité et la qualité mécanique des bouteilles recyclées.

\textbf{3. Autres produits issus des déchets plastiques.} Les \cle{pavés} de voirie (plastique + sable) ; les \cle{tuiles}, ardoises, poteaux, bancs publics, fils pour tissage ou objets d'artisanat (deux réponses suffisent).

\textbf{4. Texte argumentatif (modèle de production attendue).}

\emph{Chers habitants de notre quartier, chaque jour, nous jetons des montagnes de papiers, de cartons et de sachets plastiques qui encombrent nos rues, bouchent nos caniveaux et provoquent des inondations à la saison des pluies. Pourtant, ces déchets ne sont pas une fatalité : ils peuvent être valorisés. La valorisation consiste à redonner une utilité à un déchet, soit en le transformant en nouveau produit, soit en le recyclant, soit en en tirant de l'énergie.}

\emph{Transformer et recycler nos déchets présente des avantages considérables. D'abord, cela assainit notre environnement : moins de plastiques dans les ruisseaux signifie moins d'inondations, moins de moustiques et donc moins de paludisme. Ensuite, le recyclage du papier réduit l'abattage des arbres : recycler une tonne de papier, c'est préserver une quinzaine d'arbres de nos forêts. Enfin, la valorisation crée des emplois et des revenus : à Douala et à Buea, des jeunes fabriquent déjà des pavés, des briquettes combustibles et des bouteilles recyclées qu'ils vendent sur le marché local.}

\emph{Chacun de nous peut agir par deux gestes simples : trier ses déchets à la maison en séparant les papiers et les plastiques des ordures fermentescibles, et les déposer ou les vendre aux points de collecte au lieu de les brûler ou de les jeter dans la nature. Ensemble, faisons de nos déchets une richesse pour notre quartier, notre santé et notre économie.}

\textbf{Barème indicatif (sur 10 pts) :} organigramme complété et étapes décrites (3 pts) ; justification du tri à partir du document 2 (2 pts) ; deux autres produits (1 pt) ; texte argumentatif respectant les trois consignes (4 pts : définition de la valorisation 1 pt, trois avantages 2 pts, deux gestes 1 pt).

\textbf{Points de vigilance :} dans le texte, le jury attend explicitement la définition de la valorisation, \emph{au moins trois} avantages distincts (environnementaux \emph{et} économiques) et \emph{deux} gestes concrets ; veiller à la longueur (15 à 20 lignes) et à la qualité de l'argumentation (connecteurs logiques, exemples locaux).
\end{corrige}

\newpage

\coursec{Fiche bilan}

\courssub{Carte mentale de la leçon}

\begin{popfigure}
\begin{tikzpicture}[mindmap, grow cyclic, every node/.style={concept}, text width=2.4cm, align=center,
  root concept/.append style={concept color=popblue, font=\bfseries, text=white, text width=2.8cm},
  level 1/.append style={level distance=3.6cm, sibling angle=60, font=\small},
  level 2/.append style={level distance=2.2cm, font=\scriptsize, text width=1.9cm}]
\node[root concept]{Transformation et recyclage des déchets}
  child[concept color=popgreen] { node {Généralités}
    child { node {Déchet, tri, recyclage, valorisation} } }
  child[concept color=poporange] { node {Énergie à partir du papier}
    child { node {Incinération avec récupération de chaleur, méthanisation (biogaz), cogénération électricité/vapeur} } }
  child[concept color=poppurple] { node {Recyclage du papier}
    child { node {Serviettes jetables, papier hygiénique, papier-cadeau} } }
  child[concept color=poppink] { node {Plastiques en bouteilles}
    child { node {Tri par résine, broyage, lavage, extrusion, soufflage} } }
  child[concept color=popgold] { node {Plastiques en pavés}
    child { node {Fusion avec sable, moulage compression, refroidissement} } }
  child[concept color=popblue!70] { node {Avantages et sensibilisation}
    child { node {Environnement, économie circulaire, emplois verts} } };
\end{tikzpicture}
\legende{Carte mentale de la leçon : autour du nœud central, six branches résument les notions essentielles — (1) généralités sur les déchets, (2) valorisation énergétique du papier (incinération, méthanisation), (3) recyclage du papier en produits d'hygiène et d'emballage, (4) recyclage des plastiques (PET, PEHD) en nouvelles bouteilles, (5) transformation des plastiques mélangés en pavés composites sable-plastique, (6) avantages écologiques, économiques et actions de sensibilisation communautaire.}
\end{popfigure}

\courssub{L'essentiel à retenir}

\begin{aretenir}[Fiche bilan — Transformation et recyclage des déchets]
\textbf{Définitions clés}
\begin{itemize}
  \item \cle{Déchet} : tout résidu de production, de transformation ou d'utilisation, abandonné ou destiné à l'abandon.
  \item \cle{Tri} : opération consistant à séparer les déchets selon leur nature (papiers/cartons, plastiques, matière organique, métaux, verres) avant tout traitement.
  \item \cle{Recyclage} : réintroduction d'un déchet dans le cycle de production dont il est issu, en remplacement total ou partiel de la matière première vierge.
  \item \cle{Transformation} : modification de la forme ou de la nature d'un déchet pour obtenir un nouveau produit à valeur ajoutée (ex. pavés, objets moulés).
  \item \cle{Valorisation} : ensemble des opérations qui donnent une seconde vie au déchet, soit comme \textbf{matière} (recyclage, transformation), soit comme \textbf{énergie} (incinération avec récupération d'énergie, méthanisation).
\end{itemize}

\textbf{Techniques à connaître par cœur}
\begin{center}
\begin{tabularx}{\linewidth}{|l|X|}
\hline
\rowcolor{popblueL} \textbf{Filière} & \textbf{Étapes essentielles} \\
\hline
Valorisation énergétique du papier & Collecte et tri $\rightarrow$ séchage et broyage (combustible solide de récupération) $\rightarrow$ combustion en incinérateur à grille ou lit fluidisé $\rightarrow$ récupération de la chaleur (vapeur) $\rightarrow$ cogénération : électricité (turbine/alternateur) + chaleur (chauffage urbain, processus industriels) ; \emph{alternative} : méthanisation en digesteur anaérobie $\rightarrow$ biogaz (CH\textsubscript{4} 50--65\,\%) $\rightarrow$ cogénération ou injection réseau. \\
\hline
Recyclage du papier (produits d'hygiène/emballage) & Tri et élimination des impuretés (agrafes, plastiques) $\rightarrow$ mise en pâte (défibrage hydropulpeur dans l'eau) $\rightarrow$ désencrage (flottation, lavage) et blanchiment (peroxyde d'hydrogène, hydrosulfite) $\rightarrow$ épuration, égouttage, pressage et séchage (machine à papier) $\rightarrow$ bobines mères $\rightarrow$ conversion : serviettes jetables, papier hygiénique, papier-cadeau. \\
\hline
Recyclage bouteilles plastiques (PET, PEHD) & Tri manuel/optique par type de résine (code 1 PET, 2 PEHD) $\rightarrow$ broyage en paillettes (flakes) $\rightarrow$ lavage à chaud (soude, détergent) et rinçage $\rightarrow$ séchage $\rightarrow$ extrusion et granulation (granulés r-PET, r-PEHD) $\rightarrow$ injection-soufflage (PET) ou extrusion-soufflage (PEHD) en nouvelles bouteilles. \\
\hline
Transformation plastiques en pavés (valorisation matière) & Collecte plastiques souples/rigides mélangés (PE, PP, PS) $\rightarrow$ broyage $\rightarrow$ fusion dans extrudeur/malaxeur $\rightarrow$ incorporation de sable (charge minérale, 60--70\,\% masse) $\rightarrow$ malaxage homogène $\rightarrow$ moulage par compression dans moules métalliques chauffés $\rightarrow$ refroidissement sous pression $\rightarrow$ démoulage : pavés résistants (compression > 20 MPa) pour cours, trottoirs, voirie légère. \\
\hline
\end{tabularx}
\end{center}

\textbf{Avantages de la valorisation}
\begin{itemize}
  \item \emph{Pour l'environnement} : réduction drastique des volumes d'ordures enfouis ou brûlés à l'air libre ; limitation de la pollution des sols (lixiviats), des eaux (microplastiques, métaux lourds) et de l'air (dioxines, furannes, particules) ; économie des ressources naturelles (bois, eau, pétrole, sable) ; diminution des émissions de gaz à effet de serre (évitement CH\textsubscript{4} en décharge, substitution énergie fossile).
  \item \emph{Pour l'économie} : création d'emplois locaux non délocalisables (collecte sélective, tri manuel/mécanique, transformation artisanale ou industrielle) ; production de biens vendables (granulés, pavés, papier, électricité, biogaz) générant des revenus ; réduction des coûts municipaux de collecte et d'assainissement ; développement de l'économie circulaire et de l'entrepreneuriat jeune (start-up de recyclage à Yaoundé, Douala, Bafoussam).
\end{itemize}

\textbf{Méthodes express}
\begin{itemize}
  \item Pour \emph{expliquer une technique} : toujours ordonner les étapes chronologiquement (collecte/tri $\rightarrow$ préparation/conditionnement $\rightarrow$ traitement principal $\rightarrow$ finition/produit fini) et nommer précisément le produit final obtenu.
  \item Pour \emph{sensibiliser une communauté} : identifier le problème local concret (ex. bouchons plastiques obstruant les caniveaux à Douala, déchets papiers brûlés dans les cours d'écoles), proposer une filière de valorisation adaptée (bac de tri, point d'apport volontaire, micro-usine de pavés), énumérer les bénéfices directs (revenus, cadre de vie sain, prévention inondations) et inviter à l'action (tri à la source, compostage, achat produits recyclés).
\end{itemize}
\end{aretenir}

\begin{aretenir}[Checklist avant le Bac]
\begin{itemize}
  \item $\square$ Je sais définir rigoureusement : déchet, tri, recyclage, transformation, valorisation (matière vs énergie).
  \item $\square$ Je distingue valorisation \emph{matière} (recyclage papier, bouteilles, pavés) et valorisation \emph{énergétique} (incinération avec cogénération, méthanisation/biogaz).
  \item $\square$ Je sais expliquer, étape par étape, la production d'énergie (électricité + chaleur) à partir du papier/carton.
  \item $\square$ Je connais les étapes du recyclage du papier en pâte secondaire : défibrage, désencrage, blanchiment, séchage $\rightarrow$ serviettes, papier hygiénique, papier-cadeau.
  \item $\square$ Je sais décrire le recyclage \emph{bouteille-à-bouteille} des plastiques (tri résine, broyage, lavage, extrusion, granulés, soufflage).
  \item $\square$ Je sais expliquer la transformation des plastiques mélangés en pavés composites (fusion, adjonction sable, moulage compression, refroidissement).
  \item $\square$ Je cite au moins trois avantages environnementaux (réduction pollution, économie ressources, climat, biodiversité).
  \item $\square$ Je cite au moins trois avantages économiques (emplois verts, revenus, produits marchands, baisse coûts assainissement).
  \item $\square$ Je sais rédiger une réponse ordonnée avec un vocabulaire scientifique précis (pâte à papier, désencrage, granulés r-PET, extrusion, malaxage, cogénération, méthanisation, économie circulaire).
  \item $\square$ Je sais proposer des actions concrètes de sensibilisation de ma communauté au tri à la source et au recyclage (ex. éco-délégués, clubs environnement, partenariat mairie/ONG).
\end{itemize}
\end{aretenir}

\findecours

\end{document}
